\documentclass[parskip=half,a4paper,11pt]{scrreprt}
\usepackage[T1]{fontenc}
\usepackage[utf8]{inputenc}
\usepackage[ngerman]{babel}
\usepackage{lmodern}
\usepackage{comment}
\usepackage{microtype}
\usepackage{ulem}
\usepackage{geometry}
\usepackage{graphicx}
\usepackage{titling}
\usepackage{listings}
\usepackage{csquotes}
\usepackage{tabularx}
\usepackage{booktabs}
\usepackage{amsmath}
\usepackage{amssymb}
\usepackage[round,authoryear]{natbib}
\lstset{
    basicstyle=\setstretch{1}\normalfont\ttfamily\small\mdseries,
    keywordstyle=\normalfont\ttfamily\small\mdseries,
    identifierstyle=\normalfont\ttfamily\small\mdseries,
    commentstyle=\normalfont\ttfamily\small\mdseries,
    stringstyle=\normalfont\ttfamily\small\mdseries,
    breaklines=true,
    breakatwhitespace=true,
    columns=fixed,
    basewidth=0.53em,
    keepspaces=true,
    showstringspaces=false,
    captionpos=b,
    aboveskip=0.6\baselineskip,
    belowskip=0.6\baselineskip,
    frame=single,
    framerule=0.4pt,
    rulecolor=\color{gray!50},
    backgroundcolor=\color{gray!6},
    xleftmargin=0.5em,
    xrightmargin=0.5em,
    framexleftmargin=0.4em,
    framexrightmargin=0.4em
}
\geometry{margin=2.8cm}
\usepackage{xcolor}
\usepackage{setspace}
\usepackage{array}
\usepackage[
hidelinks,
colorlinks=true,
linkcolor=black,
urlcolor=blue,
citecolor=lightgray,
pdftitle={Logische Labyrinthe -- Einleitung},
pdfauthor={CH}
]{hyperref}
\setcounter{secnumdepth}{-1}
\setstretch{1.3}
\begin{document}
\title{%
  \includegraphics[width=0.6\textwidth]{images/labyrinth.png}\\[1cm]
    Logische Labyrinthe
}
\author{Carsten Heisterkamp \\ Sommersemester 2026 \\ HSD PBSA FB 2}
\date{}
\maketitle
\tableofcontents
\chapter{Logische Labyrinthe}

Ausgehend von dem gedruckten Buch \textit{Maze} von \textit{Christopher Manson} und Computerspielen wie \textit{Colossal Cave}, \textit{Myst} oder \textit{Blue Prince} setzen wir uns in diesem Kurs mit dem Genre der \textit{Interactive Fiction} (IF) auseinander und entwerfen Welten, Storylines und Rätsel, untersuchen diese dann mit Hilfe der Logik und übersetzen sie in die deklarative Programmiersprache Prolog. Dabei lernen wir Code nicht als technische Hürde, sondern als gestalterisches Werkzeug für narrative Systeme zu verstehen. Ohne Vorkenntnisse entwickeln wir Schritt für Schritt Welten, die wir mit Hilfe von Logik, Regeln und Bedingungen in lebendige Systeme verwandeln, die uns spielerisch nicht nur durch gedankliche Labyrinthe führen.

\section{Einleitung}

In einer Zeit visueller Überreizung und immer kürzer werdender Aufmerksamkeitsspannen besinnt sich der Kurs bewusst auf die Möglichkeiten des Textes. Ausgehend von Roman Ingardens und Wolfgang Isers Konzept der „Unbestimmtheitsstellen“ \citep{Ingarden1972,Iser1976}, die durch die Lesenden aktiv mit eigenen Erfahrungen und Bildern besetzt werden müssen, betrachten wir die rein textbasierte Narration und Interaktion als gestalterische Möglichkeit.

Ganz im Sinne von Johan Huizingas Konzept des \textit{Homo Ludens} verwenden wir dabei das Spiel als die fundamentale Form menschlichen Lernens und den Ursprung kultureller Praxis. Für den Kurs fungiert das Spiel damit zugleich als epistemische Praxis, in der Wissen nicht nur aufnehmen, sondern entwerfen und erzeugen \citep{Shaffer2006}. 

Indem wir Geschichten als \textit{regelbasierte Systeme} gestalten, entwickeln wir nicht nur Narrative, sondern erschaffen \textit{Magic Circles} \citep{Huizinga1938}, in denen Logik, Spiel und Fantasie untrennbar miteinander verwoben sind. Die Brücke zwischen Literatur, Fiktion, Spiel und Informatik schlagen wir dabei, indem wir die Logik und logische Programmierung als Werkzeug für systematisches Design verwenden. Da IF im deutschsprachigen Raum wissenschaftlich noch kaum erschlossen ist, untersuchen wir dabei gleichzeitig auch ihr spezifisches narratives Potenzial als Form des Kommunikationsdesigns. 

\section{Kursübersicht}

Der Kurs gliedert sich in drei Kapitel, die jeweils aber auch als eigenständige, abgeschlossene Abschnitte gelesen werden können. 

\textbf{Narrative Strukturen in Interactive Fiction}
Wir beginnen mit der Definition eine Narrativen Systems und der Geschichte der interaktiven Literatur und ihren Eigenschaften und Eigenheiten. Hier lernen wir auch was es bedeutet interessante, spannende und mehrschichtige Geschichten zu erzählen, oder besser zu generieren, die interaktiv, nur mit Mitteln des Texts, erzählt und erschlossen werden können.

\textbf{Grundlagen der formalen Logik}
Um diese Geschichte programmieren zu können, müssen wir wissen, wie die Spielmechaniken und Logik übersetzt werden können. Die formale Logik dient dabei der Formalisierung der Welt und ihrer Probleme und Strukturen und ist gleichzeitig auch die Grundlage für das Programmieren mit Prolog.

\textbf{Grundlagen der logischen Programmierung mit Prolog}
Hier lernen wir in der Programmiersprache Prolog zu programmieren, und die Logik unseres Spiels, der Story, ihre Spielmechaniken und Rätsel in Prolog formulieren, sprich zu programmieren und zum Leben erwecken.

\chapter{Narrative Strukturen in Interactive Fiction}

Wer eine Geschichte erzählt, trifft strukturelle, kausale, ästhetische und damit gestalterische Entscheidungen. Was kommt zuerst? Was folgt woraus? Was bleibt ungesagt? Bereits in der \textit{Poetik} (ca. 335 v. Chr.)  unterscheidet \textit{Aristoteles} zwischen bloßer \textit{Ereignisfolge }und \textit{Mythos}, einer \textit{strukturierten Zusammenfügung} der Geschehnisse, also einer \textit{logisch nachvollziehbaren} Handlungsstruktur, die Ereignissen erst \textit{Sinn }verleiht.

Diese Entwicklung der Erzählform lässt sich als eine Entwicklung begreifen, die von der kollektiven \textit{mündlichen} und später \textit{schriftlichen} \textit{Überlieferung} über die Ausdifferenzierung spezialisierter \textit{Künste} wie Theater, Musik und Literatur bis hin zum heutigen \textit{World Building} führt. In dieser aktuellen Entwicklungsstufe verschmelzen ehemals getrennte Disziplinen in interaktiven und hybriden Umgebungen (VR/AR). Das Ziel ist nicht mehr die einzelne, lineare Erzählung, sondern die Erschaffung eines konsistenten narrativen Raums.

\section{Narrative Systeme}

Narrative Systeme sind, damit keine systematische Sammlung von Inhalten, sondern ein \textit{Regelwerk} ihrer \textit{Verknüpfung}, vergleichbar einem \textit{Designsystem}, das einzelne Elemente nicht einfach definiert, sondern die \textit{Logik ihrer Beziehungen }festlegt. 

Der Begriff ist dabei kein disziplinär einheitlicher Terminus. Aus \textit{narratologischer Perspektive} \citep{Genette1972,Schmid2005} bezeichnet er die Struktur des Diskurses, also die Art und Weise, wie eine Geschichte erzählt wird, etwa hinsichtlich Erzählstimme, Zeitstruktur und Perspektive sowie der Formierung von Fakten zu narrativen Einheiten. Der \textit{spieltheoretische (ludologische) Ansatz} \citep{Aarseth1997, Murray1997} begreift es als \textit{ergodischen} \textit{Möglichkeitsraum}, dessen Regelwerk nicht eine fixierte Geschichte vorgibt, sondern die \textit{Bedingungen} erzeugt, unter denen Narrative entstehen können. Die \textit{systemtheoretische Sicht} \citep{Luhmann1984,Bruner1990} versteht Erzählungen als Rahmen, die \textit{Sinn stiften} und \textit{Ordnung schaffen}. Sie helfen uns, die Komplexität der Welt zu bewältigen, indem sie \textit{Ereignisse auswählen} und logisch miteinander \textit{verknüpfen}.

Für diesen Kurs ist vor allem die \textit{Zusammenführung dieser Zugänge} relevant, das narrative System als Regelwerk, das Bedeutung nicht speichert, sondern strukturiert.

Diese relationale Logik findet sich bereits bei \textit{Roland Barthes} und \textit{Algirdas Julien Greimas}. Sie definieren narrative Systeme als formale Systeme aus funktionalen Einheiten. Ereignisse erhalten Bedeutung nicht durch ihr bloßes Vorkommen, sondern durch ihre Relation zueinander \citep{Barthes1966,Greimas1966}. 

Seymour Chatman präzisiert diese Unterscheidung als Differenz zwischen Story (was wir erzählen) und discourse (wie wir es erzählen) \citep{Chatman1978}. Damit folgt Chatman den russischen Formalisten und trennt \textit{Fabula} und \textit{Sjužet}. Was ein narratives System von einer bloßen Enumeration (Aufzählung) unterscheidet, sind \textit{Kausalität}, \textit{Selektion} und \textit{Transformation}. Das bedeutet praktisch, dass dasselbe Material \textit{anders strukturiert} eine \textit{andere Geschichte} und damit eine \textit{andere Erfahrung} und eine \textit{andere Wirkung} erzeugt.

Marie-Laure Ryan erweitert diesen Rahmen auf digitale und interaktive Umgebungen und zeigt, dass narrative Systeme graduell variieren, d. h. sich nach dem Grad der \textit{Interaktivität}, der \textit{Linearität} und der \textit{Verteilung narrativer Agency} zwischen System und Rezipierenden \citep{Ryan2001} unterscheiden. 

Genau diese Variabilität macht den Systembegriff für Interactive Fiction und damit auch für unseren Kurs und die \textit{regelbasierte Formulierung in Logik} und \textit{Programmierung} mit Hilfe von Prolog interessant. Hier ist das narrative System kein geschlossener Text, sondern ein Regelwerk, das Erzählung erst in der Ausführung erzeugt.

Eine  weitere interessante Perspektive liefert der Production Designer \textit{Alex McDowell}, bekannt vor allem durch seine Arbeit an \textit{Minority Report} (2002). McDowell versteht World Building nicht als nachträgliche Ausgestaltung einer fertigen Handlung, sondern als deren eigentlichen Ausgangspunkt. Die Welt, ihre Regeln, Infrastrukturen und sozialen Ordnungen, wird zuerst entworfen und Geschichten entstehen dann organisch aus ihr heraus. Obwohl dieser Ansatz vorwiegend aus Interviews, Workshops und Produktionspraktiken am \textit{World Building Media Lab} (USC) bekannt ist und keine klassische akademische Publikationsform hat, formuliert er einen Begriff des narrativen Systems, der mit den oben genannten Perspektiven strukturell übereinstimmt. Die \textit{Welt }und ihre \textit{Regeln} sind die Basis, die \textit{Geschichte }ist nur ein \textit{Ergebnis }davon. Wir nutzen McDowells Methode später, um eigene Welten und Stoffe zu entwerfen.

\section{Interaktive Fiktion}

Interactive Fiction, im engeren Sinne, bezeichnet eine text- und regelbasierte, üblicherweise einzeln rezipierte Erzählform, deren Interaktionslogik primär sprachlich und deren Rätselprinzip strukturell verankert ist. Im deutschsprachigen Raum ist diese Form wissenschaftlich bislang kaum erschlossen und wird häufig mit \textit{Text-Adventure} gleichgesetzt oder als Unterform digitaler Spiele behandelt, obwohl sie nach Nick Montfort einer eigenständigen literarischen und kulturellen Tradition angehört \citep{Montfort2003}. Besonderes Potenzial liegt dabei in der Anwendung literaturtheoretischer Konzepte von Isers \textit{Leerstellen} über Genettes \textit{Narratologie} bis zu Ryans \textit{Storyworld}-Begriff auf eine Form, die \textit{Lesen}, \textit{Denken} und \textit{Entscheiden} vereint.

Es ist ein System aus \textit{Wegen}, \textit{Entscheidungen }und\textit{ Konsequenzen}, das seinen Bewohnenden\textit{ keine Übersicht} gewährt, sondern nur den nächsten Schritt. Es ist kein Zufall, dass Interactive Fiction dabei immer wieder auf ein Bild zurückgreift, das des \textit{Labyrinths}. Denn IF ist ein Zwitterwesen, das in diese Labyrinthe erschafft und dabei weder reines Spiel noch reine Literatur, weder Informatik noch Narratologie ist, sondern etwas, das sich erst im Spannungsfeld dieser Disziplinen vollständig erkennen lässt.

\section{Ursprünge}

Die Wurzeln der Interactive Fiction reichen weit vor den Computer zurück. Erzähltheoretisch lässt sich ihre Grundstruktur , die Verzweigung, die Entscheidung, die Konsequenz,  bis zu den \textit{kombinatorischen} Erzählexperimenten des Barock zurückverfolgen. Georg Philipp Harsdörffers Poetischer Trichter \citep{Harsdoerffer1647} beschreibt mechanische Verfahren zur Textgenerierung, und Raymond Queneaus Cent mille milliards de poèmes \citep{Queneau1961} realisiert das Prinzip konsequent als kombinatorisches Objekt. 

Jorge Luis Borges' Essay El jardín de senderos que se bifurcan \citep{Borges1941} formuliert das narrative Prinzip der Verzweigung schließlich als Poetologie: ein \textit{Labyrinth aus Möglichkeiten, in dem jede Entscheidung eine andere Welt öffnet} und eine \textit{andere schließt}. Analog und materiell manifestiert sich IF in Rätselbüchern, Wargames und Tabletop-Rollenspielen. Gary Gygax' und Dave Arnesons Dungeons and Dragons \citep{Gygax1974} ist dabei besonders relevant. Es überträgt die Logik der \textit{regelbasierten Welterkundung} in ein kollektives, \textit{improvisiertes Erzählsystem} und etabliert damit Konventionen, \textit{Inventar}, \textit{Raumnavigation}, \textit{Statuswerte}, die die digitale Interactive Fiction unmittelbar prägen werden. 

Johan Huizingas Homo Ludens \citep{Huizinga1938} liefert den anthropologischen Rahmen. Spiel ist nicht Unterhaltung, sondern die \textit{fundamentale Form}, in der Menschen Regeln erproben, Rollen einnehmen und Bedeutung produzieren. Der Magic Circle markiert den vom Alltag abgegrenzten Handlungsraum, in dem diese Regeln gelten und Sinn erzeugt wird. IF bewegt sich genau innerhalb dieses Rahmens, indem es Spiel, Narration und Entscheidung in einer textbasierten Form verbindet.

Mit dem Computer entstehen neue Möglichkeiten der Formalisierung. Will Crowthers \textit{Colossal Cave Adventure} (1976), ursprünglich als Höhlenkartierung für seine Kinder entworfen, gilt als erster digitaler Text, der Navigation, Rätsellösung und Narration in einem Programm vereint. Die daraus folgende Gründung von Infocom (1979) und Werke wie \textit{Zork} (1980) oder \textit{A Mind Forever Voyaging} (1985) etablieren Interactive Fiction als eigenständiges kulturelles Feld. Mit der Verbreitung des Internets entstehen schließlich\textit{ kollaborative} und \textit{hypermediale }Varianten, die die Form erneut erweitern und die Grenzen zwischen Autor:innen- und Leser:innenschaft grundlegend verschieben.

\section{Ergodische Literatur}

Der Begriff, der Interactive Fiction am präzisesten von benachbarten Formen abgrenzt, stammt nicht aus der Literaturwissenschaft, sondern aus der statistischen Physik. Espen Aarseth entlehnt in Cybertext: Perspectives on Ergodic Literature \citep{Aarseth1997} den Terminus "ergodisch" aus der Thermodynamik (griech. ergon: Arbeit, hodos: Weg) und bezeichnet damit Texte, deren Lektüre \textit{nicht-triviale Arbeit} erfordert. Es benötigt Entscheidungen, Navigation, Eingaben, die den Text erst\textit{ konstituieren}. Der Begriff ist bewusst weiter gefasst als „interaktiv" oder „digital": Er beschreibt eine Eigenschaft des \textit{Textverhaltens}, nicht seines Trägermediums, und schließt damit analoge Formen wie das Buch MAZE (Christopher Manson, 1985) ebenso ein wie digitale Textsysteme.

Aarseth unterscheidet dabei verschiedene Dimensionen der Textualität: \textit{Traversalität}, \textit{Determinismus}, \textit{Transienz}, \textit{Perspektive}, die in unterschiedlichen Konfigurationen unterschiedliche Formen ergodischer Literatur erzeugen. Interactive Fiction stellt in diesem Modell einen spezifischen Typ dar: einen ergodischen Text mit \textit{hohem Grad an Traversalität} (der Lesende wählt einen \textit{Weg durch ein Möglichkeitsfeld}), narrativer Kohärenz und sprachlicher Primärität (Sprache allein baut und steuert die Welt).

Es ist diese Kombination, die IF von Hypertext, Netzliteratur und anderen digitalen Erzählformen unterscheidet. Nicht das Medium, sondern die spezifische Struktur der Leseerfahrung.

\section{Subgenres}
Interactive Fiction unterteilt sich heute in ein Spektrum von Subgenres, die sich nach Medium, Interaktionsform und narrativer Struktur unterscheiden und die sich zum Teil mit oft synonym verwendeten Begriffen überlappen, ohne mit ihnen identisch zu sein.

Das \textbf{Text-Adventure} (oder \textit{Parser-based Interactive Fiction}, auch \textit{Parser-Fiction}) bezeichnet die historisch früheste digitale Form von Interactive Fiction und stellt dessen reinste Form dar. Ein Programm, das Spielenden eine Textsituation beschreibt und auf natürlichsprachliche Eingaben wie »GO NORTH« oder »TAKE LAMP« reagiert. Es bezeichnet eine spezifische Interaktionsform der \textit{parser-basierten} Eingabe, nicht die erzählerische Breite der Form, und markiert den Punkt, an dem Computer erstmals nicht mehr nur als Rechenmaschinen, sondern als Werkzeug für \textit{generative Erzählräume} gesehen wurden.

\textbf{Hyperfiction} und \textbf{Netzliteratur} hingegen bezeichnen Formen, die Verzweigungen primär als Hyperlink-Struktur realisieren. Deren literarischer Schwerpunkt ist oft experimenteller und weniger spielerisch-rätselhaft. Sie entstammen einer anderen Tradition (der digitalen Poetik und Medienliteratur) und weisen eine andere Rezeptionsgeschichte auf.

\textbf{Rollenspiele} (Tabletop-RPG) und \textbf{LARP} (Live Action Role-Playing) teilen mit IF die regelbasierte Weltlogik und die Verteilung narrativer Agency, realisieren diese jedoch in kollektiven, performativen Kontexten. Ihre Nähe zur IF ist strukturell, nicht medial.

\textbf{MUD} (Multi-User Dungeon) bezeichnet die netzwerkbasierte, kollektive Variante der textbasierten IF. Ursprünglich 1978 von Roy Trubshaw und Richard Bartle an der University of Essex entwickelt, übertragen MUDs die Parser-Logik der klassischen IF in einen gemeinsamen, persistenten Raum. \textit{Mehrere Spielende} erkunden dieselbe Textwelt \textit{gleichzeitig}, interagieren miteinander und gestalten die Welt durch ihre Aktionen mit. MUDs sind damit der direkte Vorläufer moderner Online-Rollenspiele (MMORPG), unterscheiden sich von der klassischen IF aber grundlegend in der sozialen Dimension. Das Rätsel ist nicht mehr nur das System, sondern auch die anderen Spielenden. Als Zwischenform zwischen Parser-IF, Tabletop-RPG und vernetztem Spiel markieren MUDs den Punkt, an dem Interactive Fiction von einer solitären zu einer kommunikativen Praxis wird.

Die \textbf{Visual Novel} kombiniert vor allem in der japanischen Spielkultur, verzweigte Narrative mit statischer oder animierter Bildsprache und bildet damit einen multimedialen Sonderfall, der zwar narrativ verwandt, sich ästhetisch und konzeptuell aber deutlich unterscheidet. 

Nicht im engeren Sinne Interactive Fiction, jedoch narrativ verwandt sind Genres und Formen der grafischen, nicht rein textbasierten Computerspiele, die aber in der Tradition der Interactive Fiction stehen, insbesondere:

\textbf{Grafische Interactive Fiction} oder \textbf{Illustrated Parser Fiction} repräsentieren die \textit{britisch-europäische} Tradition der Form des Text-Adventures, die sich vom amerikanischen Infocom-Stil deutlich unterscheidet. Werke wie \textit{The Pawn (1985)}, \textit{Guild of Thieves (1987)} oder \textit{Corruption (1988)} zeichnen sich durch literarisch ambitioniertere Texte, einen der damals leistungsfähigsten Parser überhaupt und später die Integration von Grafik bei gleichzeitig beibehaltenem Fokus auf den Text aus. Sie stehen damit genau an der Schwelle zwischen reiner IF und dem sich entwickelnden Grafik-Adventure.

\textbf{Grafik-Adventure und Point-and-Click-Adventures} teilen die narrative Grundstruktur mit IF und verwenden eine starke Narration, Raumnavigation, Inventarrätsel und verzweigte Handlungsoptionen. Der entscheidende Unterschied liegt im Interaktionsmodus. Statt sprachlicher Eingabe operiert man über grafische Interfaces und visuelle Räume. Das Bild ersetzt den Text als primäres Bedeutungsträgermedium, womit, im Sinne Isers, die produktive Leerstelle weitgehend geschlossen wird. Sie sind damit narrativ verwandt, aber ästhetisch und rezeptionstheoretisch verschieden.

\textbf{Narrative Games} und \textbf{Interactive Drama} stellen die Maximierung narrativer \textit{Immersion} bei gleichzeitiger Minimierung spielerischer Agency dar. Die \textit{Illusion der Entscheidung} überwiegt deren tatsächliche \textit{strukturelle Konsequenz} und deutet damit auf ein zentrales ungelöstes Problem interaktiver Narration. Die Spannung zwischen \textit{agency} und \textit{authorial control}, also die Spannung zwischen dem Wunsch, sinnvolle Maßnahmen ergreifen zu können und die Ergebnisse unserer Entscheidungen und Wahlmöglichkeiten zu sehen, und der Notwendigkeit, dass Erzähler:innen und Designer:innen die Struktur und Kohärenz der Geschichte sicherstellen und Interaktionsmöglichkeiten geben, ohne die narrative Qualität und Steuerbarkeit zu verlieren.

\textbf{Environmental Storytelling} und \textbf{Atmospheric Puzzle Games} teilen die Vorstellung, dass die Spielwelt selbst als Text und primäre Erzählfläche fungiert. Die Narration erschließt sich nicht durch Dialoge oder Exposition, sondern durch Umgebung, Raum, Objekte, Spuren und Stimmung. Spielende werden zum aktiven Dechiffrierer oder Archäologen einer impliziten Geschichte, die aus Fragmenten und räumlichen Hinweisen eine Handlung rekonstruieren. Der Unterschied liegt in der Rolle der Aufgabe. Bei Environmental Storytelling steht die frei zugängliche räumliche Lesbarkeit im Vordergrund, d.h. die Geschichte entsteht aus der Deutung dessen, was man findet und sieht. In Atmospheric Puzzle Games hingegen ist die implizite Narration eng an eine puzzleförmige Struktur gekoppelt. Konkrete Rätsel müssen gelöst werden, um die Welt weiter zu öffnen und die Erzählung freizulegen, womit sie eher digitalen Escape Rooms ähneln.

\textbf{Systemische Interactive Fiction} (oder \textbf{Procedural Storytelling}) bricht mit der traditionellen Autorenkontrolle zugunsten einer radikalen, regelbasierten Agency. Diese Form definiert sich durch prozedurale Generierung, Permadeath und eine strikte systemische Logik. Das Narrative ist hier nicht vorab als Plot fixiert, sondern entsteht als \textbf{emergent narrative}, ein Konzept, das Salen und Zimmerman als Geschichten beschreiben, die nicht geschrieben, sondern aus der \textit{Interaktion} von \textit{Spielregeln} geboren werden \citep{Salen2004}. Das Programm fungiert dabei als \textbf{AI Storyteller} oder \textbf{Story Generator}. Die Geschichte ist das Resultat der Simulation, wobei die Spielenden die oft fragmentierten Systemereignisse (Logs, Statusänderungen) eigenständig zu einer kohärenten Chronik verknüpfen.

\section{Hervorzuhebende Werke}

\textbf{Colossal Cave Adventure} (Will Crowther / Don Woods, 1976) – \textbf{Genre: Parser-Fiction (Urform)}. Es ist der Gründungstext der digitalen IF: eine Höhlenerkundung, die Kartografie, Rätsel und Narration erstmals in einem Programm verbindet und nahezu alle Konventionen der Form etabliert: Raumnavigation, Inventar und knappe Zustandsbeschreibungen. Es definiert IF als System von Befehl und Simulation.

\textbf{Zork} (Infocom, 1980) – \textbf{Genre: Parser-Fiction (Klassik)}. Dieses Werk professionalisiert und popularisiert das Format: reichere Texte, komplexere Welten und ein erheblich verbesserter Parser. Als kommerziell erfolgreichstes IF-Produkt seiner Zeit definiert es die Erwartungen einer ganzen Generation an die logische Tiefe und die literarische Qualität des Genres.

\textbf{The Mask of the Sun} (Ultrasoft, 1982) – \textbf{Genre: Illustrated Parser Fiction}. Ein Pionierwerk, das die Textprimärität der IF erstmals mit statischen Grafiken und einfachen Animationen ergänzt. Es steht an der entscheidenden Schwelle, an der das Bild beginnt, die „Leerstelle“ des Textes zu füllen, während die Interaktion (Parser-Eingabe) noch rein der IF-Tradition verhaftet bleibt.

\textbf{A Mind Forever Voyaging} (Steve Meretzky / Infocom, 1985) – \textbf{Genre: Literarische IF / Parser-Fiction}. Markiert eine entscheidende Erweiterung des narrativen Horizonts. Erstmals steht nicht die Lösung von Rätseln, sondern politische und emotionale Erzählung im Zentrum. Ein Beleg dafür, dass IF als literarische Form weit über ihre spielerischen Ursprünge hinausreicht und gesellschaftskritische Themen behandeln kann.

\textbf{The Pawn} (Magnetic Scrolls, 1985) – \textbf{Genre: Illustrated Parser Fiction (Britische Schule)}. Repräsentiert die britisch-europäische Tradition, die sich durch literarisch ambitioniertere Texte und hochwertige Grafik vom US-Stil unterschied. Es hebt sich durch einen extrem leistungsfähigen Parser ab, der komplexe Satzstrukturen verstand, und markiert den ästhetischen Höhepunkt der hybriden Form.

\textbf{The Secret of Monkey Island} (Lucasfilm Games, 1990) – \textbf{Genre: Point-and-Click Adventure}. Es ist der Antagonist zur klassischen IF: Die Sprache wird hier vom Interface (Parser) zum reinen Dialogelement degradiert. Die Interaktion erfolgt über Mausklicks auf Verben und Grafiken. Es markiert den historischen Moment, in dem die visuelle Repräsentation die Vorherrschaft über die textliche Simulation übernahm.

\textbf{Myst} (Cyan, 1993) – \textbf{Genre: Atmospheric Puzzle Game / Environmental Storytelling}. Hier wird die Welt selbst zum Text. Die Narration vollzieht sich nicht durch Dialoge, sondern durch Raum, Objekte und Spuren. Es hebt sich von der IF durch den Verzicht auf Sprache in der Interaktion ab, bleibt ihr aber strukturell durch das „nicht-triviale“ Durchwandern und Entschlüsseln eines komplexen Systems treu.

\textbf{Photopia} (Adam Cadre, 1998) – \textbf{Genre: Experimental / Post-Classic IF}. Das vielleicht meistdiskutierte Werk der Post-Infocom-Ära: eine fragmentierte, nicht-lineare Erzählung, die klassische Spielkonventionen (wie das Sterben oder Inventarrätsel) bewusst unterläuft. Es verhandelt das literarische Potenzial von IF in einem rein akademischen und künstlerischen Kontext neu.

\textbf{Anchorhead} (Michael S. Gentry, 1998) – \textbf{Genre: Parser-Fiction / Cosmic Horror IF}. Ein kanonisches Werk der postklassischen Parser-IF, stark von Lovecrafts kosmischem Horror inspiriert. Die Geschichte verknüpft Ermittlung, Ortsgeschichte und eskalierende Bedrohung in einer dichten, atmosphärischen Welt mit zeitlicher Dramaturgie über mehrere Tage. Besonders hervorzuheben sind das präzise Worldbuilding, die literarische Prosa und die enge Verzahnung von Narration und Rätselstruktur; in der IF-Rezeption gilt das Spiel bis heute als Referenztext für Setting und Stimmung.

\textbf{Heavy Rain} (Quantic Dream, 2010) – \textbf{Genre: Interactive Drama}. Es überträgt die verzweigte Erzählstruktur der IF in die Ästhetik des Kinos. Die Agency der Spielenden ist hier weniger logisch-rätselhaft als vielmehr reaktiv und emotional. Es hebt sich durch die Maximierung der visuellen Immersion bei gleichzeitiger Minimierung der tatsächlichen spielerischen Freiheit von der klassischen IF ab.

\textbf{Detroit: Become Human} (Quantic Dream, 2018) – \textbf{Genre: Interactive Drama / Choice-based IF}. Es erweitert die Prinzipien von \textit{Heavy Rain} durch eine noch dichtere Verzweigung der Erzählstränge und eine stärkere Betonung moralischer Entscheidungen. Die Spielenden navigieren durch ein komplexes Netz von Konsequenzen, das die narrative Kohärenz wahrt und gleichzeitig eine Vielzahl von möglichen Enden ermöglicht.

\textbf{Howling Dogs} (Porpentine, 2012) – \textbf{Genre: Hyperfiction / Choice-based IF}. Ein Schlüsselwerk der „Twine-Revolution“. Es befreit die IF vom technisch hürdenreichen Parser und ersetzt ihn durch Hyperlinks. Es hebt sich durch eine traumartige, hochgradig subjektive Prosa ab und zeigt, wie IF durch Reduktion auf bloße Wahlmöglichkeiten an emotionaler Wucht gewinnt.

\textbf{Dwarf Fortress} (Bay 12 Games, 2006) ist ein Beispiel für ein hochkomplexes Simulationsspiel, in dem narrative Bedeutung aus dem Zusammenspiel von Regeln, Zufall und Systemdynamiken hervorgeht. Die Darstellung bleibt visuell reduziert, während die eigentliche Erzählung aus emergenten Spielereignissen entsteht. Damit verschiebt sich die Produktion narrativer Bedeutung von vorgegebenen Handlungssträngen hin zur Interpretation eines komplexen Systems.

\textbf{RimWorld} (Ludeon Studios, 2013) – \textbf{Genre: Systemic IF / Story Generator}. Dieses Werk perfektioniert das Konzept des \textbf{AI Storytellers}, einer KI, die gezielt dramaturgische Ereignisse in die prozedurale Welt einspeist, um erzählerische Spannungsbögen zu erzeugen. In der Tradition von \citet{Salen2004} steht hier nicht das spielerische Ziel, sondern das Entstehen einer einzigartigen, oft tragischen Chronik im Vordergrund. Es ist das moderne Äquivalent zur IF-Simulation, bei dem das System die Rolle des Autors übernimmt.

\textbf{80 Days} (Inkle, 2014) – \textbf{Genre: Narrative Strategy / Choice-based IF}. Eine moderne Evolution, die hochqualitative Literatur mit strategischen Spielmechaniken verbindet. Es hebt sich von der klassischen IF durch eine dynamische Weltkarte und Ressourcenmanagement ab, bleibt aber im Kern ein Werk, in dem das Lesen des Textes und das Treffen von Entscheidungen die primäre Handlungsebene bilden.

\textbf{Blue Prince} (Dogubomb, 2025) – \textbf{Genre: Modern Environmental Puzzle}. Ein aktuelles Beispiel für die Rückbesinnung auf die rätselhaft-narrative Form. Es kombiniert prozedurale Architektur mit narrativer Archäologie. Es hebt sich von Myst durch eine strengere, fast mathematische Systemlogik ab und wird als Beleg für eine Renaissance der anspruchsvollen, systemorientierten Interactive Fiction gewertet. Christopher Manson hat das Spiel \textit{Blue Prince} maßgeblich als kreatives Vorbild inspiriert und steuerte dafür vier exklusive, für ein komplexes Rätsel genutzte Gemälde im ikonischen Stil seines Buches \textit{MAZE} bei.

\textbf{MAZE} (Christopher Manson, 1985) – \textbf{Genre: Choose Your Own Adventure / Analoge IF / Puzzle Book}. Ein gedrucktes Buch, das die Prinzipien der IF in ein physisches Objekt übersetzt. Es besteht aus 45 Seiten, die nicht linear gelesen werden können, sondern durch eine Kombination aus Text, Illustration und räumlicher Anordnung ein Labyrinth aus Möglichkeiten bilden. Es zeigt, dass die Logik der IF nicht an digitale Medien gebunden ist, sondern als strukturelles Prinzip auch in analogen Formen wirksam sein kann. Es ist auch Vorbild von \textit{Blue Prince} (2025), die die Prinzipien von Maze in digitale Umgebungen übertragen. MAZE ist auch der Ausgangspunkt unseres Kurses, da es ein anschauliches Beispiel für die Prinzipien der mehrschichtigen Bedeutungsebenen als Werkzeug der Weltenkonstruktion und Verrätselung und ihrer Grenzen bietet.

\section{Mögliche Welten}

Der Begriff der \textit{possible worlds} in der analytischen Philosophie von Gottfried Wilhelm Leibniz bis Saul Kripke als logisches Instrument entwickelt, wurde von der modernen \textit{Narratologie} als zentrales Modell für die Erzähltheorie adaptiert. 

\textbf{Gottfried Wilhelm Leibniz} (\textit{Theodizee}, \citep{Leibniz1710}) begründete das Konzept als Entwurf widerspruchsfreier Systeme im göttlichen Verstand. Für das Handwerk der IF ist sein Prinzip der \textit{Kompossibilität} (\textit{Compossibilité}) entscheidend: Eine Welt ist keine bloße Ansammlung von Objekten, sondern eine Menge von Elementen, die logisch miteinander vereinbar sind. IF-Design bedeutet in diesem Sinne, die interne \textit{Konsistenz der Welt} so zu gestalten, dass Regeln und Fakten niemals kollidieren (z. B. die Unmöglichkeit, dass ein Objekt gleichzeitig an zwei Orten existiert). In \textit{Zork} (1980) etwa darf die Lampe entweder im Inventar oder im Raum liegen, aber niemals an beiden Orten zugleich, eine \textit{elementare Kompossibilitätsbedingung}, die das Weltsystem \textit{konsistent} hält.

\textbf{Saul Kripke} formalisierte dieses Modell in der Modallogik durch die Einführung der \textit{Erreichbarkeitsrelation} (\textit{Accessibility Relation}) \citep{Kripke1972}. In einem Kripke-Modell wird bestimmt, welche Welten von einem gegebenen Standpunkt aus \textit{logisch} oder\textit{ physisch }zugänglich sind. Auf die Interactive Fiction übertragen bedeutet dies: Das Programm definiert den \textit{Zustandsraum} aller \textit{potenziellen Weltmomente}. Jede \textit{Interaktion }der Spielenden ist ein \textit{Navigationsakt durch diese Relationen}. 

Kripkes Konzept der \textit{starren Designatoren} (\textit{Rigid Designators}) sichert dabei die \textit{Identität der Spielobjekte} über verschiedene\textit{ Zustände} hinweg. Ein Schwert bleibt logisch dasselbe Objekt, unabhängig davon, ob es im Inventar liegt oder in einem Stein steckt. In \textit{Colossal Cave Adventure} (1976) ist der Raum hinter dem Gitter nur erreichbar, wenn der Schlüssel im Inventar ist: Die Erreichbarkeitsrelation zwischen den Weltzuständen \textit{vor dem Gitter ohne Schlüssel} und \textit{hinter dem Gitter} ist an eine Bedingung geknüpft. Kripkes Rigid Designator zeigt sich darin, dass der \textit{Schlüssel seine Identität }behält, egal ob er am Boden liegt, im Inventar steckt oder bereits benutzt wurde. 

Das scheint auf den ersten Blick trivial und intuitiv. Kripke überwindet mit den starren Designatoren jedoch das Problem des klassischen \textit{Deskriptivismus}, wie er von \textit{Gottlob Frege und Bertrand Russell} geprägt wurde \citep{Frege1892,Russell1905}. In deren Tradition galt ein \textit{Name} (oder ein Objekt) lediglich als ein „\textit{Bündel von Beschreibungen}“. Identität war demnach zwingend an spezifische Eigenschaften geknüpft. Ohne Kripkes logische Verankerung würde ein Objekt in einem dynamischen System seine Identität verlieren, sobald sich sein Zustand ändert, etwa wenn ein Schlüssel vom Boden ins Inventar wandert und die Russell’sche Beschreibung „das goldene Ding an Position X“ nicht mehr zutrifft.

Durch die \textit{Entkoppelung} von \textit{Referenz} (dem fixierten Verweis auf die Entität) und \textit{Beschreibung} (den veränderlichen Attributen) stellt Kripke sicher, dass ein Objekt über alle möglichen Welten und Spielzustände hinweg logisch dasselbe bleibt.

Für die \textbf{Design Studies} bildet dies das theoretische Fundament für \textbf{Objektpermanenz} und Systemkohärenz. Die Essenz eines Design-Elements bleibt durch diesen starren Verweis stabil, während seine grafische Darstellung, Position oder Funktion beliebig variieren kann. Dies ermöglicht es Designern, komplexe interaktive Logiken zu entwerfen, in denen die Identität eines Objekts trotz radikaler Transformationen seiner Erscheinungsform , etwa bei einem Redesign oder einem Statuswechsel im Interface, konsistent gewahrt bleibt.

Als strukturelle Analogie lässt sich in der Informatik das Prinzip des starren Designators anhand der \textbf{UUID (Universally Unique Identifier)} nennen, die einem Objekt eine unveränderliche, abstrakte Adresse zuweist, welche dessen Identität im System festschreibt, ungeachtet wechselnder Attribute oder Speicherorte.

\textbf{Marie-Laure Ryan} überträgt diese Logik schließlich auf fiktionale Welten. Eine Erzählung konstruiert statt einer einzelnen Welt, ein dynamisches Feld von Möglichkeiten. Sie unterscheidet zwischen der \textit{tatsächlichen Welt} (\textit{Actual World}) innerhalb der Fiktion und den \textit{alternativen möglichen Welten} (\textit{Alternative Possible Worlds}), die von den Figuren (oder Spielenden) geglaubt, gewünscht oder befürchtet werden \citep{Ryan1991}. In \textit{Photopia} (1998) wechselt die Erzählung zwischen verschiedenen Zeitebenen und Figurenperspektiven Die Spielenden navigieren nicht durch eine einzige Welt, sondern durch ein Feld möglicher Welten, die erst rückblickend zu einer kohärenten Geschichte konvergieren. Ryans Modell der Actual World und ihrer Alternativen wird hier zur Spielstruktur selbst.

Für IF bedeutet dieses Modell, dass das Programm das Possible Worlds System \textit{ist}.\textit{ Es definiert den Zustandsraum aller erreichbaren Weltmomente und damit implizit alle Geschichten, die in ihm gespielt werden können}. Gestalten von IF  in diesem Sinne bedeutet, \textit{mögliche Welten zu entwerfen}, ihre Grenzen zu setzen und ihre \textit{innere Logik} \textit{konsistent} zu halten, so dass sie als \textit{bedeutungsvolles Spielfeld} erfahrbar werden können.

\section{Storyworld building}

Interactive Fiction hat also eine ausgeprägte Konventionsgeschichte, die sowohl Ressource als auch Falle ist. Die dominanten Stoffe der klassischen Ära, wie Dungeon-Erkundung, Schatzsuche, Detektivfall, Science-Fiction-Raumstation, sind nicht zufällig entstanden. Sie liefern natürliche räumliche Strukturen, eine klare Spiellogik und etablierte Rätselkonventionen. Wer sich zu starr an diese Muster klammert, reproduziert oft nur klischeebehaftet Mechanismen und Stereotypen. Was macht heute also gutes und historisch informiertes Worldbuilding aus?

Nick Montfort zeigt, dass das Rätsel in der Interactive Fiction strukturell mit dem Ort verknüpft ist \citep{Montfort2003}. Die Topografie der Welt \textit{ist} die Topografie des Problems. Stoffe eignen sich für IFInteractive Fiction, wenn sie \textit{räumlich denkbar}, \textit{kausal strukturierbar} und \textit{perspektivisch erfahrbar} sind. Wenn die Welt Geheimnisse haben kann, die sich durch Erkundung und Handlung erschließen. Klischee entsteht nicht durch den Stoff selbst, sondern durch unreflektierte Reproduktion seiner Konventionen, ein Dungeon ohne innere Logik, der \textit{MacGuffin} ohne narrativen Grund, die Figur ohne Eigenlogik. 

Espen Aarseth (\textit{Cybertext}, \citep{Aarseth1997}) unterscheidet dabei zwischen \textit{Aporia}, dem Moment des Nicht-Weiterkommens, der produktiven Blockade und \textit{Epiphany} , dem Durchbruch, der die Welt neu lesbar macht. Gute Stoffe für Interactive Fiction erzeugen diesen Rhythmus strukturell, nicht zufällig.

Marcel Danesis (The Puzzle Instinct: The Meaning of Puzzles in Human Life \citep{Danesi2002}) Perspektive auf den kulturellen und kognitiven Reiz von Rätseln stärkt diesen Punkt. Vor allem die kontrollierte Verzögerung zwischen Problem und Lösung erzeugt den Spannungsbogen, der Aporia in  Erkenntnis überführt.

Der Begriff \textit{Storyworld}, von David Herman (\textit{Storytelling and the Sciences of Mind}, \citep{Herman2011}) als kognitives Modell entwickelt, bezeichnet die mentale Repräsentation einer narrativen Welt, die Lesende oder Spielende aus \textit{Textsignalen} konstruieren. Für IF ist diese Konstruktion eindeutig. Die \textit{Raumtexte} sind die \textit{primären Bausteine des Weltmodells}, und ihre Qualität bestimmt unmittelbar die Qualität der \textit{Immersion}. Das Prinzip \textit{show don't tell} , in der IF-Literatur seit Infocom kanonisch ist hier nicht nur stilistisches Gebot, sondern strukturelle Notwendigkeit: Ein Raum , der behauptet, bedrohlich zu sein ("du betrittst einen bedrohlich wirkenden Raum"), erzeugt keine Bedrohung. Ein Raum, dessen Beschreibung Rostspuren, abgesperrte Türen und das Geräusch tropfenden Wassers enthält, erzeugt sie. Mark Bernstein und Diane Greco (\textit{Reading Hypertext}, \citep{Bernstein2009}) beschreiben \textit{Raumtexte} in IF als \textit{game assets} im vollständigen Sinne. Sie sind gleichzeitig ästhetisches Objekt, Informationsträger und Zustandsanzeiger. \textit{Atmosphäre }ist in diesem Modell keine Dekoration, sondern \textit{Informationsvermittlung}.

Aus McDowells Perspektive lässt sich die Storyworld als \textit{systemischer Entwurfsraum} bestimmen. Weltbau umfasst dann nicht nur Schauplätze, sondern die relationale Ordnung von Infrastruktur, Materialität, sozialer Organisation und technischen Bedingungen. Räume und Objekte fungieren in diesem Rahmen als \textit{epistemische Träger}, sie \textit{vermitteln Weltwissen}, bevor es narrativ deutlich wird. Für Interactive Fiction folgt daraus, dass Weltkohärenz primär über konsistente Raumtexte, Objektlogiken und Zustandsrelationen hergestellt wird.

Für die Konstruktion einer Storyworld ist dabei \textit{Polysemie} entscheidend, die wir später noch genauer betrachten werden. Räume, Objekte und Weltzeichen tragen \textit{mehrere mögliche Lesarten}, ohne ihre systemische Funktion zu verlieren. Bei Barthes (\textit{S/Z}, \citep{Barthes1970}) entsteht Bedeutung dadurch, dass beim Lesen \textit{verschiedene Deutungsmöglichkeiten} gleichzeitig wirksam werden. Bei Eco (\textit{Opera aperta}, \citep{Eco1962}; \textit{The Role of the Reader}, \citep{Eco1979b}) bleibt diese Offenheit jedoch an \textit{formale} und \textit{interpretative Grenzen} gebunden, sodass nicht jede Deutung gleichermaßen gültig ist, sondern durch die Struktur des Werks gelenkt wird. Für Interactive Fiction folgt daraus, dass die Storyworld \textit{semantisch mehrdeutig} sein darf,  sie muss aber \textit{entscheidbar} und \textit{ausführbar} bleiben. Genau diese Spannung erzeugt Verrätselung, ohne die Fairness der Weltlogik aufzugeben.

\subsection{Nichtlinearität und Verzweigung}

Die Nichtlinearität der Welt ist nicht willkürlich, sondern eine direkte Folge davon, wie der Einfluss zwischen Akteuren verteilt ist. Ryan (\textit{Narrative as Virtual Reality}, \citep{Ryan2001}) unterscheidet verschiedene Topologien narrativer Verzweigung: die einfache\textit{ Baumstruktur }(jede Entscheidung öffnet einen neuen, von anderen unabhängigen Ast), \textit{das Netz} (Äste konvergieren wieder), die \textit{Flaschenhalsstruktur} (Freiheit zwischen obligatorischen Knotenpunkten) und die \textit{rhizomatische Struktur }(kein privilegierter Ausgangspunkt, keine notwendige Reihenfolge). 

Jede Topologie erzeugt eine andere Erfahrung von Freiheit und Konsequenz. Die kombinatorische Herausforderung, dass Verzweigung exponentiell wächst und damit den Autorenaufwand multipliziert, ist das zentrale praktische Problem der IF-Produktion und hat historisch zur Entwicklung von \textit{Flaschenhals-} und \textit{Konvergenzstrukturen} geführt. Prozedurale Generierung (\textit{procedural generation}) löst dieses Problem algorithmisch, statt vorgefertigter Äste werden Räume, Ereignisse und Texte aus Regelsets \textit{generiert}. Narrativ bedeutet das den Übergang von \textit{autorische Erzählung (authorial narrative)} zu \textit{emergente Erzählung (emergent narrative)}. Salen und Zimmerman (\textit{Rules of Play}, \citep{Salen2004}) beschreiben diesen Übergang als fundamentalen Paradigmenwechsel im Verhältnis von Autor, System und Spielenden.\\

\textbf{AUFGABE: Storyworld building}

Wir starten mit der Gestaltung unserer Welt, die wir im Laufe des Kurses immer wieder überarbeiten werden. In der Welt der Interactive Fiction ist die Karte das Gerüst, die Texte, die Rätsel und die Handlung richten sich daran aus.,

\begin{itemize}
    \item Sammle Ideen für verschiedene Settings und interessante Orte und räumlichen Strukturen 
    \item Erstelle die Karte: Skizziere die Räume und deren Verbindungen ("Raum1 ist östlich von Raum2").
    \item Objekte und Interaktion: Identifiziere 5 Gegenstände in den Beschreibungen, die greifbar sein sollen.
    \item Narration: Welche spielmechanischen, oder erzählerischen Möglichkeiten bieten sie
    \item Erstelle atmosphärische Beschreibungen für die Orte und Gegenstände. 
    \item Achte darauf, dass sie keine Rückschlüsse auf das Aussehen der Spielfigur zulassen (siehe Abschnitt AFGNCAAP-Prinzip).
    \item Das Raumdiagramm lässt sich als formales Modell $G=(V,E,W)$ schreiben: $V$ sind eure Räume, $E$ eure Verbindungen, $W$ eure Bedingungen für Übergänge (etwa: „Tür nur mit Schlüssel passierbar"). Versucht, euer Diagramm in diese Form zu übersetzen. Wir werden diese Schreibweise im Laufe des Kurses noch genauer betrachten.
\end{itemize}

\section{Erzählung und Subjekt}

Wenn ein Erzählsystem den Lesenden direkt adressiert, stellt sich eine grundlegende Frage: Wer oder was spricht hier, und wen spricht es an? In der klassischen Literatur ist die Erzählinstanz oft eine distanzierte, allwissende Stimme oder ein Ich-Erzählender mit einer klaren Perspektive. In der IF hingegen ist die Erzählinstanz strukturell an die Spielfigur gebunden, die gleichzeitig als Adressat und als handelndes Subjekt fungiert. Die Erzählperspektive ist keine stilistische Wahl, sondern eine Setzung darüber, wie ein Subjekt in der Welt erscheint und handelt. Die folgenden Abschnitte untersuchen, wie IF die Beziehung zwischen Erzählinstanz, Spielfigur und Spielenden organisiert, welche Rolle der Parser als Interaktionsmodus dabei spielt und wie Immersion und Handlungsfähigkeit miteinander vermittelt werden.

\subsection{Erzählperspektive}

\textit{„You are standing in an open field west of a white house"} Die zweite Person ist in der Interactive Fiction keine stilistische Entscheidung, sondern eine \textit{strukturelle Notwendigkeit}. Sie löst das Problem der \textit{Identifikation ohne Identitätsfixierung}. Der Spielende ist weder Beobachtender (dritte Person) noch explizit positioniertes Ich (erste Person), sondern ein unbestimmtes \textit{Du}, das die eigene Figur und den eigenen Körper im Text erst durch Handlung konstituiert. Marie-Laure Ryan (\textit{Narrative as Virtual Reality}, \citep{Ryan2001}) beschreibt dies als \textit{immersive second person}: Die Anrede erzeugt einen Sog in die Welt, weil sie die Distanz zwischen Lesendem und Figur grammatikalisch aufhebt. Gegenbeispiele existieren, einige IF-Werke operieren in der ersten Person (\textit{Photopia}, Adam Cadre, 1998, beispielsweise wechselt systematisch die Perspektive) oder experimentieren mit der dritten Person, um Entfremdungseffekte zu erzeugen, Bertolt Brechts Prinzip des Verfremdungseffekts ist hier narratologisch anwendbar. Experimentelle Formen wie \textit{Photopia} zeigen darüber hinaus, dass eine an eine Figur oder ein Bewusstsein gebundene Systemstimme auch \textit{unzuverlässig} sein kann (Booth, \textit{The Rhetoric of Fiction}, \citep{Booth1961}): nicht als klassische Erzählerfigur, sondern als perspektivisch gefärbte Adressierungsform, deren Darstellung die Spielenden durch Handlung und Exploration korrigieren müssen. Die zweite Person bleibt jedoch der konventionelle Standard, weil sie die Spannung zwischen Figur und Spielenden offen hält.

\subsection{Tod des Autors und die Geburt des Spielers}

Roland Barthes' Essay \textit{La mort de l'auteur} \citep{Barthes1968} und Michel Foucaults Antwort \textit{Qu'est-ce qu'un auteur?} \citep{Foucault1969} beschreiben eine Verschiebung der Bedeutungsproduktion vom Autor zum Lesenden. Der Text bedeutet nicht, was der oder die Autor:in intendierte, sondern was der Lesende in der Begegnung mit ihm \textit{produziert}. In der Interactive Fiction radikalisiert sich diese Verschiebung strukturell. Die schreibende Person hat die Welt gebaut, aber die spielende Person entscheidet, welche Geschichte in ihr entsteht. Die \textit{auteur}-Funktion, wie Foucault sie beschreibt, bleibt als \textit{Systementwurf }erhalten , die Regeln, die Grenzen, die Möglichkeiten sind autorisiert, aber die narrative Realisierung ist eine co-produktive Leistung. Henry Jenkins (\textit{Game Design as Narrative Architecture}, \citep{Jenkins2004}) beschreibt den IF-Autor deshalb als \textit{environmental storyteller}, also nicht \textit{Erzähler:in}, sondern \textit{Welterbauer:in}, der Bedingungen für Erzählung schafft, ohne sie zu kontrollieren.

\subsection{Ageless Faceless Gender-Neutral Culturally Ambiguous Adventure Person}

Das AFGNCAAP-Prinzip wurde ursprünglich von Infocom, den Schöpfern von \textit{Zork}, geprägt. Da der Text den Spieler direkt mit „Du“ anspricht, wollte man vermeiden, Details wie Alter, Geschlecht oder Aussehen festzulegen. So konnte jeder die Rolle ohne Reibungspunkte ausfüllen.

\textbf{Maximale Immersion:} Ohne festes Aussehen wird die Barriere zwischen dir und der Spielfigur abgebaut – \textit{du} bist die Person im dunklen Verlies.

\textbf{Zork-Tradition:} Besonders in den späteren \textit{Zork}-Teilen wurde der AFGNCAAP fast schon parodiert, da die Figur oft Dinge tat oder überlebte, die völlig absurd waren, ohne jemals ein Wort zu sagen.

\textbf{Gegenteil zum Charakter:} Während moderne Spiele oft voll vertonte Charaktere mit eigener Hintergrundgeschichte nutzen (wie Geralt in \textit{The Witcher}), ist der AFGNCAAP ein völlig leeres Gefäß.\\

\textbf{AUFGABE: AFGNCAAP}

Entwickel eine kurze Einstiegsszene (ca. 200 Wörter), mit der die Spielfigur im Spiel erscheint und wie sie sie wahrnimmt. 

\begin{itemize}
    \item Verwende das Setting aus Aufgabe 1.
    \item Nutze strikt die 2. Person Singular („Du“).
    \item Beschreibe keinerlei physische Merkmale der Figur (kein Geschlecht, kein Alter, keine Ethnie).
    \item Beschreibe die Umgebung ausschließlich über die Sinne der Spielfigur.
    \item Entwerfe die Reaktion des Spiels auf den Standard-Befehl UNTERSUCHE MICH 
    \item Erstelle eine Liste von drei funktional wichtigen Gegenständen, die die Figur bei sich trägt
    \item Das AFGNCAAP-Prinzip darf nicht verletzt werden und es dürfen keine Rückschlüsse auf die Identität der Figur möglich.
    \item Wo liegen die Grenzen des AFGNCAAP-Modells, wenn man eine tiefgreifende, emotionale Charakterentwicklung erzählen möchte
    \item Die verändert sich die Spielerfahrung, wenn der Parser plötzlich doch ein Pronomen verwenden würde?
\end{itemize}

\subsection{Die Welt ist das Spiel und das Spiel weiß mehr}

IF-Welten behandeln Wissen asymmetrisch. Das System kennt seinen gesamten Zustandsraum, der Spielende kennt nur den aktuellen Moment. Diese Asymmetrie ist die Grundlage jedes Rätsels und zugleich ein ethisches Designprinzip. Graham Nelson (\textit{The Craft of Adventure}, \citep{Nelson1994}) formuliert es als Pflicht des Systems, fair zu spielen: Hinweise müssen vorhanden, zugänglich und interpretierbar sein. Die Narratologie kennt dieses Prinzip als \textit{restricted narration} (Genette, \textit{Discours du récit}, \citep{Genette1972}): Die Erzählinstanz weiß mehr als sie zeigt, und diese Differenz erzeugt Spannung. Interactive Fiction formalisiert dieses Prinzip als Spielmechanik, der Spielende bewegt sich auf einem \textit{Wissensgradienten} durch die Welt, und das \textit{sukzessive Schließen} der epistemischen Lücke ist die\textit{ Erzählung}. McDowells Designperspektive präzisiert die epistemische Asymmetrie der IF weiter. Spannung entsteht nur dann dauerhaft, wenn das System mehr weiß als die Spielenden, dieses Mehrwissen jedoch regelgeleitet und nachvollziehbar organisiert ist. Fairness ist damit keine externe Zusatznorm, sondern eine Eigenschaft der Weltkonstruktion selbst. Rätsel, Hinweise und Hindernisse sind nur dann überzeugend, wenn sie als \textit{Konsequenz konsistenter Weltregeln} lesbar bleiben.

\subsection{Der Parser als Interaktionsmodus}

Die klassische Interaktionsform der IF ist der Parser, ein Programm, das natürlichsprachliche Eingaben (Verb-Objekt-Konstruktionen wie TAKE LAMP, GO NORTH, EXAMINE DOOR) analysiert und in Systemzustände übersetzt. Nick Montfort (\textit{Twisty Little Passages}, \citep{Montfort2003}) beschreibt den Parser als konstitutives Element der Form der die Illusion eines sprachfähigen Gegenübers  erzeugt und damit eine andere Immersionsqualität als jedes grafische Interface.

Alternativen zum Parser sind seit den 1990er Jahren systematisch erprobt worden. Hyperlink-basierte IF (realisiert etwa in Twine) ersetzt die offene Eingabe durch eine Auswahl vordefinierter Optionen und verschiebt damit Agency von der \textit{Produktion} zur\textit{ Selektion}. Menübasierte Systeme (Choice of Games, Inkle Studios) maximieren \textit{Zugänglichkeit} auf Kosten sprachlicher \textit{Offenheit}.

Die Wahl des \textit{Interaktionsmodus} ist dabei keine rein technische, sondern eine \textit{gestalterische Entscheidung }mit direkten \textit{narrativen Konsequenzen}.Während Hyperlink-Systeme ihre Handlungsmöglichkeiten explizit offenlegen (Affordanz), operiert der Parser durch die bewusste \textit{Verschleierung }der Optionen. Dies erzeugt eine spezifische Spannung zwischen der Intention der Spielenden und dem Weltmodell des Programms, ein Phänomen, das Montfort als strukturelles „\textit{vocabulary problem}" beschreibt \citep[Kap.~3]{Montfort2003} und in der IF-Design-Praxis als „\textit{Gap of Incompatibility}“ bekannt ist. Diese Lücke erzwingt jene nicht-triviale Anstrengung, die das Wesen der Ergodik (wie im vorangegangenen Abschnitt dargelegt) ausmacht. Das „Lesen“ der Interactive Fiction besteht hier nicht nur in der Rezeption des Textes, sondern im tastenden, explorativen \textit{Erproben des Systems}, um dessen\textit{ verborgene Grenzen zu dechiffrieren}. Ein Parser-System ist damit\textit{ aktiver}, \textit{fehleranfälliger} und in seiner Simulation von Handlungsunsicherheit \textit{authentischer} als selektionsbasierte Formen.

\subsection{Immersion und Handlungsfähigkeit (Agency)}

Janet Murray beschreibt in \textit{Hamlet on the Holodeck}\textbf{ }Immersion als das Gefühl des \textit{Eintauchens in eine andere Welt}, im Sinne einer\textit{ haptischen Erfahrung} des Eintauchens, nicht als die visuelle Illusion des Films. Agency definiert sie als die Erfahrung, dass eigene Handlungen bedeutungsvolle Konsequenzen in der Welt haben \citep{Murray1997}. 

Beide Qualitäten stehen in einem Spannungsverhältnis. Maximale Immersion kann Agency untergraben (die Welt zieht mich in eine vorbestimmte Geschichte) und maximale Agency wiederum kann die Immersion untergraben (ich erlebe das System als System, nicht als Welt). Ryan beschreibt in \textit{Narrative as Virtual Reality} \citep{Ryan2001}: Die Illusion der Agency erfordert, dass Entscheidungen wahrnehmbar nachwirken. Dies muss nicht notwendigerweise sofort, aber erfahrbar geschehen. Das \textit{verzögerte Echo }einer frühen Entscheidung ist narrativ oft \textit{wirkungsvoller }als unmittelbares Feedback, weil es \textit{Kausalität} und damit\textit{ Bedeutung} erzeugt. 

(\textit{Non-Player Characters (NPCs)}) sind in diesem Rahmen Weltbewohner, deren Verhalten die innere Logik des Systems verkörpert. Sie übernehmen eine strukturelle Vermittlungsfunktion zwischen Regelwerk und Erfahrung und transformieren abstrakte Systemzustände in sozial lesbare Konsequenzen und stabilisieren so die Wahrnehmung von Agency. Zugleich wirken sie als epistemische Instanzen, indem sie Information selektiv freigeben, verzerren oder verweigern und damit den Wissensgradienten der Spielenden modellieren. Zur Immersion tragen NPCs vor allem dann bei, wenn ihr Verhalten regelkonsistent, zeitlich nachvollziehbar und kontextsensitiv bleibt. Ihre Reaktion auf den Spielenden ist der primäre Test dafür, ob die Welt als kohärent erfahren wird.

\subsection{Tod, Scheitern und das Speicherproblem}

In der Interactive Fiction ist der Tod der Spielfigur kein erzählerisches Ereignis, sondern ein epistemischer Vorgang. Er teilt mit, was nicht funktioniert hat. Das System setzt die Spielenden an einen früheren Zustand zurück, aber mit \textit{neuem Wissen}. Diese Tür ist tödlich, dieser Gegenstand war notwendig, diese Reihenfolge falsch. Montfort (\textit{Twisty Little Passages}, \citep{Montfort2003}) beschreibt diesen Mechanismus als \textit{Aporia}, die produktive Blockade, die erst durch Scheitern überwunden werden kann. Das \textit{Sterben} ist in der klassischen Interactive Fiction kein Spielfehler, sondern die \textit{primäre Lernform}.

Die Funktionen \texttt{SAVE} und \texttt{RESTORE} unterbrechen diesen Mechanismus und erzeugen dabei ein narratives Problem, denn wenn Entscheidungen jederzeit rückgängig gemacht werden können, verlieren sie ihre Konsequenz und damit ihre Bedeutung. Agency setzt jedoch voraus, dass Handlungen wirklich zählen. Murray (\textit{Hamlet on the Holodeck}, \citep{Murray1997}) beschreibt die Spannung zwischen \textit{safety} und \textit{risk} wie folgt: Das Speichern schafft Sicherheit, neutralisiert aber zugleich die Kausalstruktur der Erzählung. Manche IF-Werke reagieren darauf bewusst. So erlaubt \textit{Photopia} (1998) beispielsweise kein Sterben, weil die Erzählung ihre Wirkung aus Unausweichlichkeit zieht. \textit{Zork} (1980) macht Tod und \texttt{RESTORE}  zur kalkulierten Designgröße. Die Entscheidung, wie mit Scheitern umgegangen wird, ist damit keine technische, sondern eine erzählerische. Wer im Dunkeln weiterläuft, wird vom \textit{\href{https://zork.fandom.com/wiki/Grue}{Grue}} gefressen und weiß danach, was Licht bedeutet.

\section{Narrative Strukturen entwerfen}

Die bisherige Diskussion narrativer Systeme beschreibt den theoretischen Rahmen. Für den konkreten Entwurf von IF-Welten stellt sich nun die Frage, welche Elemente Handlung (Plot) in einer nicht-linearen, regelbasierten Form erzeugen. In der Interactive Fiction entsteht Handlung nicht nur durch Reihenfolge, sondern durch \textit{Zustandsänderung}, Zugriff auf\textit{ Information} und wahrnehmbare \textit{Konsequenz} auf die Handlung.

\subsection{Makrostruktur}

\textbf{Exposition}: Die Ausgangssituation klärt, wer die Spielfigur ist, wo sie sich befindet und welches Problem den Handlungsraum strukturiert. In der IF liegt diese Funktion meist im ersten Raumtext, im initialen Inventar und im Anfangszustand der Welt.

\textbf{Konflikt}: Neben klassischen Konfliktformen (gegen Umwelt, Figur oder sich selbst) tritt in der IF ein systemspezifischer Konflikt hinzu, die Reibung zwischen Intention der Spielenden und Regellogik des Systems.

\textbf{Steigerung und Komplikation}: Spannung wächst in der IF häufig t\textit{opografisch} und \textit{systemisch}. Räume sind gebunden an \textit{Bedingunge}n, Objekte an \textit{Zustände}, Fortschritt an gelöste Abhängigkeiten, wie Rätseln, oder Ressourcen.

\textbf{Peripetie und Erkenntnis}: Der Wendepunkt ist oft ein Informationsereignis. Neues Wissen verändert nicht nur die Deutung der Geschichte, sondern auch die verfügbaren Handlungsoptionen.

\textbf{Auflösung}: IF erlaubt unterschiedliche Ending-Typen (einheitlich, multiple, offen, Scheitern, zirkelhaft). Die Wahl des Endings ist eine Designentscheidung über Bedeutung, nicht nur über Abschluss.

\subsection{Mikrostruktur}

\textbf{Raum und Szene}: Der Raum ist in der IF die kleinste plot-relevante Einheit. Er ist gleichzeitig Ort, Informationsfläche und Handlungsrahmen.

\textbf{Ereignis}: Narrativ relevant ist eine Aktion dann, wenn sie einen Zustand verändert. Für die Modellierung bedeutet das, dass nicht jede Interaktion bereits ein Ereignis ist.

\textbf{Zeitstruktur}: Reihenfolge, Dauer und Frequenz sind in der IF gestaltbar. Die Zeit der Erzählung entsteht aus Systemlogik und Traversal, nicht aus einer einzigen fixierten Sequenz.

\subsection{IF-spezifische Plotelemente}

\textbf{Zustandsgraph}: Das erzählerische Potenzial der IF liegt im Raum möglicher Zustände und Übergänge. Der gespielte Plot ist eine konkrete Traversal dieses Möglichkeitsraums.

\textbf{Flags und Trigger}: Flags markieren Zustände, Trigger lösen Ereignisse bei erfüllten Bedingungen aus. Sie sind das zentrale Werkzeug narrativer Progression in regelbasierten Welten.

\textbf{Pacing und Progression}: Rhythmus entsteht aus dem Wechsel zwischen Blockade und Durchbruch. Zu wenig Widerstand entwertet Entscheidungen, zu viel Widerstand erzeugt Frustration.

\textbf{Roter Faden}: In nicht-linearen Strukturen ist Kohärenz emergent. Sie entsteht über thematische, kausale und figurale \textit{Kontinuität}, in Form von Wiederholung und Entwicklung von sprachlichen Bildern, Metaphern und Symbolen.

\textbf{Geheimnis}: Das Geheimnis organisiert die \textit{Informationsverteilung}. Es motiviert Erkundung, strukturiert Neugier und steuert den epistemischen Abstand zwischen Welt und Spielenden.

\textbf{Konsequenz und Kontinuität}: Agency wird erfahrbar, wenn Handlungen die Welt verändern und diese Veränderung wahrnehmbar bleibt, auch verzögert.

\section{Metafiktion}

Metafiktion bezeichnet literarische Verfahren, die den fiktionalen Status des Textes thematisieren und damit die Grenze zwischen Welt und Darstellung sichtbar machen. Patricia Waugh (\textit{Metafiction}, \citep{Waugh1984}) definiert sie als \textit{systematische Selbstreflexivität}: Der Text spricht über sich als Text. In der IF nimmt Metafiktion spezifische Formen an, die über die literarische Tradition hinausgehen: Das System kann die Spielenden direkt adressieren, Konventionen des Genres kommentieren oder die Mechaniken des Parsers selbst zum Rätsel machen. Michael Mateas und Andrew Stern (\textit{Facade}, \citep{Mateas2005}) haben gezeigt, wie die Exposition des Systems als narratives Mittel eingesetzt werden kann. Das kanonischste Beispiel bleibt jedoch Zork Zero (Infocom, 1988) und vor allem Hitchhiker's Guide to the Galaxy (Douglas Adams / Steve Meretzky, 1984), in dem die Frustration des Parsers selbst zur komischen und narrativen Strategie wird. Mark McGurl (\textit{The Program Era}, \citep{McGurl2009}) beschreibt diese Selbstreflexivität als konstitutives Merkmal postmoderner Fiktion generell. IF radikalisiert sie, weil das System, über das reflektiert wird, zugleich das Medium der Reflexion ist.

\section{Erzählerische Vielschichtigkeit}

Bedeutung entsteht in Erzählungen selten eindeutig. Schon klassische Rhetorik und moderne Texttheorie haben beschrieben, wie sprachliche Zeichen mehrere Bedeutungsebenen gleichzeitig tragen können. Für die IF ist diese \textit{Vielschichtigkeit} keine dekorative Qualität, sondern ein \textit{gestalterisches Prinzip}. Ein Element, das auf \textit{mehreren Ebenen zugleich} gelesen werden kann, als Hinweis, als Symbol, als Handlungsanweisung, erzeugt jene \textit{hermeneutische Offenheit}, die das \textit{Erkunden} und\textit{ Interpretieren} zur eigentlichen \textit{Spielerfahrung} macht. Die folgenden Abschnitte beschreiben die theoretischen Grundlagen dieser Mehrdeutigkeit und zeigen, welche konkreten Mittel ihr im Rätseldesign zur Verfügung stehen.

\subsection{Polysemie}

Polysemie ist die gezielte, \textit{kontrollierte Offenheit} von Bedeutungen. Roland Barthes beschreibt mit der \textit{Pluralität} des Textes in \textit{S/Z} \citep{Barthes1970}, dass Bedeutung in Codes organisiert ist und erst im Lesevollzug entsteht, also dass Bedeutung nicht fest im Text steckt, sondern beim Lesen entsteht, weil wir verschiedene „Codes“ entschlüsseln. Umberto Eco schreibt in \textit{Opera aperta} \citep{Eco1962}, dass Offenheit kein Verzicht auf Form ist, sondern eine strukturierte Ermöglichung von Interpretation. Barthes unterscheidet fünf solcher Codes, die gleichzeitig aktiv sind: der \textit{hermeneutische Code} organisiert Geheimnisse und offene Fragen, der \textit{proairetische Code} strukturiert Handlungssequenzen und ihre Konsequenzen, der \textit{semantische Code} trägt Konnotationen und Stimmungen, der \textit{symbolische Code} verweist auf tiefer liegende Oppositionen und Bedeutungsfelder, und der \textit{kulturelle Code} verankert das Element im geteilten Weltwissen der Lesenden. Für die Interactive Fiction ist dieses Modell besonders gewinnbringend, weil es beschreibt, auf welcher Ebene ein Element gerade gelesen wird: Ein Schlüssel ist gleichzeitig Handlungsobjekt (proairetisch), Rätselhinweis (hermeneutisch), Symbol für Zugang und Ausschluss (symbolisch) und kulturell vertrautes Requisit. Die Mehrfachcodierung eines Elements ist damit keine Zufälligkeit, sondern das normale Funktionieren von Bedeutung in narrativen Systemen. Offenheit ist kein Chaos, sondern eine bewusst gestaltete Ermöglichung von mehreren Interpretationen. Als Gestalter:innen legen wir fest, wie weit Interpretationen gehen dürfen. Mit seinem Konzept des „Modelllesers“ in dem Text \textit{Lector in fabula} \citep{Eco1987} macht Eco klar, dass diese Offenheit immer auch Grenzen hat. Der Text konstruiert seine Modellleser durch eingeschriebene Codes und Signale, sodass die Freiheit in der Interpretation stets eine gerahmte Freiheit bleibt.

Ein konkretes Beispiel verdeutlicht, wie Polysemie auf verschiedenen Ebenen einer IF-Welt gleichzeitig operieren kann. Betrachten wir eine verlassene Bibliothek, in der ein Tagebuch liegt. Als Objekt ist das Tagebuch ein Gegenstand, der aufgenommen und untersucht werden kann. Als Figurenartefakt gehört es einer abwesenden Person und trägt deren Perspektive in die Welt. Als Hinweis verbirgt es eine Information, die erst durch Kombination mit einem anderen Gegenstand ihren Sinn enthüllt. Als Ereignisträger löst das Lesen eines bestimmten Eintrags einen Zustandswechsel aus, eine Tür öffnet sich, ein NPC verhält sich anders. Kein einzelnes Element ist eindeutig, aber die Spielwelt liefert durch \textit{Exploration} genug \textit{Kontext}, um die richtige\textit{ Lesart} aus dem \textit{Geflecht der Bedeutungen} herauszuarbeiten. Polysemie meint hier also nicht Beliebigkeit, sondern eine \textit{organisierte Offenheit}, die Neugier erzeugt, ohne die Lösbarkeit preiszugeben.

\subsection{Polysemie als Verrätselungsprinzip}

Für Interactive Fiction und Rätseldesign bedeutet das, dass ein Element mehrere Bedeutungen und Interpretationen zulässt, jedoch nicht beliebig viele Lösungen. Ein Hinweis darf verschiedene plausible Lesarten haben, die richtige Lösung muss sich jedoch klar aus den Informationen in der Spielwelt ableiten lassen. So entstehen Rätsel, die hermeneutisch spannend und interpretativ vielfältig sind, ohne frustrierend oder willkürlich zu wirken. Die Herausforderung für das Design besteht darin, die Balance zwischen Offenheit und Lösbarkeit zu halten. Ein Element, das zu viele mögliche Interpretationen zulässt, kann die Spielenden überfordern und den Spielfluss stören. Ein Element, das zu eindeutig ist, verliert seine interpretative Tiefe und damit auch seinen Reiz als Rätsel. Die Wirksamkeit von Rätseln beruht nicht auf bloßer Schwierigkeit, sondern auf einer semantisch gerahmten Unsicherheit, die zu aktiver Deutung motiviert \citep{Danesi2002}.

\subsection{Mittel um Mehrdeutigkeit zu erreichen}

\textbf{Analogie}: Zwei verschiedene Dinge folgen derselben Ordnung. Wer diese gemeinsame Struktur erkennt, versteht auch das Rätsel. 

Beispiel: In einem Tempel stehen vier Schalter. Erst später merkt man, dass ihre richtige Reihenfolge dem Lauf der Sonne entspricht: Morgen, Mittag, Abend, Nacht.

\textbf{Symbol}: Ein Ding taucht mehrfach auf und bekommt nach und nach eine tiefere Bedeutung. 
Es ist nicht nur nützlich, sondern sagt auch etwas über die Welt. 

Beispiel: Ein rostiger Schlüssel öffnet nicht nur eine Tür, sondern steht zugleich für eine verdrängte Vergangenheit, zu der man erst Zugang finden muss.

\textbf{Allegorie}: Ein Abschnitt des Spiels kann gleichzeitig wörtlich und übertragen gelesen werden. 

Beispiel: Eine Figur steigt durch einen dunklen Tunnel immer tiefer hinab. Das ist einerseits ein Weg durch den Raum, andererseits kann es auch als Bild für Angst, Erinnerung oder Selbstverlust gelesen werden.

\textbf{Synekdoche (pars pro toto)}: Ein kleines Detail lässt auf ein größeres Ganzes schließen. 

Beispiel: In einem verlassenen Zimmer liegt nur eine halb verbrannte Liste mit Namen. Daraus wird klar, dass hier nicht nur eine Person verschwunden ist, sondern dass es ein größeres, verborgenes Geschehen gibt.

\textbf{Personifikation}: Die Welt wirkt so, als würde sie selbst handeln oder sprechen. Dadurch werden Hinweise nicht nur beschrieben, sondern fast wie Botschaften erfahrbar. 

Beispiel: Eine Tür knarrt nicht einfach, sondern klingt jedes Mal warnend, wenn man sich dem falschen Raum nähert.

Was diese Mittel verbindet, ist ein gemeinsames Prinzip: Bedeutung entsteht nicht durch Eindeutigkeit, sondern durch \textit{Schichtung}. Ein Symbol trägt mehr als seine Funktion, eine Allegorie mehr als ihren wörtlichen Sinn. Für das Rätseldesign bedeutet das: Je mehr Lesarten ein Element zulässt, desto stärker bleibt es im Gedächtnis, vorausgesetzt, die Welt liefert durch Interaktion und Exploration genug Anhaltspunkte, um die Schichten auch tatsächlich freizulegen

Nicht jedes rhetorische Mittel eignet sich gleichermaßen für interaktive Rätsellogik. Die zuvor genannten sind in der Welt selbst verankert, Spielende erschließen ihre Bedeutung durch Exploration und Interaktion, unabhängig von einer sprechenden Instanz. Die folgenden Hyperbel und Euphemismus hingegen setzen eine solche voraus. Sie funktionieren nur, wenn jemand etwas berichtet. Ein NPC, ein Tagebucheintrag, eine Inschrift. Fehlt diese Erzählstimme, wird die Hyperbel entweder als wörtliche Mechanik missverstanden oder als Fülltext ignoriert. Der Euphemismus verliert seine Wirkung, wenn die Täuschung bereits durch Exploration aufgedeckt wurde, bevor sie greifen konnte. Sind NPCs oder eine Narration vorhanden, können beide Mittel gezielt eingesetzt werden, besonders dann, wenn die Erzählinstanz selbst unzuverlässig ist und die Diskrepanz zwischen dem Gesagten und der erfahrbaren Welt zum eigentlichen Rätsel wird.

\textbf{Die Hyperbel} übertreibt eine Aussage so stark, dass sie als Mechanik wörtlich genommen werden muss und kann Erwartungen an die Mechanik erzeugen, die das System dann einlösen muss.

Beispiel: Eine Inschrift warnt: „Wer diese Schwelle überschreitet, verliert alles." Im Spiel stellt sich heraus, der Satz stimmt buchstäblich, jedes Objekt im Inventar verschwindet beim Betreten des Raums.

\textbf{Der Euphemismus}, etwas Bedrohliches wird verharmlosend umschrieben und verschleiert dadurch einen Hinweis. Er kann vor allem dann sinnvoll, wenn eine unzuverlässige Erzählinstanz Teil des Rätsels ist.

Beispiel: Ein Schild neben einer verriegelten Tür trägt die Aufschrift „Mitarbeiterpause, bitte nicht stören". Dahinter hört man Schreie.

\textbf{Unzuverlässige Erzählinstanz}: Die Stimme, die die Welt erklärt, lügt oder irrt sich und die Welt selbst liefert die Gegenbeweise. Sie kann als interpretativ reiches Rätselprinzip eingesetzt werden, setzt aber voraus, dass die Korrektur des verzerrten Bildes durch die Weltlogik selbst ermöglicht wird. Das heißt, die Welt muss so gestaltet sein, dass die Spielenden durch Exploration und Interaktion die Diskrepanz zwischen der Erzählung und der Realität erkennen können. 

Beispiel: Eine KI-Assistenz begleitet die Spielenden durch eine verlassene Forschungsstation und erklärt, alle Experimente seien „planmäßig abgeschlossen worden". Doch zerstörte Labore, versiegelte Notausgänge und unleserlich gemachte Akten erzählen eine andere Geschichte. Wer sie zusammensetzt, erkennt, dass die KI selbst das Unglück ausgelöst hat.

\subsubsection{ÜBUNG: Polysemie}

Entwerfe ein Rätsel, das auf einem einzigen Element basiert, das gleichzeitig als Objekt, Hinweis und Symbol fungiert. 
\begin{itemize}
    \item Beschreibe die verschiedenen Bedeutungsebenen dieses Elements und wie sie miteinander verknüpft sind. 
    \item Achte darauf, dass die Welt genügend Kontext bietet.
    \item Erstelle eine klare,  nicht triviale Lösung, die sich aus der Kombination der Bedeutungsebenen ergibt.
    \item Überlege, wie die Offenheit des Elements die Spielerfahrung bereichert, ohne sie zu frustrieren.
\end{itemize}

\section{Rätseldesign}

Rätsel sind keine arbiträren Hindernisse. Gut gestaltete Rätsel folgen einer inneren Logik, die für die Spielenden erschließbar ist, und erzeugen dabei jene Spannung zwischen Wissen und Nicht-Wissen, die Csikszentmihalyi (\textit{Flow}, \citep{Csikszentmihalyi1990}) als Voraussetzung tiefer Konzentration beschreibt: die Aufgabe überfordert nicht, unterfordert aber auch nicht. Graham Nelson (\textit{The Craft of Adventure}, \citep{Nelson1994}) formuliert für die IF ein Fairnessprinzip, das diese Balance normativ erfasst: Alle zur Lösung notwendigen Informationen müssen in der Welt vorhanden, zugänglich und interpretierbar sein. Marc Leblanc, Brian Hunicke und Robert Zubek (\textit{MDA: A Formal Approach to Game Design}, \citep{Hunicke2004}) ergänzen die Designperspektive um eine systematische Beschreibung der Wechselwirkungen zwischen Regeln, Systemmechanik und Spielerfahrung.\textit{ Rätseldesign ist demnach immer auch Weltdesign}. Ein Rätsel überzeugt nicht als isolierte Aufgabe, sondern als \textit{lesbare Konsequenz der Weltregeln}, in die es eingebettet ist.

\subsection{Spielmechanische Techniken}

Spielmechaniken sind die operativen Einheiten, durch die eine Spielwelt mit den Spielenden interagiert. Sie definieren, was getan werden kann, unter welchen Bedingungen eine Handlung eine Wirkung hat und welche Zustände erreichbar sind. In der IF ist das Repertoire solcher Mechaniken durch die Sprachlichkeit der Interaktion geprägt. Aktionen sind Sätze, Bedingungen sind Weltregeln, Konsequenzen sind Texte. Die folgenden Abschnitte unterscheiden zunächst grundlegende Typen spielerischer Herausforderungen, die quer durch Genres auftreten, und beschreiben dann die spezifischen Verrätselungstechniken, die für das Medium der IF besonders charakteristisch sind.

\subsubsection{Spielerische Herausforderungen}

\textbf{epistemisch}: Unsicherheit über Fakten der Spielwelt.  \\
\textit{Beispiel: Ist der Zeuge glaubwürdig, oder lügt er?}

\textbf{inferentiell}: Eine verborgene Regel muss aus Beobachtungen erschlossen werden.\\
\textit{Beispiel: Alle roten Türen öffnen sich mit Metallgegenständen, warum?}

\textbf{operativ}: Eine bekannte Regel muss korrekt ausgeführt werden.\\
\textit{Beispiel: Den richtigen Code in der richtigen Reihenfolge eingeben.}

\textbf{strategisch}: Teilprobleme müssen in der richtigen Reihenfolge gelöst werden. \\
\textit{Beispiel: Erst Wasser holen, dann Feuer löschen, dann Raum betreten.}

\textbf{interpretativ}: Die gefundene Lösung muss innerhalb der Storyworld gedeutet werden. \\
\textit{Beispiel: Das Muster ist gelöst, aber was bedeutet es für die Geschichte?}

\textbf{parser-linguistisch}: Die Formulierung der Eingabe ist selbst das Problem. Der Spielende kennt die Lösung, findet aber nicht die Syntax, die das System versteht. Es benötigt genaues Lesen der Weltbeschreibungen, weil die korrekten Verben und Objekte dort eingebettet sind. \\
\textit{Beispiel: \texttt{UNTERSUCHE INSCHRIFT} schlägt fehl, \texttt{LIES INSCHRIFT} öffnet eine neue Hinweisebene.}

\subsection{Verrätselungstechniken}

\textbf{Mehrfachcodierung}: Ein Element funktioniert gleichzeitig als Hinweis und als Bedeutungsträger.  \\
\textit{Beispiel: Ein Gemälde zeigt die Lösung, und erzählt zugleich die Vorgeschichte.}

\textbf{Diegetische Hinweise}: Die Lösung ist in Weltzeichen eingebettet, nicht als externe Hilfe markiert. \\
\textit{Beispiel: Der Code steht auf einem alten Grabstein, nicht in einem Notizbuch.}

\textbf{Fragmentierung}: Relevante Informationen sind verteilt und müssen rekonstruiert werden. \\
\textit{Beispiel: Vier Seiten eines zerrissenen Briefes liegen in verschiedenen Räumen.}

\textbf{Metarätsel}: Einzelrätsel ergeben erst zusammen ein übergeordnetes Muster. \\
\textit{Beispiel: Jede Lösung liefert einen Buchstaben, gemeinsam bilden sie ein Lösungswort.}

\textbf{Perspektivwechsel}: Bekannte Informationen werden durch neuen Kontext anders lesbar. \\
\textit{Beispiel: Die scheinbar unwichtige Zahl vom Anfang ist plötzlich der gesuchte Code.}

\textbf{Epistemischer Dunkelraum}: Der Spielende kennt weder den Umfang der Welt noch die Liste möglicher Aktionen. Das Nicht-Wissen darüber, was existiert und was getan werden kann, ist selbst der primäre Rätselmechanismus der IF. Anders als in grafischen Interfaces gibt es kein UI, das Optionen offenlegt. \\
\textit{Beispiel: Ein Gegenstand im Inventar hat eine Funktion, die erst relevant wird, wenn man den richtigen Raum betritt, den man noch nicht kennt.}

\textbf{Inventarrätsel}: Das Inventar ist ein kombinatorischer Zustandsraum. Gegenstände haben latente Verwendungen, die nicht durch ihre Beschreibung, sondern nur durch Weltkenntnis und Kombination aktiviert werden. Nicht der Gegenstand allein, sondern seine Relation zu einem Weltzustand löst das Rätsel. \\
\textit{Beispiel: Eine leere Flasche scheint nutzlos. Erst wenn der Spielende einen Brunnen im Garten findet, ergibt sich die Kombination: Flasche füllen, Wasser auf das brennende Teil schütten, das die Kellertür blockiert.}

\subsection{Allgemeine Rätselmechaniken}

\textbf{Lock-and-key}: Fortschritt ist an ein Objekt, Wissen oder einen Weltzustand gebunden. \\
\textit{Beispiel: Die Tür öffnet sich nur mit dem Schlüssel aus Raum 3.}

\textbf{Zustandsrätsel}: Eine bestimmte Abfolge von Aktionen verändert die Welt lösungsrelevant. \\
\textit{Beispiel: Hebel A, dann B, dann A ergibt die richtige Brückenkonfiguration.}

\textbf{Kombinationsrätsel}: Verteilte Hinweise oder Objekte müssen korrekt zusammengeführt werden. \\
\textit{Beispiel: Drei Fragmente ergeben zusammen die Karte zum Versteck.}

\textbf{Navigationsrätsel}: Die Raumstruktur selbst ist das Problem. \\
\textit{Beispiel: Ein Labyrinth, das nur rückwärts durchquerbar ist.}

\textbf{Sozial-/Dialogrätsel}: Informationen und Zugänge hängen von Gesprächslogik, Rollen oder Vertrauen ab. \\
\textit{Beispiel: Der Wächter redet nur, wenn man sich als Botschafter ausgibt.}

\subsection{Typische Anti-Pattern}

\textbf{Guess-the-verb}: Das Scheitern liegt am fehlenden Parserwort, nicht an der Aufgabe, eine parser-linguistisch Verrätselung ohne ausreichend Hinweise auf  das richtige Wort.\\
\textit{Beispiel: »BENUTZE Seil« funktioniert, »NIMM Seil« nicht – ohne erkennbaren Grund.}

\textbf{Moon logic}: Die Lösung ist technisch möglich, aber nicht herleitbar. \\
\textit{Beispiel: Den Frosch mit der Geige kombinieren, weil beide »schief klingen«.}

\textbf{Hidden Softlock}: Eine frühe, unmarkierte Fehlentscheidung blockiert später die Lösung. \\
\textit{Beispiel: Ein weggeworfenes Item wird erst drei Akte später benötigt.}

\textbf{Fake Agency}: Entscheidungen wirken relevant, haben aber keine echte Auswirkung. \\
\textit{Beispiel: Zwei Dialogoptionen führen zum identischen Ergebnis.}

\section{Fairness und Schwierigkeitsgrad}

Was ein Rätsel fair macht, ist in der IF-Gemeinschaft intensiv diskutiert worden. Die einflussreichsten Fairness-Regeln stammen von Graham Nelson (\textit{The Craft of Adventure}, \citep{Nelson1994}), der auf Basis der Infocom-Tradition und eigener Arbeit am Inform-System ein normatives Regelwerk formuliert: Ein Rätsel ist fair, wenn alle notwendigen Informationen vor dem Lösungsmoment zugänglich waren, wenn seine Lösung der inneren Logik der Welt entspricht, wenn kein Fehlversuch den Spielenden in einen unlösbaren Zustand führt (\textit{unwinnable state}) und wenn der Spielende über Konsequenzen informiert wird, bevor sie irreversibel werden. Ernest Adams (\textit{Fundamentals of Game Design}, \citep{Adams2010}) ergänzt dies spielwissenschaftlich: Fairness ist eine Funktion von Informationsparität. Das System darf nicht mehr wissen, als es dem Spielenden zugänglich gemacht hat.

\subsection{Moon logic vermeiden}

„Moon logic“ ist ein Begriff aus dem Spiele‑ und Puzzledesign und beschreibt Lösungen, die scheinbar „vom Mond“ kommen, also extrem obskur, unplausibel oder intuitiv schwer nachvollziehbar sind. Die Lösung ist nicht aus den gegebenen Hinweisen oder der Spielwelt sinnvoll ableitbar, sondern erfordert völlig unerwartete oder abstrakte Gedankensprünge, beispielsweise durch Insider‑Humor oder außerhalb der Spielwelt liegende Sprachspiele.

In Monkey Island (insbesondere 1 und 2) sind viele solcher „Moon‑Logic‑Puzzles“ stilistisch und thematisch bewusst als insider Witz, Wortspiel oder absurder Gag angelegt. Beispiele sind in Monkey Island 2 die \textit{Monkey‑Wrench‑Lösung} oder das \textit{Gewichts‑Puzzle im Aufzug}. Sie sind nicht „schlecht gemacht“, sondern eher witzige Übertreibung, die erst im Nachhinein als \textit{Moon Logic} wahrgenommen wird. Sie zählt zu den sogenannten Anti-Design-Pattern, also einem Entwurfsmuster, das man meiden sollte.

Wie in vielen klassischen Adventures gibt es auch in Monkey Island Stellen, an denen ein Gegenstand nur dann funktioniert, wenn man ihn zu einem bestimmten Zeitpunkt noch nicht benutzt hat. Das ist Moon‑Logic‑typisch, weil die Lösung nicht aus der Weltlogik, sondern aus einem „System‑Bug im Spiel‑Design“ folgt: Man muss quasi im Vorhinein wissen, was später wichtig ist.\\

Nachdem wir nun wissen, was eine interessante Storyworld und gute Rätsel ausmacht, müssen ihre Logik und ihre Zustände so präzise beschrieben werden, dass ein Computer sie verarbeiten kann. Damit beschäftigen wir uns im nächsten Kapitel.

\textbf{AUFGABE: Possible Worlds}

Fertigstellung unserer Welt, damit wir sie in Logik und dann in Prolog übersetzen können.

\begin{itemize}
    \item Führt die entwickelte Inhalt aus den vorherigen Aufgaben zu einen Storyworld Board zusammen
    \item Spielkonzept, Setting
    \item Ein Diagramm aller Räume und deren Verbindungen.
    \item Akteure und Entities, AFGNCAAP, NPC und mögliche Dialoge
    \item Startinventar     
    \item Interaktive Objekte und ihre Beschreibung, Position, Zustände und Mechanismen
    \item Spielmechanik, Rätsel-Logik und Wenn-Dann“-Bedingungen, Zustände, Regeln
    \item mögliche Aktionen (GEHE, ÖFFNE) und Synonyme
    \item Dokumentation, Texte, ASCII Bilder
    \item Walkthrough: Eine Liste der Befehle, die direkt zum Ziel führt
\end{itemize}

\chapter{Logik}

Die Grundlage für die Programmierung mit Prolog ist eine bestimmte Form der Logik, Prädikatenlogik erster Ordnung, die verglichen zu der Urform der Logik, die auf die frühen griechischen Philosophen zurückgeht, noch recht jung ist. Dieses Kapitel führt in die Elemente der formalen Logik und speziell der Prädikatenlogik, ihr Geschichte und Abgrenzung zu anderen Stufen der Logik, ein. Dabei betrachten wir vor allem auch, wie wir sie zur Formulierung der inneren Logik unserer Story World, Spielmechaniken und Rätsel verwenden können, die wir dann im nächsten Schritt in die Programmiersprache Prolog übersetzen und programmieren. 

Logik ist eine Familie von Denkweisen, die sich darin unterscheiden, wie \textit{sicher} ihre \textit{Schlüsse} sind:\\
\begin{table}[h]
\centering
\begin{tabular}{>{\raggedright\arraybackslash}p{2.2cm} >{\raggedright\arraybackslash}p{3.5cm} >{\raggedright\arraybackslash}p{4.5cm} >{\raggedright\arraybackslash}p{4cm}}
\toprule
\textbf{Verfahren} & \textbf{Richtung} & \textbf{Beispiel} & \textbf{Logischer Rahmen} \\
\midrule
Deduktion &
  Regel + Fall $\rightarrow$ Ergebnis &
  \textit{Alle Räume haben einen Ausgang. Der Keller ist ein Raum.
  $\rightarrow$ Der Keller hat einen Ausgang.} &
  Klassische Logik, monoton, in Prolog nativ \\
\addlinespace
Induktion &
  Fälle $\rightarrow$ Regel &
  \textit{Keller, Wohnzimmer, Dachboden haben alle einen Ausgang.
  $\rightarrow$ Wahrscheinlich haben alle Räume einen Ausgang.} &
  Nicht-monoton, revidierbar \\
\addlinespace
Abduktion &
  Ergebnis + Regel $\rightarrow$ plausibler Fall &
  \textit{Der Keller hat einen Ausgang. Alle Räume haben Ausgänge.
  $\rightarrow$ Der Keller ist vermutlich ein Raum.} &
  Nicht-monoton, hypothetisch \\
\bottomrule
\end{tabular}
\caption{Schlussverfahren im Vergleich.}
\end{table}
\\
Nur die Deduktion ist in der klassischen formalen Logik, wie wir sie hier kennenlernen werden, direkt verankert. Sie ist \textit{wahrheitserhaltend} und \textit{monoton}, also dass hinzukommendes Wissen bereits gezogene Schlüsse nicht ungültig machen kann. Induktion und Abduktion sind dies nicht, denn sie benötigen \textit{nicht-monotone Erweiterungen}, weil ihre Schlüsse \textit{revidierbar} sind, die wir am Ende des Kapitels kurz kennenlernen werden. 

\subsubsection{ÜBUNG: Sherlock Holmes }

Für welches Verfahren ist Sherlock Holmes bekannt und welches wendet er tatsächlich an?

\section{Strukturen des logischen Denkens} 

Wir verwenden Begriffe der Logik fast täglich, doch was meinen wir eigentlich, wenn wir sagen:  ’Das ist doch wohl logisch!’, wenn eine Tatsache, oder Aussage selbstverständlich ist, oder ‘Das ist total unlogisch‘, wenn etwas nicht verstanden wird. Und wie unterscheiden sich die mathematische Logik und die Alltagslogik voneinander?

\subsection{Informelle Logik, oder die Logik des Alltags}

Die \textit{informelle} oder \textit{natürliche} Logik ist die Form des Schließens, die wir im Alltag verwenden. Sie beruht nicht primär auf formaler Gültigkeit, sondern auf Plausibilität, Erfahrungswissen, Situationsdeutung und sozial geteilten Bewertungsmaßstäben. Manfred Kienpointner zeigt in \textit{Alltagslogik} \citep{Kienpointner1992}, dass Argumente im Alltag oft nicht streng nach formaler Logik funktionieren, sondern davon abhängen, was im jeweiligen Kontext plausibel erscheint. Stephen Toulmin erklärt in \textit{The Uses of Argument} \citep{Toulmin1958}, dass Alltagsargumente meist aus drei Teilen bestehen: Daten, einer Schlussregel und einer Behauptung. Wie überzeugend sie sind, hängt vor allem davon ab, ob die Schlussregel im jeweiligen Zusammenhang akzeptiert wird. Douglas Walton \citep{Walton1998} ergänzt, dass viele scheinbare Fehlschlüsse nur im Kontext richtig beurteilt werden können. Dieselbe Argumentstruktur kann je nach Situation falsch oder durchaus sinnvoll sein. 

Beispiel: 
\begin{displayquote}
»Vor dem Restaurant stehen viele Menschen an, also muss das Essen gut sein.«
\end{displayquote}
Formal logisch ist dies ein Fehlschluss (Argumentum ad populum), alltagslogisch aber oft eine funktionale Heuristik zur Reduktion von Unsicherheit und Entscheidungsaufwand.

Dieser Unterschied lässt sich mit Daniel Kahnemans Unterscheidung zwischen \textit{System 1} und \textit{System 2} erklären \citep{Kahneman2011}: \textit{System~1} ist das schnelle, automatische Denken, es erkennt Muster, reagiert auf Erfahrung und trifft Entscheidungen ohne bewusste Anstrengung, arbeitet dabei aber nicht mit strikten Wahrheitswerten, sondern mit graduellen Plausibilitäten. \textit{System~2} ist das langsame, bewusste Denken, es prüft Schritt für Schritt, wendet Regeln an und kontrolliert, ob eine Schlussfolgerung tatsächlich gültig ist. Im Alltag dominiert System~1, formale Logik und Programmierung eher System~2. 

In der informellen Logik reichen implizite Annahmen und pragmatische Plausibilität häufig aus. In der formalen Logik und in der Logikprogrammierung müssen Annahmen \textit{explizit als Fakten} oder\textit{ Regeln} repräsentiert werden, mit denen man rechnen kann.

Für Interactive Fiction bedeutet das methodisch, dass Weltzustände, Raumrelationen, Objekteigenschaften, NPC-Verhalten und Rätsellogik nicht nur intuitiv stimmig sein sollten, sondern müssen als\textit{ konsistentes System formulierbar} sein. Die Übersetzung von alltagslogischem in formallogisches Denken ist damit kein Verlust an Kreativität, sondern die Voraussetzung dafür, narrative Komplexität in ausführbare, überprüfbare Regeln zu überführen.

\subsection{Grenzen der Logik}

Das bedeutet auch, dass bestimmte erzählerische Formen an die Grenzen dieser Übersetzbarkeit stoßen. Bewusstseinsströme (Ulysses, James Joyce), abstrakte Lyrik (Das Ästhetische Wiesel, Christian Morgenstern), der (extrem) unzuverlässige Erzähler (The Tell-Tale Heart, Edgar Allan Poe.),  fatalistische Handlungslosigkeit (König Ödipus, Sophokles) oder simultane Multiperspektivität (Mrs. Dalloway, Virginia Woolf ) lassen sich nur schwer mit den Grundprinzipien von Interactive Fiction, also der Kausalität und Handlungsfähigkeit vereinbaren. Was in der Literatur durch Assoziation, Klang oder Unveränderlichkeit wirkt, kollidiert mit dem impliziten Versprechen der Eingabezeile: "Du kannst handeln und die Welt verändern und die Welt reagiert konsistent." 

Dies kann aber auch als Herausforderung begriffen werden und der bewusste Einsatz solcher Formen in Interactive Fiction zum experimentellen, gestalterischen Mittel werden.

\section{Grundlagen der formalen Logik}

Bevor wir uns mit der technischen Umsetzung digitaler Welten befassen, müssen wir das Fundament verstehen, auf dem diese Welten stehen. Die \textbf{formale Logik}. Als Designer:innen verwenden wir Logik nicht als trockenes Regelwerk der Mathematik, sondern als ein Werkzeug zur \textbf{Strukturierung von Wahrheit}. Es geht darum, die \textit{Komplexität der Welt zu zerlegen}, sie zu \textit{formalisieren} und in \textit{Symbole überführen}, um\textit{ Beziehungen zwischen Dingen} präzise beschreibbar und damit\textit{ anwendbar }zu machen.

\subsection{Die Syllogistik des Aristoteles}

Aristoteles formuliert im 4. Jahrhundert v.\,Chr. das erste systematische Modell gültigen Schließens, den \textbf{Syllogismus}. Ein Syllogismus besteht aus genau \textit{zwei Prämissen} und einer \textit{Konklusion}, wobei alle drei Aussagen eine Beziehung zwischen Kategorien herstellen.

\subsubsection{Das historische Beispiel: Aristoteles´ Ur-Syllogismus}

\textbf{Prämisse A:} Alle Menschen sind sterblich.\\
\textbf{Prämisse B:} Sokrates ist ein Mensch.\\
\textbf{Konklusion:} Also ist Sokrates sterblich.\\

Die Stärke des Syllogismus liegt in seiner \textit{Formgültigkeit}. Ob die Konklusion korrekt ist, hängt \textit{nicht vom Inhalt} ab, sondern allein von der \textit{Struktur}. Wenn die Prämissen wahr sind und die Form stimmt, \textit{muss} die \textit{Konklusion} wahr sein. Aristoteles identifiziert dabei gültige und ungültige Schlussfiguren. Ein frühes Beispiel dafür, dass rationales Denken \textit{Regeln} hat, die sich \textit{systematisch} beschreiben lassen.

Ein Parser-Spiel ist größtenteils nichts anderes als eine riesige Sammlung solcher Syllogismen.

\textbf{Regel:} Wenn ein Objekt verschlossen ist, kann man es nicht öffnen. \\
\textbf{Fakt:} Die Truhe ist verschlossen. \\
\textbf{Logik:} Der Befehl ÖFFNE TRUHE schlägt fehl. \\

Für die spätere Arbeit mit Prolog ist das relevant, weil Prolog genau dieses Prinzip automatisiert. Aus gegebenen Fakten (Prämissen) und Regeln (Schlussformen) leitet das System Konklusionen ab, rein mechanisch, ohne Intuition.

Die Syllogistik hat allerdings klare Grenzen, denn sie kann nur Beziehungen zwischen \textit{Kategorien} ('Mensch') und ihren Eigenschaften ('sterblich') ausdrücken (\textit{Alle X sind Y}), aber keine Aussagen über \textit{Einzelobjekte}, ihre \textit{Eigenschaften (Prädikate)} oder \textit{Relationen} zwischen ihnen. Der Syllogismus tut sich schwer mit Sätzen wie: „Sokrates ist größer als Platon.“, weil 'größer als' eine Relation (Beziehung) zwischen zwei Individuen ist und keine einfache Klassenzugehörigkeit. Genau diese Erweiterung leistet die moderne Logik.

\subsection{Aussagenlogik}

Die Aussagenlogik ist die einfachste Stufe der formalen Logik. Ihr Grundbaustein ist die \textbf{Aussage} (Proposition): ein Satz, der genau einen Wahrheitswert hat, \textit{wahr} oder \textit{falsch}, ohne Zwischenstufen.

\textit{Die Tür ist offen.} $\rightarrow$ wahr oder falsch.\\
\textit{Der Spieler hat den Schlüssel.} $\rightarrow$ wahr oder falsch.\\
\textit{\sout{Vielleicht scheint die Sonne.}} $\rightarrow$ keine Aussage, da kein eindeutiger Wahrheitswert

Solche Aussagen lassen sich mit \textbf{logischen Operatoren} verknüpfen:

\textbf{Konjunktion} (UND, $\land$): \textit{Die Tür ist offen UND der Spieler hat die Lampe.} – Nur wahr, wenn beides gilt.\\
\textbf{Disjunktion} (ODER, $\lor$): \textit{Der Schlüssel liegt im Keller ODER im Garten.} – Wahr, wenn mindestens eines gilt.\\
\textbf{Negation} (NICHT, $\neg$): \textit{Die Tür ist NICHT offen.} – Kehrt den Wahrheitswert um.\\
\textbf{Implikation} (WENN\dots DANN, $\rightarrow$): \textit{WENN der Spieler den Schlüssel hat, DANN kann die Tür geöffnet werden.} Nur falsch, wenn die Bedingung erfüllt ist, aber die Folge nicht eintritt.

Für IF-Design ist die Implikation besonders wichtig: Jede Bedingung im Spiel, \textit{wenn Objekt X im Inventar, dann Aktion Y möglich}, ist eine Implikation. In Prolog wird genau diese Struktur später dann zur Regel.

\subsubsection{ÜBUNG Aussagenlogik}
Formuliere für die folgenden Situationen jeweils die logische Bedingung als Formel. Formalisiere sie, idem du sie in ihre formal logischen Bestandteile zerlegst.

Bestimme zuerst die Teilaussagen (A, B, C,) und verknüpfe sie dann mit den Symbolen ($\land$, $\lor$, $\neg$, $\rightarrow$)  zu  WENN...UND($\land$)...ODER($\lor$)...DANN ($\rightarrow$) und die aussagenlogische Formel (A $\land$ B) $\rightarrow$ C bzw. (A $\lor$ B) $\rightarrow$ C\\

\textbf{Beispielsatz} \\
„Wenn AFGNCAAP im Maschinenraum ist UND der Hebel auf 'An' steht, dann ist die Maschine aktiviert.“\\

\textbf{Logische Teilaussagen}\\
A: AFGNCAAP befindet sich im Maschinenraum. \\
B: Der Hebel hat den Zustand „An“. \\
C: Die Maschine ist aktiviert. \\

FGNCAAP befindet sich im Maschinenraum UND($\land$) der Hebel hat den Zustand „An“ DANN ($\rightarrow$)  Die Maschine ist aktiviert.

\textbf{Aussagenlogische Formel}\\
(A $\land$ B) $\rightarrow$ C

\textbf{Aufgaben:}
\begin{itemize}
    \item Wenn AFGNCAAP die Kerze hält und die Kerze angezündet ist, leuchtet der Altar auf.
    \item Wenn AFGNCAAP im Verlies ist und der Wächter in diesem Verlies ist, spricht der Wächter AFGNCAAP an.
    \item Wenn AFGNCAAP eine Fackel hält und diese Fackel nicht erloschen ist, kann AFGNCAAP die Inschrift lesen.
    \item Wenn AFGNCAAP den Dietrich (Nachschlüssel) hat oder AFGNCAAP ein Brecheisen besitzt, kann AFGNCAAP die Verliestür öffnen.
    \item Überlege 3 eigene Aussagen und versuche sie so zu formalisieren.
    \item Wähle 3 Aussagen aus unserer aktuellen Storyworld und versuche sie in diese Form zu bringen.
\end{itemize}

\subsection{Mit Logik rechnen: Wahrheitswerte}

Der entscheidende Schritt der formalen Logik ist, dass Wahrheitswerte \textit{berechenbar} werden. Jeder logische Operator hat eine eindeutige Definition, die sich als \textbf{Wahrheitstabelle} darstellen lässt. Statt A und B verwenden wir hier $p$ und $q$, die in der Mathematik üblicherweise für Aussagen verwendet werden.

\begin{center}
\begin{tabular}{cc|cccc}
$p$ & $q$ & $p \land q$ & $p \lor q$ & $\neg p$ & $p \rightarrow q$ \\
\hline
w & w & w & w & f & w \\
w & f & f & w & f & f \\
f & w & f & w & w & w \\
f & f & f & f & w & w \\
\end{tabular}
\end{center}

\subsubsection{Leere Wahrheit}
Die letzte Spalte der Tabelle irritiert oft. \textit{WENN falsch, DANN immer wahr}, klingt kontraintuitiv, ist aber formal korrekt. 

Dass die \textbf{Implikation} ($p \rightarrow q$) wahr ist, wenn die Bedingung $p$ falsch ist, wirkt auf unser intuitives Sprachempfinden oft seltsam, weil wir im Alltag meist eine \textit{kausale Verbindung }erwarten. 
In der Logik nennen wir das 'leere Wahrheit' ('vacuous truth').

'Wenn du dein Zimmer aufräumst $p$, dann bekommst du ein Eis $q$.'

Wir prüfen nun die vier Fälle der Wahrheitstabelle anhand der Frage:

\textbf{Sagt der Vater die Wahrheit?}

\begin{itemize}
    \item Zimmer aufgeräumt ($w$), Eis bekommen ($w$): Das Versprechen wurde gehalten ($w$). $w, w \rightarrow w$
    \item Zimmer aufgeräumt ($w$), kein Eis bekommen ($f$): Der Vater hat gelogen. Das Versprechen wurde gebrochen ($f$). $w, w \rightarrow f$
    \item Zimmer \textit{nicht} aufgeräumt ($f$), \textit{trotzdem} ein Eis bekommen ($w$): Der Vater hat nicht gelogen. Er hat nur eine Aussage darüber getroffen, was passiert, wenn aufgeräumt wird. Er hat nicht eingeschränkt, dass es kein Eis gibt wenn der Sohn nicht aufräumt. Die \textit{Regel wurde daher nicht verletzt} ($w$). $f, w \rightarrow w$
    \item Zimmer \textit{nicht} aufgeräumt ($f$), kein Eis bekommen ($f$): Auch hier hat der Vater \textit{nicht gelogen}. Da die \textit{Bedingung nicht erfüllt} wurde, war er zu nichts verpflichtet ($w$). $w, f \rightarrow w$ \\
\end{itemize}

\begin{table}[h] 
    \centering 
    \begin{tabular}{lllp{5cm}} \toprule
    \textbf{p (Aufräumen)} & \textbf{q (Eis)} & \textbf{Implikation} & \textbf{Grund} \\ \midrule 
    w & w & w & Versprechen gehalten \\ 
    w & f & f & Versprechen gebrochen \\ 
    f & w & w & Leere Wahrheit (Bedingung nicht verletzt, da kein Verbot für Eis ohne Aufräumen) \\ 
    f & f & w & Leere Wahrheit (Bedingung für Eis nicht erfüllt) \\ \bottomrule 
    \end{tabular}
\end{table}

\newpage
\subsubsection{Logische Gleichheit}

Zwei Ausdrücke heißen \textbf{logisch gleich äquivalent (gleich)}, wenn sie in jeder Zeile denselben Wahrheitswert haben. Die wichtigsten Äquivalenzen, etwa die \textbf{De Morganschen Regeln} ($\neg(p \land q) \equiv \neg p \lor \neg q$) sind nicht nur mathematisch relevant, sondern helfen beim Umformulieren von Spielbedingungen: \textit{„Nicht beides gleichzeitig"} ist dasselbe wie \textit{„mindestens eines fehlt"}.

Das Prinzip dahinter: Wenn du eine verneinte Klammer auflöst, „kippt“ der logische Operator in sein Gegenteil um:

Aus $\neg(p \land q)$ wird $\neg p \lor \neg q$ \\
Aus \textbf{NICHT (p UND q)} wird \textbf{(NICHT p) ODER (NICHT q)}.\\

Aus $\neg(p \land q)$ wird $\neg p \land \neg q$ \\
Aus \textbf{NICHT (p ODER q)} wird \textbf{(NICHT p) UND (NICHT q)}.\\

\textbf{ÜBUNG Logische Äquivalenz}

\begin{itemize}
    \item Zeige, dass die Aussagen „Nicht (Die Sonne scheint und es regnet)“  äquivalent ist zu „Die Sonne scheint nicht oder es regnet nicht“.
    \item Identifiziere jeweils die beiden Teilaussagen, und weise sie $p$ und $q$ zu. Was passiert wenn ich die Klammer der einen Aussage auflöse bzw. bei der anderen setze?

HINWEIS: Dass die Sonne scheint UND es regnen kann ist kein echter logischer Widerspruch (Kontradiktion), das gaukelt uns unsere Alltagslogik vor. Beide Aussagen beschreiben unterschiedliche meteorologische Zustände, die sich nicht direkt logisch ausschließen. Ein echter Widerspruch würde nur dann vorliegen, wenn eine Aussage und ihre eigene Verneinung gleichzeitig wahr sein sollen, ( $p$  $\land$ $\neg p$) , was bedeuten würde:  'es regnet UND es regnet NICHT'
\end{itemize}

\subsection{Prädikatenlogik}

Die Aussagenlogik kann nur über ganze Sätze urteilen,  sie sieht jedoch nicht die \textit{Struktur} einer Aussage. Während die Aussagenlogik einen Satz nur als Ganzes (wahr oder falsch) betrachtet, erlaubt die Prädikatenlogik (erster Ordnung) eine Analyse seiner Bestandteile. Sie löst diese Einschränkung, indem sie Aussagen  \textit{Objekte}, \textit{Eigenschaften} und \textit{Relationen} zerlegt. Sie kann ausdrücken, dass eine Eigenschaft für \textit{alle} oder \textit{mindestens ein} Objekt gilt. In der Aussagenlogik müsste man für jedes einzelne Objekt (z. B. Designstudierenden) eine eigene Variable einführen, was allgemeine Schlussfolgerungen unmöglich macht. Sie macht sichtbar, dass verschiedene Aussagen dasselbe Objekt oder dieselbe Eigenschaft betreffen und kann Beziehungen zwischen mehreren Objekten präzise formulieren, wie zum Beispiel „x mag y“ oder „x ist größer als y“. Dafür führt sie \textbf{Variablen} ein, wodurch \textbf{Aussageformen} entstehen. Diese sind wie \textit{Schablonen,} die erst durch das\textit{ Einsetzen konkreter Werte} zu einer \textit{prüfbaren Aussage} werden.

Statt \texttt{position(laterne, eingang)} als unteilbare Aussage $p$ zu behandeln, schreiben wir in Prädikatenlogik: \texttt{position} ist ein \textit{Prädikat} (eine Relation), \texttt{laterne} und \texttt{eingang} sind \textit{Terme} (die Objekte, über die gesprochen wird). Diese Notation ist \textit{direkt} die Syntax von Prolog: \texttt{position(laterne, eingang).}

Die Prädikatenlogik führt die \textbf{Variablen} und \textbf{Quantoren} ein:

\begin{itemize}
    \item \textbf{Allquantor} ($\forall$, „für alle"): $\forall X\colon \text{raum}(X) \rightarrow \text{betretbar}(X)$ – Jeder Raum ist betretbar.
    \item \textbf{Existenzquantor} ($\exists$, „es gibt"): $\exists X\colon \text{position}(X, \text{eingang}) \land \text{gegenstand}(X)$ – Es existiert etwas im Eingang, das ein Gegenstand ist.\\
\end{itemize}

\begin{tabular}{lll}
\toprule
\textbf{Aspekt} & \textbf{Aussagenlogik} & \textbf{Prädikatenlogik} \\
\midrule
Analyse & Ganze Sätze (wahr/falsch) & Interne Struktur (Objekte, Relationen) \\

Generalisierung & Jede Aussage einzeln & $\forall/\exists$ für alle oder einige Objekte \\

Beispiel & ``S1 $\wedge$ S2'' & $\forall x$ (Mensch($x$) $\rightarrow$ Sterblich($x$)) \\
\bottomrule
\end{tabular}
\vspace{1cm}
Die Prädikatenlogik erster Ordnung ist genau die Logik, auf der Prolog basiert.\textit{ Fakten }sind quantorenfreie Aussagen, \textit{Regeln }sind Implikationen mit\textit{ implizitem Allquantor,} und \textit{Abfragen} sind \textit{existenzquantifizierte} Fragen an die \textit{Wissensbasis}.

\subsection{Mit Quantoren Rechnen}

Quantoren interagieren mit logischen Operatoren nach festen Regeln. Zwei Grundfälle zeigen zunächst, wie sich Quantoren über Verknüpfungen \textit{verteilen}:

$\forall X\colon (P(X) \land Q(X)) \equiv \forall X\colon P(X) \land \forall X\colon Q(X)$\\
„Alle Räume sind dunkel \textit{und} feucht" ist gleichbedeutend damit, dass alle Räume dunkel sind \textit{und} alle Räume feucht sind.

$\exists X\colon (P(X) \lor Q(X)) \equiv \exists X\colon P(X) \lor \exists X\colon Q(X)$\\
„Es gibt einen Raum, der dunkel \textit{oder} feucht ist" ist gleichbedeutend damit, dass es einen dunklen Raum gibt \textit{oder} einen feuchten Raum gibt.

\textbf{Nicht intuitiv:} Der Allquantor verteilt sich \textit{nicht} über $\lor$, und der Existenzquantor verteilt sich \textit{nicht} über $\land$:

$\forall X\colon (P(X) \lor Q(X)) \not\equiv \forall X\colon P(X) \lor \forall X\colon Q(X)$\\
„Alle Räume sind dunkel oder feucht" bedeutet \textit{nicht}, dass alle Räume dunkel sind oder alle Räume feucht sind. Es könnte sein, dass die Hälfte der Räume nur dunkel und die andere Hälfte nur feucht ist. Dann ist die linke Aussage wahr, die rechte aber falsch, weil weder alle Räume dunkel noch alle Räume feucht sind.

$\exists X\colon (P(X) \land Q(X)) \not\equiv \exists X\colon P(X) \land \exists X\colon Q(X)$\\
„Es gibt einen Raum, der dunkel \textit{und} feucht ist" bedeutet \textit{nicht} dasselbe wie „Es gibt einen dunklen Raum und es gibt einen feuchten Raum". Im zweiten Fall könnten das zwei völlig verschiedene Räume sein: einer dunkel, ein anderer feucht, aber keiner von beiden beides zugleich.

\textbf{De Morgansche Regeln für Quantoren.} Dasselbe Prinzip, das wir bei den De Morganschen Regeln der Aussagenlogik gesehen haben ($\neg(p \land q) \equiv \neg p \lor \neg q$), gilt analog für Quantoren: Negation verwandelt jeden Quantor in sein Gegenteil.

$\neg \forall X\colon P(X) \equiv \exists X\colon \neg P(X)$\\
„Nicht alle Räume sind dunkel" bedeutet: „Es gibt mindestens einen Raum, der nicht dunkel ist."

$\neg \exists X\colon P(X) \equiv \forall X\colon \neg P(X)$\\
„Es gibt keinen Schlüssel" bedeutet: „Für jedes Objekt gilt: Es ist kein Schlüssel."

\textbf{Reihenfolge der Quantoren.} Die Reihenfolge ist nicht beliebig: $\forall X \exists Y\colon \text{verbindet}(X, Y)$ (jeder Raum hat mindestens einen Ausgang) ist etwas anderes als $\exists Y \forall X\colon \text{verbindet}(X, Y)$ (es gibt einen Raum, der von überall erreichbar ist). Solche Unterschiede sind beim Modellieren einer IF-Welt direkt relevant.

\textbf{Verbindung zu Prolog.} Der Allquantor wird implizit durch Variablen in Regeln ausgedrückt:

\texttt{gefaehrlich(X) :- raum(X), dunkel(X).}

bedeutet: \textit{Für jedes X, das ein Raum und dunkel ist, gilt: X ist gefährlich.} Der Existenzquantor entspricht der Abfrage:

\texttt{?- position(X, eingang).} fragt: \textit{Gibt es ein X im Eingang?}

Prolog kennt keine echte logische Negation, sondern nur die \textbf{Negation als Fehlschlag} (\textit{Negation as Failure}, NAF), ausgedrückt durch \texttt{\textbackslash+}. 

Sie basiert auf der \textbf{Closed-World-Annahme}: Was nicht beweisbar ist, gilt als falsch. 

\texttt{\textbackslash+ position(schluessel, \_)} bedeutet nicht logisch „Es gibt keinen Schlüssel", sondern: „Prolog kann keinen Ort für den Schlüssel beweisen." Das entspricht der zweiten De Morganschen Regel ($\neg \exists \equiv \forall \neg$) nur dann korrekt, wenn die Wissensbasis vollständig ist.

\subsubsection{ÜBUNG: Quantoren und De Morgansche Regeln}

\begin{itemize}
    \item Bestimme für jede der folgenden Aussagen, ob $\forall$ oder $\exists$ verwendet wird, und bilde jeweils die De Morgansche Gegenaussage:
    \begin{itemize}
        \item „Alle Türen im Schloss sind verschlossen."
        \item „Es gibt einen Raum, der keinen Ausgang hat."
        \item „Kein Gegenstand im Inventar ist ein Schlüssel."
    \end{itemize}
    \item Welche dieser Aussagen lässt sich direkt als Prolog-Regel formulieren, welche als Abfrage mit \texttt{\textbackslash+}? Schreibe die Prolog-Entsprechungen auf.
    \item Wähle zwei Zustände aus deiner eigenen Storyworld und formuliere sie einmal mit $\forall$ und einmal mit $\exists$. Was ändert sich in der Bedeutung?
\end{itemize}

\subsection{Logik höherer Ordnung}

Die Prädikatenlogik \textit{erster} Ordnung quantifiziert über Objekte: Räume, Schlüssel, Figuren. Logik \textit{höherer} Ordnung erlaubt darüber hinaus, über Prädikate selbst zu quantifizieren, also über ihre Eigenschaften und Relationen. 

Ein Satz wie: \textit{„Es gibt eine Eigenschaft, die alle Rätsel gemeinsam haben"} ($\exists P\colon \forall X\colon \text{raetsel}(X) \rightarrow P(X)$) ist erst in höherer Ordnung formulierbar.

Für die Praxis im Kurs ist das nicht unmittelbar relevant: Prolog arbeitet standardmäßig in der Prädikatenlogik erster Ordnung, und die allermeisten IF-Mechaniken lassen sich darin vollständig beschreiben. 

Der Hinweis ist dennoch wichtig, weil er die \textit{Grenze} sichtbar macht: Was in der Prädikatenlogik erster Ordnung nicht ausdrückbar ist, wie etwa Aussagen über Regelstrukturen selbst, berührt das Gebiet der \textit{Metaprogrammierung,} das wir im Prolog-Kapitel noch aufgreifen werden.

\subsection{Logik des Zufalls}

Die formale Logik ist \textit{deterministisch}: Aus gegebenen Fakten und Regeln folgt\textit{ immer dasselbe Ergebnis}. Aussagen sind entweder \textit{wahr} oder \textit{falsch}, es existiert kein Dazwischen. In der Praxis, etwa im interaktiven Design oder in Spielen, wollen wir jedoch \textit{Zufallselemente }einbauen, ein NPC erscheint mit 50\,\% Wahrscheinlichkeit, ein Rätsel wird zufällig generiert. Das lässt sich mit klassischer Logik allein nicht ausdrücken und erfordert Erweiterungen wie die probabilistische Logik.


\subsection{Probabilistische Logik}

Probabilistische Logik erweitert die klassische Logik, indem sie Aussagen nicht als absolut wahr oder falsch bewertet, sondern mit Wahrscheinlichkeitswerten zwischen 0 und 1 versieht. Sie basiert häufig auf dem \textbf{Satz von Bayes} und wird in KI, Statistik und generativen Systemen eingesetzt.

\textbf{Aussagenlogik:} Verknüpft atomare Aussagen mit Operatoren wie UND, ODER, NICHT – immer binär. Beispiel: „Es regnet oder es schneit."

\textbf{Prädikatenlogik:} Ergänzt die Aussagenlogik um Prädikate und Quantoren ($\forall$, $\exists$), um Aussagen über Individuen zu machen. Beispiel: $\forall x: \text{Vogel}(x) \rightarrow \text{fliegt}(x)$.

\textbf{Probabilistische Erweiterung:} Beide Varianten können um Wahrscheinlichkeitsgrade ergänzt werden. Beispiel: $P(\text{Regen}) = 0{,}7$ statt nur „Es regnet." Für Prädikate: $P(\text{fliegt}(x) \mid \text{Vogel}(x)) = 0{,}9$.

\begin{center}
\begin{tabular}{lll}
\toprule
\textbf{Eigenschaft} & \textbf{Klassische Logik} & \textbf{Probabilistische Logik} \\
\midrule
Wahrheitswerte & $\{0, 1\}$ -- binär & Spektrum $[0, 1]$ \\
Unsicherheit & Nicht abbildbar & Ja, via Wahrscheinlichkeit \\
Inferenz & Deduktiv, monoton & Bayesianisch -- Schlüsse werden bei neuer Evidenz \textit{aktualisiert}, nicht aufgehoben \\
Anwendung & Exakte Beweise & KI, Statistik, generatives Design \\
\bottomrule
\end{tabular}
\end{center}

\subsection{Zufallsvariablen}

Eine \textit{Zufallsvariable} ist eine messbare Funktion $X: \Omega \to \mathbb{R}$ auf einem Wahrscheinlichkeitsraum $(\Omega, \Sigma, P)$, die jedem möglichen Versuchsausgang $\omega \in \Omega$ einen reellen Wert zuordnet. Sie modelliert zufällige Ereignisse und kann \textbf{diskret} (z.\,B.\ Würfelwurf: $X \in \{1,\ldots,6\}$, je mit $P(X=k) = \tfrac{1}{6}$) oder \textit{stetig} sein (z.\,B.\ Pixelhelligkeit in $[0, 255]$).

In der \textit{klassischen Logik} gibt es keine Zufallsvariablen, alles ist deterministisch binär. In der probabilistischen Logik ermöglichen sie, Aussagen Wahrscheinlichkeitsverteilungen zuzuweisen, etwa:
\[
    P(\text{NPC erscheint}) = X(\omega) = 0{,}5
\]
Dies bildet die Grundlage für \textit{Bayessches Schließen}\textbf{ }und für generative Systeme im Design, bei denen Ausgaben nicht fest definiert, sondern wahrscheinlichkeitsgesteuert erzeugt werden.

\subsection{Monotone und nicht-monotone Logik}

\textbf{Monotone Logik} bedeutet, dass abgeleitete Schlüsse bei Hinzufügen neuer Annahmen erhalten bleiben, sie können nur wachsen, nie zurückgenommen werden. 

Formal: $\Gamma \vdash A \Rightarrow \Gamma \cup \{B\} \vdash A$. 

Bedeutet: "Wenn aus einer Menge von Voraussetzungen $\Gamma$ eine Aussage $A$ folgt, dann folgt diese Aussage $A$ auch noch, nachdem eine weitere Annahme $B$ zur Menge der Voraussetzungen hinzugefügt wurde."

Klassische formale Logik ist stets monoton und typisch für mathematische Beweise.

\textbf{Nicht-monotone Logik} erlaubt hingegen, dass neue Informationen frühere Schlüsse ungültig machen. Dies entspricht dem Alltagsdenken und der Arbeit mit Standardannahmen (sog.\ \textit{Defaults}). Klassisches Beispiel: Die Annahme „Tweety fliegt" gilt zunächst, weil Tweety ein Vogel ist, wird aber revidiert, sobald bekannt wird, dass Tweety ein Pinguin ist. Dies wird formal oft über \textit{Default-Regeln} realisiert, wie etwa:
\[ \text{Fliegt}(X) \leftarrow \text{Vogel}(X) \land \neg \text{Abnormal}(X) \]
Hierbei bezeichnet $\neg \text{Abnormal}(X)$ eine Bedingung, die als erfüllt gilt, solange kein Gegenbeweis vorliegt. 

Wird das System um das Wissen $\text{Pinguin}(X) \rightarrow \text{Abnormal}(X)$ erweitert, führt dies dazu, dass der ursprüngliche Schluss „Tweety fliegt“ bei bekanntem Pinguin-Status automatisch entfällt. Dieser Mechanismus spiegelt menschliches\textit{ Alltagsdenken} wider, bei dem Vorannahmen so lange gelten, bis spezifische Ausnahmen oder neue Fakten ihre Korrektur erzwingen.

\begin{center}
\begin{tabular}{lll}
\toprule
\textbf{Eigenschaft} & \textbf{Monoton} & \textbf{Nicht-monoton} \\
\midrule
Neue Annahmen & Erweitern Schlüsse nur & Können Schlüsse zurücknehmen \\
Beispiel & Mathematische Beweise & „Tweety fliegt" $\to$ Pinguin-Revision \\
Anwendung & Deduktion & KI, Defaults, Alltagsdenken \\
\bottomrule
\end{tabular}
\end{center}

\subsection{Epistemische Einordnung der Schlussverfahren}

Die logischen Systeme, die wir zu Beginn kennengelernt haben (Deduktion, Induktion und Abduktion) lassen sich nun wie folgt charakterisieren: Während die \textit{Deduktion als monotoner Prozess innerhalb der klassischen Logik wahrheitserhaltend operiert,} sind Induktion und Abduktion \textit{erweiternde} (ampliative) Verfahren. Sie \textit{erweitern} den \textit{Wissensbestand über den Informationsgehalt der Prämissen hinaus}, erfordern jedoch aufgrund ihrer probabilistischen oder auf \textit{Defaults} basierenden Natur einen nicht-monotonen Rahmen, um \textit{Revisionen} bei widersprüchlichen \textit{neuen Informationen} zu ermöglichen.
\begin{table}[h]
\centering
\begin{tabular}{lll}
\toprule
\textbf{Verfahren} & \textbf{Logischer Typ} & \textbf{Charakteristik} \\
\midrule
Deduktion & Monoton & Wahrheitserhaltend, sicher \\
Induktion & Nicht-monoton & Verallgemeinernd, revidierbar \\
Abduktion & Nicht-monoton & Hypothesenbildung, revidierbar \\
\bottomrule
\end{tabular}
\caption{Logische Schlussverfahren und ihre Monotonie-Eigenschaften.}
\end{table}
Standard Prolog, das wir in diesem Kurs nutzen,  ist grundsätzlich ein \textit{deduktives System} und leitet als Programmiersprache aus \textit{Fakten} und \textit{Regeln} \textit{automatisiert }Schlüsse ab, führt jedoch weder Induktion noch Abduktion (oder nur mit Hilfe von  Erweiterungen) durch. Das werden wir uns nun im nächsten Kapitel  genauer ansehen.

\subsection{Mit der Tupelnotation die Welt beschreiben}

Die Tupelnotation dient nicht nur der abstrakten Beschreibung eines Welt- oder Rätselmodells, sondern auch als Zwischenschritt zur formalen Logik. Ein Tupel ordnet die zentralen Bestandteile eines Systems, die Prädikatenlogik macht diese Bestandteile und ihre Relationen als Aussagen formulierbar. Sie hilft uns also bei der Formalisierung unserer Storyworld und ihrer Mechaniken und Zustände und wir sehen so auch früh, ob es überhaupt formalisierbar ist und implementiert werden kann.

\textbf{Kompakte Gesamtbeschreibung:} Ein einziges Symbol ($G$) steht für das gesamte Modell. Statt eine Welt in Prosa zu beschreiben, lässt sie sich als mathematisches Objekt handhaben und benennen.

\textbf{Klare Komponentenstruktur:} Statisches und Dynamisches wird explizit getrennt. Die Topologie der Welt (Räume, Verbindungen, Bedingungen) ist von ihrem aktuellen Zustand (was gerade gilt) klar unterschieden.

\textbf{Brücke zur Implementierung:} Jede Komponente des Tupels entspricht direkt einer Gruppe von Prolog-Prädikaten. Der Tupel ist damit nicht nur Denkwerkzeug, sondern Entwurfsplan für das Programm.

\textbf{Die vier Leitfragen sind immer:}

Das Gesamtmodell $G$ für ein Raummodell.

\begin{enumerate}
    \item Was sind die Objekte (Räume)? → $V$ (Vertices)
    \item Wie hängen sie zusammen? → $E$ (Edges)
    \item Was schränkt diese Beziehungen ein? → $W$ (Weights)
    \item Was kann sich zur Laufzeit ändern? → $S$ (State)
\end{enumerate}

Bei der Wahl der Symbole (Buchstaben) gibt es lose Konventionen aus der Mathematik und Informatik. Domänenübergreifende Konventionen gibt es nicht und Lesbarkeit schlägt dabei immer Konvention. Wenn man von den Konventionen abweicht, sollte man es jedoch explizit definieren. $Proc$ für Prozesse ist verständlicher als ein obskures $P$, das mit Wahrscheinlichkeit oder Pfaden verwechselt werden könnte.

\textbf{Etablierte Buchstaben:}

\begin{itemize}
    \item $V$ für \textit{Vertices} (Knoten), $E$ für \textit{Edges} (Kanten) – aus der Graphentheorie, quasi universell
    \item $S$ für \textit{State}/Zustand – aus der Automatentheorie
    \item $\Sigma$ für Alphabet oder Aktionsmenge – aus der formalen Sprachtheorie
    \item $Q$ für Zustände in Automaten, $\delta$ für Übergangsfunktion
    \item $C$ für \textit{Conditions}/\textit{Constraints}
    \item $\Lambda$ oder $L$ für \textit{Labels} (beschriftete Kanten)
\end{itemize}

\vspace{1cm}

Das Schema lässt sich auf beliebige Domänen übertragen. Zur Illustration ein spielfremdes Beispiel — Terminplanung:

$G = (T, D, K, S)$
\begin{itemize}
    \item $T$ = Termine
    \item $D$ = Abhängigkeiten (B erst nach A)
    \item $K$ = Konflikte/Bedingungen (Raum belegt, Person verfügbar)
    \item $S$ = aktuell gebuchte Termine
\end{itemize}

\vspace{1cm}

Für unser Spiel:

Ein graphenbasiertes Weltmodell kann etwa als $G=(V,E,W)$ notiert werden. Dabei ist $G$ das Gesamtmodell, $V$ die Menge der Räume, $E$ die Menge der Verbindungen zwischen Räumen und $W$ die Menge der Bedingungen, unter denen bestimmte Übergänge möglich sind. Dieses Tupel beschreibt die \textit{statische Topologie} der Welt, also welche \textit{Räume} es gibt, welche \textit{Verbindungen} existieren, welche Übergänge \textit{Bedingungen} haben.

Was sich im \textit{Spielverlauf verändert}, also ob der Spieler einen Schlüssel trägt oder eine Tür geöffnet hat, gehört nicht in dieses statische Modell. Dafür wird eine zweite Ebene eingeführt: der \textit{Zustand} $S$. Er hält fest, welche Bedingungen im \textit{aktuellen Spielmoment erfüllt} sind.

Das Gesamtmodell lautet dann:

$$G=(V,E,W,S)$$

Beispielsweise:

$$V=\{eingang,korridor,archiv\}$$
$$E=\{(eingang,korridor),(korridor,archiv)\}$$
$$W=\{(korridor,archiv,tuer\_offen)\}$$
$$S=\{tuer\_offen\}$$

Das bedeutet: Das Modell $G$ besteht aus drei Räumen mit zwei Verbindungen. Der Übergang von \textit{korridor} nach \textit{archiv} ist an die Bedingung \texttt{tuer\_offen} geknüpft. Der aktuelle Zustand $S$ enthält genau diese Bedingung, also ist der Übergang im Moment erfüllt. Die Tupelschreibweise sagt damit noch nicht, dass etwas gilt, sondern welche \textit{Bestandteile} ein Modell hat, und was davon \textit{unveränderlich} ist und was sich \textit{ändern} kann.

Erst in der Prädikatenlogik werden diese Bestandteile als formalisierte Aussagen ausgedrückt. Jedes Element $v \in V$ wird zum Faktum \texttt{raum(v)}, jedes Paar $(x,y) \in E$ zum Faktum \texttt{verbindet(x,y)}, und jedes Tripel $(x,y,z) \in W$ zum Faktum \texttt{bedingt(x,y,z)}.

In Prolog wird ein Fakt dabei automatisch zu einem Teil eines Prädikats, sobald er in der Wissensbasis definiert ist. Ein Prädikat ist in Prolog die Gesamtheit aller Fakten und Regeln mit demselben Funktor (Namen) und derselben Stelligkeit (Anzahl der Argumente). Die drei Fakten \texttt{raum(eingang)}, \texttt{raum(korridor)} und \texttt{raum(archiv)} bilden zusammen das Prädikat \texttt{raum/1}.

Das Prädikat \texttt{erreichbar(X,Y)} ist dagegen kein direkter Bestandteil des Tupels, sondern ein abgeleitetes Prädikat, das aus Verbindungen und Bedingungen folgt. Es wird auf der \textit{Regelebene} als Implikation definiert. Damit wird aus der bloßen Modellbeschreibung ein logisches System. Wir können nun etwa formulieren:

„Wenn eine Tür offen ist und zwei Räume verbindet, ist der Zielraum erreichbar."
$$\forall X\forall Y\,(verbindet(X,Y) \land zustand(tuer, offen) \rightarrow erreichbar(X,Y))$$

Dasselbe Prinzip gilt für die Navigation im Raum. Ob ein Übergang möglich ist, hängt davon ab, welche Bedingungen im aktuellen Zustand $S$ erfüllt sind:

\textit{„Für alle Räume $X$ und $Y$ und jede Bedingung $Z$ gilt: Wenn $X$ mit $Y$ verbunden ist, dieser Übergang die Bedingung $Z$ voraussetzt, und $Z$ im aktuellen Zustand $S$ erfüllt ist, dann ist $Y$ von $X$ aus erreichbar."}
$$\forall X\forall Y\forall Z\,((verbindet(X,Y) \land bedingt(X,Y,Z) \land erfuellt(Z)) \rightarrow erreichbar(X,Y))$$

$S$ kann dabei viele verschiedene Arten von \textit{Zustandsatomen} enthalten. Das Argument eines Prädikats ist immer das \textit{Subjekt} des Fakts: \texttt{inventar(spieler, schluessel)} liest sich als ,,der Spieler hat den Schlüssel``, nicht ,,der Schlüssel hat einen Spieler``. Die Reihenfolge der Argumente trägt die Bedeutung.

Typische IF-Kategorien:

\begin{itemize}
    \item \textbf{Inventar:} \texttt{inventar(spieler, schluessel)}, \texttt{inventar(spieler, laterne)}
    \item \textbf{Zustand von Objekten:} \texttt{zustand(archivtuer, verschlossen)}, \texttt{zustand(archivtuer, offen)}
    \item \textbf{Position:} \texttt{position(spieler, korridor)}, \texttt{position(akte, archiv)}
    \item \textbf{Erlerntes Wissen:} \texttt{wissen(spieler, schlosstechnik)}, \texttt{wissen(spieler, zugangscode)}
    \item \textbf{Ereignisse:} \texttt{untersucht(spieler, aktenschrank)}, \texttt{geloest\_durch(archivtuer, schluessel)}
\end{itemize}

Was nicht in $S$ enthalten ist, gilt als falsch. Das entspricht der sogenannten \textit{Closed-World-Assumption}: Alles, was nicht explizit als wahr vermerkt ist, wird als nicht zutreffend angenommen.

Wir werden sehen, dass in Prolog diese Annahme der Standard ist. Fakten wie \texttt{inventar(spieler, schluessel).} oder \texttt{zustand(archivtuer, offen).} in der Wissensbasis sind damit nichts anderes als die Elemente von $S$ zur Laufzeit des Programms.

Die Tupelnotation ist also die \textit{Entwurfsnotation}, die Prädikatenlogik ihre \textit{Explikation (Präzisierung)} als Menge formaler Aussagen. Auf diese Weise wird aus einem informellen Konzept ein logisch prüfbares Weltmodell.

\subsubsection{ÜBUNG: Bidirektionaler Durchgang}

Formuliere als prädikatenlogische Regel:

\textit{„Ein Durchgang zwischen $X$ und $Y$ ist genau dann bidirektional, wenn $X$ mit $Y$ verbunden ist und $Y$ mit $X$ verbunden ist."}
$$\forall X \forall Y\, \underline{\hspace{2cm}} \leftrightarrow \underline{\hspace{2cm}}$$

Sie beschreibt das Prinzip der Bidirektionalität in der Graphentheorie oder Logik. Eine Verbindung ist genau dann zweiseitig (bidirektional), wenn sie in beide Richtungen existiert.

\subsection{Das STRIPS-Modell, oder Interaktion als Zustandsübergang}

Bisher beschreibt unser System $G=(V,E,W,S)$ einen \textit{statischen Moment} der Spielwelt. Um zu definieren, wie Aktionen (wie das Aufheben eines Objekts oder das Öffnen einer Tür) die Welt verändern, müssen wir unsere Tupelnotation um eine dynamische Komponente erweitern: die menge der Aktionen $\Sigma$. Unser Modell wächst damit organisch zu $G=(V,E,W,S,\Sigma)$.

Dieses formale Vorgehen greift auf das klassische \textbf{STRIPS-Modell} (Stanford Research Institute Problem Solver, \citealp{Fikes1971}) zurück. Es gilt als Meilenstein der Künstlichen Intelligenz zur automatischen Handlungsplanung. Eine Aktion $a \in \Sigma$ wird nicht durch prozeduralen Programmcode beschrieben, sondern rein \textit{deklarativ} durch drei Mengen, die den Zustand $S$ transformieren:

\textbf{$\text{pre}(a)$ (Preconditions):} Vorbedingungen, die erfüllt sein müssen, damit die Aktion ausführbar ist. Sie müssen in der aktuellen Welt wahr sein (mathematisch: $\text{pre}(a) \subseteq S$).

\textbf{$\text{del}(a)$ (Delete List):} Fakten, die durch die Handlung ihre Gültigkeit verlieren und aus dem Weltzustand gelöscht werden.

\textbf{$\text{add}(a)$ (Add List):} Neue Fakten, die durch die Handlung wahr werden und zur Welt hinzugefügt werden.

Mathematisch lässt sich der neue Zustand der Welt $S'$ nach der Ausführung einer Aktion $a$ elegant als reine Mengenoperation ausdrücken:

\begin{equation}
S' = (S \setminus \text{del}(a)) \cup \text{add}(a)
\end{equation}

Ein formalisiertes Minimalbeispiel aus unserem IF-Kontext (Interactive Fiction) verdeutlicht das Prinzip. Die \textit{Aktion} $a$ (\enquote{Nimm den Schlüssel}) lässt sich als Manipulation des Zustandsraums folgendermaßen definieren:
\begin{align*}
    a &= \text{\texttt{nimm(schluessel)}} \\
    \text{pre}(a) &= \{ \text{\texttt{position(spieler, korridor)}}, \text{\texttt{position(schluessel, korridor)}} \} \\
    \text{del}(a) &= \{ \text{\texttt{position(schluessel, korridor)}} \} \\
    \text{add}(a) &= \{ \text{\texttt{inventar(spieler, schluessel)}} \}
\end{align*}

Prüft das System bei einer Eingabe, ob die Vorbedingungen ($\text{pre}(a)$) im aktuellen Zustand $S$ als wahr gelten, kann die Aktion ausgeführt werden. Durch die Mengendifferenz $\setminus \text{del}(a)$ verschwindet der Schlüssel aus dem Korridor. Durch die Vereinigungsmenge $\cup \text{add}(a)$ taucht er stattdessen im Inventar auf. Alle anderen Fakten des Zustands $S$, etwa die verriegelte Tür oder die Position einer brennenden Fackel, bleiben unangetastet.

Dass Prolog und STRIPS so ineinandergreifen, ist kein historischer Zufall. Beide Konzepte entstanden unabhängig voneinander in den frühen 1970er Jahren und teilen gemeinsame Wurzeln in der Prädikatenlogik erster Stufe (\textit{First-Order Logic}). Die theoretische Grundlage für die logische Modellierung von Handlungen und sich wandelnden Welten liefert das \textit{Situationskalkül} (Situation Calculus, \citealp{McCarthy1969}). Darin untersuchten McCarthy und Hayes, wie man mathematisch präzise über Zustände sprechen kann. 

STRIPS entstand als pragmatische Lösung für das dort beschriebene \textit{Frame-Problem}, also die programmtechnische Schwierigkeit, bei einer Handlung nicht nur angeben zu müssen, was sich \textit{ändert}, sondern auch endlos auflisten zu müssen, \textit{was} sich alles nicht ändert, also \textit{gleich} bleibt. Indem STRIPS dieses Problem über \textit{Preconditions} und \textit{Add/Delete-Listen} umgeht, schuf es eine Brücke zur \textit{ausführbaren Logik}. 

Prolog ist daher eine ideale Sprache, um solche logischen Planungsmodelle operativ umzusetzen, was detailliert in den großen Standardwerken der logischen Programmierung \citep{Bratko2001} oder der theoretischen Künstlichen Intelligenz \citep{Russell2020} beschrieben wird. Für unseren Kurs bedeutet das, dass wir die Mengenoperationen von STRIPS für die Formalisierung der Aktionen unserer IF verwenden werden.

\subsubsection{ÜBUNG STRIPS-Modell}

Erkläre das obige Minimalbeispiel natürlichsprachlich. Was sind $S$, $a$, $\text{pre}$, $\text{del}$ und $\text{add}$ und wie verändern sie den Zustand der Welt?

\chapter{Mit Logik programmieren} 

Wenn wir wissen, wie logische Aussagen formuliert und Wahrheitswerte für Aussagen über Objekte zugeordnet werden, können wir sie in Prolog‑Programme überführen. In diesem Kapitel lernen wir, wie wir aus diesen Prädikatenlogischen Aussagen ausführbare Programme werden, mit denen wir unsere Spielwelt erzeugen und ihr Leben einhauchen. 

Der Übergang zur Programmierung mit Hilfe von Logik erfordert einen Wechsel der Vorstellung, was Programmieren bedeutet. In den "klassischen" imperativen Sprachen (C/C++, Python, JavaScript) beschreibt man \textit{Abläufe} \textit{Schritt} für\textit{ Schritt}. Das geschieht, indem man das Problem zuerst \textit{selbst} löst und dann in Einzelschritte zur Abarbeitung der Lösung als \textit{Algorithmus} schreibt, der dann automatisiert abgearbeitet wird.

\subsection{Automatisiertes Schlussfolgern (Inferenz)}

In Prolog beschreibt man logische \textit{Zusammenhänge}, also Aussagen, was wahr ist, bzw. was gelten soll und unter welchen \textit{Bedingungen }etwas folgt. Also im Grunde so, wie wir es in der formalen Logik formuliert haben. Das Programm findet dann durch Inferenz, also durch logisches Ableiten (Rechnen mit Logik) die Lösungen, die aus diesen Zusammenhängen ableitbar sind, sofern sie existieren. 

\subsection{ Der Interpreter}

Prolog gehört damit zu den \textit{höheren} Programmiersprachen, da sie als \textit{interpretierte }Sprache eine \textit{höhere Abstraktionsebene} bietet. Mit Hilfe eines \textit{Interpreters} werden  die logischen Regeln in imperativen Code übersetzt, damit sie auf herkömmlicher Hardware ausgeführt werden kann. Technisch basiert die Ausführung der Inferenz meist auf der \textit{Warren Abstract Machine}. Die WAM ist \textit{Algorithmus} der darauf optimiert ist, den durch die definierten Fakten und Regeln aufgespannten \textit{Suchraum} mittels \textit{Unifikation} und \textit{Backtracking} möglichst effizient zu durchlaufen (traversieren). Dies ermöglicht es die Inferenz auf selbst stark verschachtelte Rätselketten auszuführen und Abhängigkeiten in Echtzeit auf Konsistenz zu prüfen.

Das Denkmodell für Prolog ist daher \textit{nicht der}\textit{ Algorithmus,} sondern die \textit{Wissensbasis}, eine Sammlung von \textit{Fakten} und \textit{Regeln}, an die man \textit{Anfragen} stellt und auf die Prolog durch automatisiertes Schlussfolgern (Inferenz) eine Antwort findet. Man programmiert nicht \textit{„Gehe in den nächsten Raum und prüfe, ob die Tür offen ist"}, sondern man formuliert: \textit{Es gibt Räume. Manche sind verbunden. Manche Verbindungen setzen Bedingungen voraus.} Und dann fragt man: \textit{Wohin kann ich von hier aus gehen?} und Prolog findet die Antwort. 

\subsection{Interactive Fiction und Prolog}

Für Interactive Fiction ist dieses Modell interessant, da es dem \textit{Entwurfsprozess }selbst entspricht. \textit{Man beschreibt eine Welt, ihre Objekte, Relationen und Regeln, und das System leitet daraus ab, was in dieser Welt möglich ist}. Die Welt wird nicht Zeile für Zeile\textit{ erzählt,} sondern als logisches System \textit{deklariert}. Dieser Unterschied zwischen \textit{deklarativ statt imperativ} ist der konzeptuelle Kern des gesamten Prolog-Kapitels.

Die Tupelnotation und STRIPS, die wir im vorherigen Kapitel  kennengelernt haben, verwenden wir als Vorstufe der Implementierung. Ein Modell wie $G=(V,E,W,S,\Sigma)$ legt fest, welche Räume, Verbindungen und Bedingungen eine Welt besitzt, in welchem Zustand sie sich befindet und welche Aktionen auf diesen Zustand wirken. In Prolog erscheinen die Elemente aus $V$ dann als Fakten von \texttt{raum/1}, die Elemente aus $E$ als Fakten von \texttt{verbindet/2} oder \texttt{verbindet/3}, und die Elemente aus $W$ als Bedingungen in Regeln. Der dynamische Zustand $S$ wird über Prolog-Fakten wie \texttt{position/2} oder \texttt{inventar/1} abgebildet. Die Aktionsmenge $\Sigma$ des STRIPS-Modells schließlich entspricht direkt den Prolog-Regeln, die mit \texttt{retract} und \texttt{assertz} den Zustand $S$ transformieren: Die Vorbedingungen ($pre(a)$) sind die Ziele im Regelrumpf (die Bedingungen, die erfüllt sein müssen), die Delete-Liste ($del(a)$) entspricht \texttt{retract}, die Add-Liste ($add(a)$) entspricht \texttt{assertz}. 

Die \textit{Tupelschreibweise} beantwortet also die \textit{strukturelle} Frage,\textit{ was} zum Modell gehört, \textit{STRIPS} beantwortet die\textit{ dynamische Frage}, \textit{wie} sich \textit{Aktionen} darauf \textit{auswirken}, und\textit{ Prolog }führt diese dann aus. Diese drei Ebenen werden uns in den nachfolgenden Aufgaben und Übungen begleiten.

\newpage
\section{Die Programmiersprache Prolog}

Prolog – \textit{PROgrammation en LOGique} – wurde 1972 von Alain Colmerauer an der Université d’Aix-Marseille entwickelt, mit dem ursprünglichen Ziel, natürliche Sprache maschinell verarbeitbar zu machen. Robert Kowalski (Universität Edinburgh) lieferte die theoretische Grundlage, die prozedurale Interpretation von \textit{Horn-Klauseln} \citep{Kowalski1974}, die wir später noch genauer kennenlernen werden.

\textit{Prolog ist damit keine Sprache, die nachträglich mit Logik ausgestattet wurde, sondern sie ist Logik, die zur Programmiersprache wurde.}

\subsection{Grundlagen}

Prolog lässt sich als formale Beschreibung von Wissen verstehen: Wir definieren, was gilt, und unter welchen Bedingungen weiteres gilt. Die mathematische Grundlage dafür ist die \textit{Mengenlehre}. Ein \textit{Prädikat }beschreibt eine \textit{Menge von Objekten} und eine \textit{Relation }beschreibt eine \textit{geordnete Beziehung zwischen Objekten}. So entsteht ein klar strukturiertes Modell, in dem Aussagen präzise formuliert und logisch überprüfbar werden.

\textbf{Prädikate} beschreiben \textit{Eigenschaften} von Objekten oder Beziehungen zwischen ihnen. Ein Prädikat wie \texttt{raum(X)} definiert die \textit{Menge} aller X, die Räume sind. Raum ist ein Subjektiv, kein Prädikat, daher wird \texttt{raum(X)} verwendet, um auszudrücken, dass X die \textit{Eigenschaft} hat, ein Raum zu sein. Ein Prädikat wie \texttt{position(X, Y)} beschreibt die Relation „X befindet sich in Y". 

\textbf{Fakten} beschreiben bedingungslose Wahrheiten innerhalb der Wissensbasis. Das folgende Faktum definiert, dass das Individuum \textit{labor} ein Raum ist:

\texttt{raum(eingang).}
\texttt{raum(korridor).}
\texttt{raum(labor).}

Ein weiteres Faktum kann festhalten, dass im Labor das Licht noch funktioniert:

\texttt{licht\_aktiviert(labor).}

\textbf{Regeln} definieren logische Implikationen und Abhängigkeiten.

\texttt{beleuchtet(X) :- raum(X), licht\_aktiviert(X).}

X ist beleuchtet, \textit{falls} X ein Raum ist \textit{und} dort das Licht aktiviert ist.

\textbf{Anfragen} bilden die Schnittstelle, um den Zustand der Welt zu prüfen. Die folgende Anfrage testet, ob die Bedingungen für \textit{beleuchtet} auf das Individuum \textit{labor} zutreffen.

\texttt{?- beleuchtet(labor).}

Wichtig: Jedes \textbf{Faktum} ist logisch gesehen eine spezielle \textbf{Regel ohne Bedingungen}.

\subsection{Reine Logik und Seiteneffekte}

Ein zentraler Aspekt bei der Programmierung in Prolog ist die Unterscheidung zwischen Programmteilen der „reinen“ deklarativen Logik und Programmteilen, die Seiteneffekt behaftet sind und prozedurale, nicht rein logische Befehle enthalten, beispielsweise für das Einlesen, oder Speichern von Dateien, oder das Ausgeben von Texten, auf dem Bildschirm, oder Drucker.

\textbf{Reine Logik:} Prädikate, die lediglich Wissen ableiten (z.\,B. Erreichbarkeit eines Raums), ohne den Zustand der Datenbank zu verändern.

\textbf{Seiteneffekte:} Um ein interaktives Spiel zu ermöglichen, nutzt Prolog Prädikate wie \texttt{write/1} oder \texttt{format/2} zur Textausgabe sowie \texttt{assertz/1} (Fakten hinzufügen) und \texttt{retract/1} (Fakten löschen). Letztere erlauben es, die Welt zur Laufzeit dynamisch zu verändern, etwa wenn ein Charakter einen Gegenstand aufnimmt und sich damit die „Wahrheit“ über dessen Aufenthaltsort permanent verschiebt. Solche Operationen liegen außerhalb der reinen Logik und können daher nicht\textit{ sicher }für ein logische Herleitung verwendet werden. 

Wir werden beide Arten kennenlernen, einerseits für die logische Formulierung unser Welt und Spielmechaniken und andererseits für das Textinterface, also die Interaktion mit den Spielenden, die trotz allem Minimalismus ja trotzdem auch Spaß machen soll.

\subsection{Prolog starten und ein Programm laden}

Für den Kurs verwenden wir \textbf{Scryer Prolog}, einen modernen, ISO-nahen Prolog-Interpreter. Nach der Installation startet man den Interpreter im Terminal mit:
\begin{lstlisting}[language=bash, caption={}]
scryer-prolog
\end{lstlisting}
Eine Programmdatei kann auch direkt beim Start als Parameter übergeben werden:
\begin{lstlisting}[language=bash, caption={}]
scryer-prolog welt.pl
\end{lstlisting}
Der Interpreter zeigt den Prompt \texttt{?-} und ist bereit für Anfragen. Ein Prolog-Programm schreibt man in einer Textdatei mit der Endung \texttt{.pl}, zum Beispiel \texttt{welt.pl}. Um es in den Interpreter zu laden, verwendet man \texttt{consult/1}:
\begin{lstlisting}[caption={}]
?- consult('welt.pl').
\end{lstlisting}
Kurzform mit eckigen Klammern:
\begin{lstlisting}[caption={}]
?- [welt].
\end{lstlisting}
Nach dem Speichern einer geänderten Datei lädt man sie erneut mit \texttt{consult/1} oder der Kurzform:
\begin{lstlisting}[caption={}]
?- consult('welt.pl').
?- [welt].
\end{lstlisting}
\subsection{Navigation im Interpreter}

\texttt{halt.} beendet den Interpreter\\
\texttt{;} nach einer Antwort fordert die nächste Lösung an\\
\texttt{.} (Enter) akzeptiert die aktuelle Lösung und bricht ab\\
\texttt{Ctrl+C} unterbricht eine laufende Anfrage\\

\subsection{Wichtige Sprachkonzepte und Begriffe}

In diesem Kurs verwenden wir bevorzugt ISO-konforme Prädikate, damit die Programme portabel bleiben. Die folgenden Konzepte definieren die grundlegende Terminologie, die wir für die weitere Arbeit benötigen.

\textbf{Syntax und Semantik}

Die \textbf{Syntax} beschreibt, welche\textit{ Zeichenfolgen} in Prolog formal gültig sind. Die \textbf{Semantik} beschreibt, welche \textit{Bedeutung} diese Ausdrücke haben. Diese Trennung ist wichtig, weil ein Programm \textit{syntaktisch korrekt} und zugleich fachlich unpassend modelliert sein kann.

\textbf{Term}

Der Grundbaustein in Prolog ist der \textbf{Term}. ISO-Prolog unterscheidet im Kern die folgenden Termtypen:
\textbf{
Variablen} beginnen mit Großbuchstaben oder Unterstrich, z.\,B. \texttt{X}, \texttt{Raum}, \texttt{\_Tmp}\\
\textbf{Anonyme Variable}: Einzelner Unterstrich \verb|_| wird  verwendet, wenn der Wert der Variable irrelevant ist.
\textbf{Atome}: symbolische Konstanten, meist kleingeschrieben, z.\,B. \texttt{eingang}, \texttt{verschlossen}\\
\textbf{Zahlen}: numerische Konstanten, z.\,B. \texttt{3}, \texttt{2.5}\\
\textbf{zusammengesetzte Terme}: \texttt{f(t1,...,tn)}, z.\,B. \texttt{position(spieler, eingang)}, oder \texttt{verbindet(eingang, korridor)}\\

Ein \textbf{Individuum} (Entität) der modellierten Welt wird in Prolog in der Regel durch ein \textit{Atom} repräsentiert, etwa \texttt{eingang} oder \texttt{laterne}. Das Atom ist dabei \textit{kein Textobjekt} im Sinne eines freien \textit{Strings}, sondern ein \textit{formales Symbol} des Modells.

\textbf{Funktor, Argument, Arität}

Im Term \texttt{beleuchtet(labor)} ist  der \textbf{Funktor}, \texttt{labor} ist das \textbf{Argument}, und hat die \textbf{Aritaet} \texttt{1}. Die Schreibweise \texttt{beleuchtet/1} ist der \textbf{Praedikatsindikator}. Er identifiziert ein Praedikat eindeutig ueber Name und Aritaet.

\textbf{Prädikat}

Ein \textbf{Prädikat} beschreibt eine Eigenschaft oder Relation, z.\,B. \texttt{raum/1} oder \texttt{verbindung/2}. Alle Klauseln mit \textit{demselben} Prädikatsindikator gehören logisch zum selben Prädikat.

\textbf{Klausel}

Ein Prolog-Programm ist eine Menge von \textbf{Klauseln}. In der grundlegenden Arbeitsform nutzen wir definite Horn-Klauseln:

\begin{itemize}
    \item \textbf{Fakt}: nur Kopf, z.\,B. \texttt{raum(eingang).}
    \item \textbf{Regel}: Kopf mit Bedingungen, z.\,B. \texttt{beleuchtet(X) :- raum(X), licht\_aktiviert(X).}
\end{itemize}

Die Regel liest sich so: \textit{Der Kopf gilt genau dann, wenn alle Ziele im Rumpf gelten.}

\textbf{Anfrage (Goal)}

Eine \textbf{Anfrage} ist ein Ziel an den Interpreter, z.\,B.:
\begin{lstlisting}[caption={}]
?- beleuchtet(labor).
\end{lstlisting}
Das System versucht dieses Ziel durch Unifikation und Rückwärtsschluss aus den vorhandenen Klauseln herzuleiten.

\textbf{Beispiel}

Unser Beispiel zeigt die gesamte Kette von der Symbolwahl (Atome), über die Struktur (Terme), bis zur Modellierung mit Praedikaten, Klauseln und Anfrage.\\

\begin{lstlisting}[caption={}]
% Kommentare beginnen mit % und werden vom Interpreter ignroriert.

raum(labor).             % Fakt: labor ist ein Raum
licht_aktiviert(labor).  % Fakt: Licht ist noch an
beleuchtet(X) :-         % Regel: X ist beleuchtet, wenn ...
    raum(X), 
    licht_aktiviert(X).

?- beleuchtet(labor).          % Anfrage
true.                          % Prolog hat eine Ableitung gefunden.
\end{lstlisting}

\subsection{Abfragen und Unifikation}
Abfragen sind die operative Schnittstelle zur Wissensbasis. In der mathematischen Sicht sind sie Ziele, deren Erfüllbarkeit aus Fakten und Regeln hergeleitet wird. In der Prolog-Praxis beginnt eine Abfrage mit \texttt{?-} und endet mit einem Punkt.

\textbf{Grundform von Abfragen}
\begin{lstlisting}
?- beleuchtet(labor).
true.
\end{lstlisting}
Diese Abfrage ist eine Wahr-Falsch(Ja-Nein)-Frage. Das System antwortet \texttt{true}, wenn das Ziel aus der Wissensbasis folgt, andernfalls \texttt{false}.

\textbf{Abfragen mit Variablen}
\begin{lstlisting}
?- beleuchtet(X).
   X = labor.
\end{lstlisting}
Hier wird nicht nur geprueft, ob ein Ziel gilt, sondern es werden auch Werte für \texttt{X} gesucht, die das Ziel erfüllen, also wahr werden lassen. Prolog liefert damit nicht nur Wahrheitswerte, sondern konkrete \textit{Bindungen}, also die \textit{Einsetzungen für X}, die die Anfragen wahr werden lassen.

\textbf{Konjunktion und Disjunktion
}
Mehrere Ziele können kombiniert werden. Die Konjunktion\texttt{,} bedeutet logisches UND. Die Disjunktion \texttt{;} bedeutet logisches ODER:
\begin{lstlisting}
?- raum(X), beleuchtet(X).
?- beleuchtet(labor); beleuchtet(korridor).
\end{lstlisting}
\subsection{Unifikation}

Unifikation ist der Prozess, zwei Terme durch geeignete Variablenbindungen gleich zu machen. Formal gesucht wird eine Substitution, die beide Terme identisch macht. Falls mehrere Substitutionen möglich sind, ist die \textit{allgemeinste Unifikator-Substitution} (most general unifier, MGU) zentral, da alle spezielleren Loesungen daraus ableitbar sind.
\begin{lstlisting}[caption={}]
% unifiziert: X wird an labor gebunden
beleuchtet(X) = beleuchtet(labor).   

 % scheitert: zwei verschiedene Atome
beleuchtet(labor) = beleuchtet(korridor).
\end{lstlisting}
\textbf{Wichtige Prädikate zur Gleichheit}

\texttt{=/2}: versucht zu unifizieren\\
\texttt{==/2}: testet strikte Termidentität ohne neue Bindungen\\
\texttt{\textbackslash =/2}: testet Nicht-Unifizierbarkeit\\

Fuer theoretisch strenge Erstordnungs-Unifikation ist der Occurs-Check relevant. Das ISO-Prädikat \texttt{unify\_with\_occurs\_check/2} erzwingt diesen Check explizit.
\begin{lstlisting}[caption={}]
beleuchtet(labor).

?- X = eingang.
true.
% bindet X die freie Variable an eingang

?- eingang = labor.
false.
% scheitert, da eingang und labor unterschiedliche Atome sind

?- beleuchtet(X), X == labor.
true.
% unifiziert X mit labor, dann testet die Identitaet, was true ergibt
\end{lstlisting}
Diese Beispiele zeigen drei unterschiedliche Ebenen: reine Variablenbindung, fehlschlagende Gleichsetzung verschiedener Atome und Kombination aus Ableitung plus Identitätsprüfung.

\subsection{Backtracking und der Cut}
Backtracking ist der grundlegende Suchmechanismus von Prolog. Wenn ein Ziel scheitert, springt das System zum letzten offenen Entscheidungspunkt zurück und versucht die nächste Alternative. Operational entspricht das der schrittweisen Traversierung eines Suchbaums mit Rücksprung auf Choice Points.

\textbf{Backtracking als Suchstrategie}

Die deklarative Lesart eines Programms bleibt unverändert, die operationale Ausführung erfolgt in den meisten Implementierungen durch Tiefensuche und links-nach-rechts Selektion. Dadurch ist die Reihenfolge von Regeln und Teilzielen für Laufzeitverhalten und Terminierung relevant.

Backtracking tritt auf, wenn ein Prädikat mehrere Klauseln hat. Prolog versucht sie der Reihe nach. Scheitert eine, springt es zurück (\textit{backtrack}) und probiert die nächste. Das wird erst sichtbar, wenn eine Variable im Kopf der Anfrage verschiedene Werte annehmen kann:

\begin{lstlisting}[caption={}]
% Welche Methoden gibt es, das Archiv zu betreten?
zugang(archiv, schluessel).
zugang(archiv, dietrich).
zugang(archiv, ausweiskarte).

% Inventar der Spielfigur:
inventar(spieler, schluessel).

% Alle Zugaenge des Archivs:
?- zugang(archiv, Methode).
Methode = schluessel ;
% ; im Terminal fordert die naechste Loesung an (Backtracking sichtbar)
Methode = dietrich ;
Methode = ausweiskarte.

% Nur die, die der Spieler gerade hat:
?- zugang(archiv, Methode), inventar(spieler, Methode).
Methode = schluessel.
% Prolog probiert alle zugang-Klauseln durch;
% nur schluessel erfuellt auch inventar(spieler, schluessel)
\end{lstlisting}

Das erste \texttt{?-} liefert alle drei Alternativen. Das zweite kombiniert Backtracking mit einer Bedingung: Prolog probiert jede Methode und prüft, ob der Spieler sie hat. Die ersten drei Klauseln sind reine Fakten ohne Variablen, aber die \emph{Anfrage} hat eine Variable (\texttt{Methode}) und zeigt damit die Alternativen. Alle Lösungen auf einmal:

\begin{lstlisting}[caption={}]
?- findall(M, zugang(archiv, M), Methoden).
Methoden = [schluessel, dietrich, ausweiskarte].
\end{lstlisting}
Das war Backtracking mit sichtbarer Variable. Oft braucht man nur eine Ja/Nein-Antwort, ohne alle Alternativen aufzuzählen. Dann hat das Prädikat kein Argument, das verschiedene Werte annimmt. Das Backtracking findet intern trotzdem statt, bleibt aber nach außen unsichtbar:
\begin{lstlisting}[caption={}]
% Spielzustand:
inventar(spieler, schluessel).
zustand(archivtuer, verschlossen).
% zustand(archivtuer, offen) existiert NICHT als Fakt (CWA)

% Klausel 1: Spieler hat den Schluessel
archiv_zugaenglich :- inventar(spieler, schluessel).

% Klausel 2: Tuer wurde vorher bereits geoeffnet
archiv_zugaenglich :- zustand(archivtuer, offen).

?- archiv_zugaenglich.
true.
% Klausel 1 greift. Prolog liefert true.
% Bei Backtracking wuerde Klausel 2 aber noch geprueft.
\end{lstlisting}
Man kann hier den Cut einsetzen. Da \texttt{archiv\_zugaenglich/0} nur \texttt{true} oder \texttt{false} liefert (nie mehrere Werte), ist es redundant, nach dem ersten Erfolg weitere Klauseln zu prüfen. Ein Cut nach dem ersten Ziel verhindert das:
\begin{lstlisting}[caption={}]
archiv_zugaenglich :- inventar(spieler, schluessel), !.
archiv_zugaenglich :- zustand(archivtuer, offen).

?- archiv_zugaenglich.
true.
% Mit Cut: Klausel 2 wird nie geprueft, wenn Klausel 1 greift.
% Ohne Cut: Prolog koennte bei Backtracking Klausel 2 auch noch ausprobieren.
\end{lstlisting}
Das ist ein \textbf{Green Cut}: Er entfernt einen Suchpfad (Klausel 2), der ohnehin nur \texttt{true} liefern könnte, das beobachtbare Ergebnis ändert sich nicht. Im Unterschied dazu würde ein Red Cut eine Lösung \emph{unterdrücken}, die ohne Cut sichtbar wäre. Weil \texttt{archiv\_zugaenglich/0} keine Variable hat, gibt es keine unterdrückbare Lösung. Der Cut hier ist reine Effizienz.

Der Unterschied zu \texttt{zugang/2} oben: \texttt{archiv\_zugaenglich/0} hat keine Variable im Kopf, die Antwort ist immer nur \texttt{true} oder \texttt{false}. Das Backtracking passiert intern zwischen den Klauseln, ist aber von außen nicht sichtbar. Die Wahl zwischen beiden Mustern hängt davon ab, ob man die \emph{gefundene Alternative} weiterverarbeiten möchte (dann Variable) oder nur wissen will, ob \emph{irgendeine} Alternative existiert (dann 0-stellig).

\textbf{Der Cut als Kontrolloperator}

Der Cut \texttt{!/0} ist ein prozeduraler Operator. Wird ein Cut erreicht, verwirft Prolog die bis dahin offenen Alternativen im aktuellen Prädikatsstufe. Der Cut beeinflusst also nicht nur Effizienz, sondern auch die beobachtbare Menge der Antworten.

\textbf{Typisches Muster}
\begin{lstlisting}[caption={}]
% labor ist nur betretbar, wenn die Laterne mitgenommen wurde
betretbar(labor) :- inventar(spieler, laterne), !.
betretbar(labor) :- !, write('Es ist zu dunkel.'), nl, fail.
betretbar(X)     :- raum(X).   % alle anderen Raeume immer betretbar
\end{lstlisting}
Wenn \texttt{inventar(spieler, laterne)} gilt, greift Klausel~1 und der Cut verhindert, dass Klausel~2 oder~3 geprüft werden. Hat der Spieler die Laterne nicht, scheitert Klausel~1 vor dem Cut, Prolog versucht Klausel~2: der Cut dort sperrt danach Klausel~3, die Fehlermeldung wird ausgegeben, und \texttt{fail} sorgt für \texttt{false}. Ohne den Cut in Klausel~2 würde Prolog nach dem \texttt{fail} zu Klausel~3 weitergehen und \texttt{betretbar(labor)} über \texttt{raum(labor)} als \texttt{true} zurückliefern, also trotz Dunkelheit Zugang gewähren. Das ist ein Red Cut: Beide Cuts zusammen verändern die sichtbare Ergebnismenge, nicht nur die Effizienz.

\textbf{Green Cut und Red Cut}

\begin{itemize}
    \item \textbf{Green Cut}: entfernt nur redundante Suchpfade, ohne die deklarative Bedeutung zu ändern.
    \item \textbf{Red Cut}: verändert die Menge der ableitbaren Lösungen und damit die Semantik des Programms.
\end{itemize}

Es sollte bevorzugt mit Green Cuts gearbeitet werden. Red Cuts sind nur dann vertretbar, wenn die beabsichtigte Verhaltensänderung explizit dokumentiert und fachlich begründet ist. Der Cut ist ein Kontrollwerkzeug für Suchraumbegrenzung. Daher gilt zuerst logische Korrektheit herstellen, danach gezielt optimieren.

\subsection{Determinismus-Angaben und Meta-Argumente}

Die Dokumentation von Prolog-Prädikaten verwendet standardisierte Präfixe, die das Interface (die Datenein- und ausgänge) beschreiben, also wie die Daten an die Parameter (Argument-Binding) gebunden werden.

\begin{center}
\begin{tabular}{ll}
\toprule
\textbf{Annotation} & \textbf{Bedeutung} \\
\midrule
\texttt{det} & Genau eine Lösung, kein Backtracking \\
\texttt{semidet} & Eine Lösung oder keine \\
\texttt{nondet} & Mehrere Lösungen möglich (auch null) \\
\texttt{multi} & Mehrere Lösungen möglich, mindestens eine \\
\texttt{failure} & Schlägt immer fehl (z.\,B.\ \texttt{repeat/fail}-Loop) \\
\texttt{+Argument} & Eingabe (das Argument ist gebunden) \\
\texttt{-Argument} & Ausgabe (das Argument wird gebunden) \\
\texttt{?Argument} & Ein- oder Ausgabe \\
\texttt{@Argument} & Eingabe, wird nicht gebunden (nur inspiziert) \\
\texttt{:Argument} & Meta-Argument: das Argument wird als Goal aufgerufen \\
\bottomrule
\end{tabular}
\end{center}

Die Argument-Präfixe erscheinen vor dem Variablennamen, z.\,B. \texttt{+NeuerRaum} oder \texttt{-Text}. \texttt{+} bedeutet, der Aufrufer muss den Wert bereitstellen; \texttt{-} bedeutet, das Prädikat bindet die Variable; \texttt{?} heißt beides möglich. Das \texttt{:}-Präfix markiert Meta-Argumente – Terme, die nicht als Daten, sondern als aufzurufende Ziele (Goals) verwendet werden. In der Signatur \texttt{forall(:Bedingung, :Aktion) is det} bedeutet \texttt{:Bedingung}, dass an dieser Stelle ein Goal erwartet wird (z.\,B. \texttt{member(X, [1,2,3])}), nicht etwa ein Atom oder eine Zahl.

Diese Annotationen finden sich im gesamten Kursmaterial: Sie machen die Schnittstelle eines Prädikats auf einen Blick lesbar, ohne den Code verstehen zu müssen.

\section{Prolog und Interactive Fiction}

Mit diesem Abschnitt beginnt der eigentliche Einstieg in das Programmieren von Prolog. Die folgenden Beispiele und Abschnitte orientieren sich an den zentralen Elementen einer IF-Welt und dessen Spielmechaniken. Die Welt selbst ist Wissensbasis, der Datenbasis mit Fakten und Regeln, der Interaktion durch Interface und Steuerung, sowie der Gestaltung von Rätseln und Puzzles. In jedem Abschnitt werden wir die logischen Strukturen der Spieleemenete und Spielmechanik kennenlernen, und sie in ihr logische Struktur zerlegen, die für ihre Umsetzung in Prolog notwendig ist und diese dann in Prolog programmieren.

\subsection{ Das Institut}

Alle Prolog-Beispiele in diesem Kapitel bauen auf derselben Welt auf, einem verlassenen Forschungsgebäude. Ziel ist es, die \textit{akte} aus einem unbekannten, v\textit{ersiegelten Archiv} zu bergen. Die Welt wird Schritt für Schritt modelliert, beginnend mit einfachen Fakten und endend mit einem spielbaren Prototyp.

\begin{center}
\begin{tabular}{lll}
\toprule
\textbf{Element} & \textbf{Beschreibung} & \textbf{Tupel-Komponente} \\
\midrule
Räume & \texttt{eingang}, \texttt{korridor}, \texttt{labor}, \texttt{archiv} & $V$ \\
Verbindungen & eingang$\leftrightarrow$korridor, korridor$\leftrightarrow$labor & $E$ \\
Tür & \texttt{archivtuer} zwischen korridor und archiv (gesperrt) & $W$ \\
Objekte & \texttt{laterne} (eingang), \texttt{schluessel} (korridor), \texttt{akte} (archiv) & $S$ \\
Aktionen & Navigation, Aufnehmen, Türöffnen & $\Sigma$ \\
\bottomrule
\end{tabular}
\end{center}

Die folgende Faktenbasis ist der Ausgangspunkt. Statische Fakten beschreiben die Struktur der Welt; dynamische Fakten den veränderlichen Spielzustand:\\
\begin{lstlisting}[language=Prolog, caption={}]
% --- Statisches Weltmodell (unveraenderlich) ---

raum(eingang).
raum(korridor).
raum(labor).
raum(archiv).

verbindet(eingang,  korridor).   % bidirektional: je Richtung ein Fakt
verbindet(korridor, eingang).
verbindet(korridor, labor).
verbindet(labor,    korridor).

tuer(korridor, archiv, archivtuer).  % Tuer zwischen zwei Raeumen
tuer(archiv,   korridor, archivtuer).

gegenstand(laterne).
gegenstand(schluessel).
gegenstand(akte).

% --- Dynamischer Spielzustand (veraenderlich) ---

:- dynamic position/2.   % position(Entitaet, Ort)
:- dynamic inventar/2.   % inventar(Spieler, Gegenstand)
:- dynamic zustand/2.    % zustand(Entitaet, Zustand)

position(spieler, eingang).
position(laterne, eingang).
position(schluessel, korridor).
position(akte, archiv).

zustand(archivtuer, verschlossen).
\end{lstlisting}

\noindent
\textbf{Prädikatübersicht:} \texttt{raum/1}, \texttt{verbindet/2}, \texttt{tuer/3} und \texttt{gegenstand/1} sind statisch und werden nie verändert. \texttt{position/2}, \texttt{inventar/2} und \texttt{zustand/2} sind dynamisch und werden durch Spielaktionen mit \texttt{retract}/\texttt{assertz} verändert.

\subsection{Die Welt als Wissensbasis}
Eine der ursprünglichen Motivationen war das Anlegen von Wissensdatenbanken und Expertinn:ensystemen, die das Wissen über eine Domäne in Form von Fakten und Regeln repräsentieren. Prolog-Programme bestehen daher immer aus einer Wissensbasis, die Fakten und Regeln enthält, die die Welt beschreiben. Damit ist die Datenbasis gleich auch das eigentliche Prolog Programm. Wenn wir Daten umschreiben, beispielsweise die Position der Spielfigur, verändern wir nicht nur den Zustand der Welt, sondern auch die Datenbasis und damit unser Prolog-Programm. Das geschieht allerdings nur im Arbeitsspeicher, nicht in der Quelldatei auf der Festplatte.

\textbf{Das Universum}
In der Logikprogrammierung nennt man die Menge aller Grundterme (also aller Terme ohne Variablen) das Herbrand‑Universum, benannt nach dem Logiker Jacques Herbrand. Für das Modellieren einer Welt bedeutet das, dass nur die Objekte „existieren“, die wir als Terme in unseren Fakten und Regeln explizit definieren. Objekte, die sich nicht als solche Terme schreiben lassen, gehören nicht zum Herbrand-Universum.

\textbf{Closed World Assumption (CWA)}
Closed World Assumption bedeutet in Prolog, dass Aussagen, die sich aus der Wissensbasis nicht herleiten lassen, als nicht wahr behandelt. Das Negationsprädikat \texttt{\textbackslash+} prüft entsprechend die Nichtableitbarkeit einer Aussage. Für Interactive Fiction ist dieses Prinzip nützlich, weil Zustände eindeutig und kontrollierbar modelliert werden können. Zugleich muss die Wissensbasis hinreichend vollständig formuliert sein, damit fehlende Fakten nicht fälschlich als inhaltliche Verneinung interpretiert werden. \\
\begin{lstlisting}[caption={}]
% Die archivtuer existiert, ist aber nicht als offen definiert.
zustand(archivtuer, verschlossen).

?- zustand(archivtuer, offen).
false.
\end{lstlisting}
Dies nutzen wir für das Institut: \texttt{zustand(archivtuer, verschlossen)} ist explizit gesetzt. Das Gegenteil, \texttt{zustand(archivtuer, offen)}, existiert nicht in der Wissensbasis und gilt damit als falsch. Das erspart einen separaten \texttt{nicht\_offen}-Fakt. Der Zustand kann später per \texttt{retract}/\texttt{assertz} geändert werden.

\begin{lstlisting}[caption={}]
% Nur passierbar, wenn die Tuer offen ist (CWA: alle anderen sind zu)
kann_gehen(X, Y) :-
    tuer(X, Y, T),
    zustand(T, offen).  % kanonisch: zustand/2, nicht offen/1
\end{lstlisting}
\textbf{Open World Assumption (OWA)}

Im Gegensatz zur Closed World Assumption geht die Open World Assumption davon aus, dass die Wissensbasis unvollständig sein kann. Eine Aussage, die sich nicht herleiten lässt, gilt dann nicht als falsch, sondern als noch unbestimmt. Das ist konzeptionell verschieden von einer Welt, in der alle Fakten von Anfang an feststehen. 

Ein sinnvoller Anwendungsfall in Interactive Fictiion: Ein unbekanntes Artefakt liegt im Raum. Seine Eigenschaften sind noch nicht Teil der Wissensbasis, weil der Spielstand sie noch nicht festgelegt hat. Weder \texttt{giftig(kelch)} noch \texttt{\textbackslash+ giftig(kelch)} ist sinnvoll ableitbar, bevor das Objekt untersucht wurde. 

\begin{lstlisting}[caption={}]
% Vor dem Untersuchen: kein artefakt_typ in der Wissensbasis
giftig(X) :- artefakt_typ(X, gift).

?- giftig(kelch).   % false unter CWA, unter OWA: unbestimmt

% Nach dem Untersuchen (assert durch Spielaktion):
% artefakt_typ(kelch, gift).

?- giftig(kelch).   % true
\end{lstlisting}

Die Entscheidung zwischen CWA und OWA ist eine Designentscheidung. Ist die Welt vollständig beschrieben und alles Ungenannte falsch, oder gibt es Spielraum für Unbestimmtheit? Aus Sicht der formalen Logik ist hier Vorsicht geboten:  Prolog mit \texttt{\textbackslash+} implementiert formal 
\textit{Negation as Failure} (NAF), keine echte logische Negation.

Klassische Negation und \textit{Negation as Failure} verhalten sich grundverschieden. \texttt{\textbackslash+ giftig(kelch)} bedeutet nicht, dass der Kelch nachweislich ungiftig ist, sondern nur, dass Prolog \texttt{giftig(kelch)} nicht beweisen kann. Die Aussage scheitert mangels Beweis, nicht weil das Gegenteil bekannt wäre. Echte logische Negation \texttt{¬giftig(kelch)} wäre dagegen eine positive Aussage in der Wissensbasis: Der Kelch ist explizit als ungiftig deklariert. Im IF-Kontext wird der Unterschied praktisch: Sobald \texttt{assertz(giftig(kelch))} einen neuen Fakt hinzufügt, kollabiert NAF still – der Kelch ist plötzlich giftig, ohne dass ein Widerspruch entsteht. Echte Negation hätte \texttt{¬giftig(kelch)} und \texttt{giftig(kelch)} als Konflikt sichtbar gemacht.

Sobald \texttt{assertz} und \texttt{retract} hinzukommen, ist das Programm keine unveränderliche Theorie mehr. Logik-Puristen ordnen das der \textit{nicht-monotonen Logik} zu (Default Logic, \citealt{Reiter1980}): Schlüsse gelten nur so lange, bis neue Fakten sie revidieren. Man verlässt damit die reine deklarative Logikprogrammierung. Für Interactive Fiction und viele reale Prolog Anwendungen ist das pragmatisch sinnvoll, aber es lohnt sich zu wissen, wo man die \textit{reine Logik} hinter sich lässt.

In diesem Kapitel arbeiten wir ausschließlich mit der \textit{Closed World Assumption}. Das hält die Modellierung klar und die Semantik nachvollziehbar. OWA-basierte Mechaniken sowie weitere fortgeschrittene Techniken behandeln wir im Kapitel \textit{Fortgeschrittene Interactive Fiction Mechaniken}.

\subsection{Datenbasis}
Die Datenbasis haben wir schon kennengelernt als den Ort, an dem unser Programm entsteht, indem wir Fakten und Regeln definieren.

Für das Design interaktiver Systeme ist die Trennung zwischen statischem Weltwissen und veränderlichem Spielzustand wichtig. Dynamisch modelliert werden alle zustandsrelevanten Elemente der Spielwelt, etwa Positionen von Figuren und Objekten, Inventar, Zugänge, Schalterzustände sowie Dialog- und Questfortschritte. 

\subsection{Daten hinzufügen, ändern und entfernen}

Technisch erfolgt das mit \texttt{dynamic/1}. Zustandsübergänge werden dann zur Laufzeit mit \texttt{assertz/1} und \texttt{retract/1} in der Wissensbasis vollzogen.

\textbf{Spielmechanik:} Der Spieler gibt einen Befehl wie \texttt{gehe\_nach(korridor)} ein. Die Spielfigur verlässt ihren aktuellen Raum und erscheint im Zielraum, sofern eine Verbindung besteht.

\textbf{Formalisierung (STRIPS):}
\begin{itemize}
    \item $pre(a)$: \texttt{position(spieler, AktuellerRaum)}, \texttt{verbindet(AktuellerRaum, NeuerRaum)}
    \item $del(a)$: \texttt{position(spieler, AktuellerRaum)}
    \item $add(a)$: \texttt{position(spieler, NeuerRaum)}
\end{itemize}

\textbf{Prolog:}
\begin{lstlisting}[language=Prolog, caption={}]
% position/2 wurde bereits als :- dynamic deklariert

gehe_nach(NeuerRaum) :-
    position(spieler, AktuellerRaum),               % pre(a)
    verbindet(AktuellerRaum, NeuerRaum),            % pre(a)
    retract(position(spieler, AktuellerRaum)),      % del(a)
    assertz(position(spieler, NeuerRaum)).          % add(a)
\end{lstlisting}
\texttt{retract/1} entfernt die alte Position (STRIPS-\textit{Delete List}), \texttt{assertz/1} schreibt die neue Position in die Wissensbasis (STRIPS-\textit{Add List}). Die beiden \texttt{verbindet}-Prüfungen sind die Vorbedingungen ($pre(a)$): Nur wenn eine Verbindung in der Faktenbasis existiert, gelingt die Aktion.
\begin{lstlisting}[language=Prolog, caption={}]
?- position(spieler, eingang).
true.

?- gehe_nach(korridor).
true.

?- position(spieler, eingang).
false.

?- position(spieler, korridor).
true.

% Kein Weg von eingang direkt nach archiv:
?- gehe_nach(archiv).
false.
\end{lstlisting}
Genau hier, in diesen \texttt{retract}- und \texttt{assertz}-Befehlen, sehen wir die direkte technische Umsetzung des \textbf{STRIPS-Modells}:

Die Vorbedingungen ($pre(a)$) sind die Ziele im Regelrumpf: Der Spieler muss irgendwo sein, und zwischen diesem Ort und dem Ziel muss eine Verbindung bestehen. Das Entfernen der alten Position aus dem Zustand $S$ (\textit{Delete List}, $del(a)$) übernimmt \texttt{retract}, das Hinzufügen der neuen Position (\textit{Add List}, $add(a)$) leistet \texttt{assertz}. 

Die Mengenoperation $S' = (S \setminus del(a)) \cup add(a)$ ist damit exakt das, was Prolog ausführt.

\subsection{Fakten und Relationen}

\textbf{Spielmechanik:} Objekte und Räume existieren in der Spielwelt. Ein Gegenstand liegt an einem bestimmten Ort und der Spieler kann ihn nur wahrnehmen oder aufnehmen, wenn er sich im selben Raum befindet.

\textbf{Formalisierung:} Jedes Objekt und jede Beziehung wird als Fakt codiert. Die Relation $\text{position}(X, Y)$ sagt: Objekt $X$ befindet sich in Raum $Y$. Ob ein Gegenstand sichtbar ist, ergibt sich dann durch Anfrage: $\text{position}(\text{spieler}, R) \land \text{position}(G, R)$.

\textbf{Prolog:} Fakten und Relationen bilden den statischen Kern der Wissensbasis.

\begin{center}
\begin{tabular}{p{0.23\textwidth}p{0.34\textwidth}p{0.34\textwidth}}
\toprule
\textbf{Aspekt} & \textbf{Einfaches Faktum} & \textbf{Verbundenes Faktum} \\
\midrule
Stelligkeit & 1-stellig (ein Argument) & 2-stellig (zwei Argumente) \\
Bedeutung & Eigenschaft eines Objekts & Beziehung zwischen Objekten \\
Typische Nutzung & Attributzuweisung & Verbindungen oder Zuordnungen modellieren \\
Beispiel & \texttt{rot(auto).} & \texttt{verbunden(raum1, raum2).} \\
\bottomrule
\end{tabular}
\end{center}

\textbf{Beispiele für unterschiedliche Stelligkeit von Fakten}

\textbf{Einfaches Faktum (1-stellig):} \texttt{raum(eingang).}
Dieses Format eignet sich, wenn genau eine Eigenschaft eines einzelnen Objekts festgehalten wird.

\textbf{Zweistelliges Faktum:} \texttt{position(schluessel, korridor).}
Dieses Format wird genutzt, wenn eine Beziehung zwischen zwei Objekten modelliert werden soll, hier Objekt und Ort.

\textbf{Mehrstelliges Faktum (3-stellig):} \texttt{tuer(korridor, archiv, archivtuer).}
Dieses Format ist sinnvoll, wenn zusätzlich zur Grundbeziehung weitere Information direkt am Fakt mitgeführt werden soll, hier die konkrete Tür als verbindendes Element.\\

\textbf{Wann verwenden wir welche Form?}

Diese kurzen Merkregeln werden im Laufe des Kurses an konkreten Beispielen klarer werden.

\begin{itemize}
    \item \textbf{Fakt, 1-stellig:} reine Eigenschaft eines Objekts.
    \item \textbf{Fakt, 2-stellig:} einfache Beziehung zwischen zwei Objekten.
    \item \textbf{Fakt, 3-stellig oder mehr:} nur wenn alle Angaben immer gemeinsam gebraucht werden.
    \item \textbf{Regel statt Fakt:} sobald etwas abgeleitet wird oder Bedingungen hat.
    \item \textbf{Liste statt vieler Einzel-Fakten:} wenn Reihenfolge, Sammlung oder Stapeloperationen wichtig sind.
\end{itemize}

\subsubsection{ÜBUNG: Architektur und Objekte}

\begin{itemize}
    \item Erstellt die ersten Räume der Welt mit drei Räumen und definiert die Verbindungen zwischen ihnen.
    \item Fügt Gegenstände und Objekte hinzu, die sich in einem der Räume befinden.
    \item Definiert Gegenstände, die sich in anderen Gegenständen befinden können, z.B. einen Schlüssel in einer Truhe.
    \item Definiert die Möglichkeit, bestimmte Gegenstände von einem Raum in einen anderen zu bewegen.
\end{itemize}

\subsubsection{Spiel Charaktere und Inventar}

Spielfiguren und NPCs lassen sich in Prolog direkt als Fakten und Regeln modellieren. Dabei stellt sich schnell die Frage, wie weit die Mechaniken gehen sollen. Klassische IF konzentriert sich auf Storytelling und Rätsel, denn Kampfmechaniken können den Lesefluss stören. Wenn eine Geschichte von Charakterentwicklung handelt, der Spieler aber nur Zahlenwerte durch Skillpunkte erhöht, entsteht eine ludo-narrative Dissonanz, die die Immersion untergräbt. Wer IF liest, möchte oft nicht gleichzeitig komplexe Kampfstatistiken verwalten. 

Sinnvoll werden solche Mechaniken, wenn sie narrativ bedeutungsvoll sind. Konflikte, die nicht als Würfelmechanik, sondern durch echte Entscheidungen dargestellt werden (z.B. \enquote{einen Trick anwenden} statt \enquote{angreifen}), erhöhen die Immersion. Skilltrees können als Teil des Puzzles funktionieren, wenn sie den passenden Aufbau für ein konkretes Hindernis erfordern. Und NPCs, die je nach Spielerfahrung oder erworbenem Wissen unterschiedlich reagieren, machen Progression spürbar ohne sie auf Zahlen zu reduzieren.

\subsubsection{ÜBUNG: Spiel Charaktere und Inventar}
\begin{itemize}
    \item Fügt einen Charakter hinzu, der sich in einem der Räume befindet, und modelliert die Möglichkeit, den Charakter von einem Raum in einen anderen zu bewegen.
    \item Erstellt einen weiteren non-player character (NPC) und definiert eine Beziehung zwischen den beiden Charakteren, z.B. Freundschaft oder Feindschaft. 
    \item Hat der Charakter bestimmte Eigenschaften? Wenn, ja erstellt diese als Fakten.
    \item Erstellt ein Inventar für die Spielfigur, in das Objekte aufgenommen und wieder abgelegt werden können (\texttt{nimm(Gegenstand)} und \texttt{lege\_ab(Gegenstand)}).
    \item Verbindet euren Prolog-Code mit dem \textit{STRIPS-Modell}: Schreibt für eure Aktion \texttt{nimm(Gegenstand)} explizit auf, wie die Mengen $\text{pre}(a)$, $\text{del}(a)$ und $\text{add}(a)$ aussehen.
    \item Hat der Charakter bestimmte Fähigkeiten, die gerlent werden können, im Laufe des Spiels? 
    \item Soll der Charakter kämpfen können? Wenn ja, wie könnte eine einfache Kampfmechanik aussehen, die auf Fakten und Regeln basiert?
\end{itemize}

\subsection{Interface und Steuerung}

\textbf{Spielmechanik:} Der Spieler schaut sich um. Das System liest den aktuellen Aufenthaltsort der Spielfigur und gibt eine Beschreibung des Raums sowie die darin liegenden Gegenstände aus.

\textbf{Formalisierung:} Kein Zustandsübergang (keine STRIPS-Aktion), sondern ein lesender Zugriff: $\text{position}(\text{spieler}, R) \land \text{raum}(R)$ liefert den aktuellen Raum; $\text{position}(G, R) \land \text{gegenstand}(G)$ liefert alle Gegenstände darin.

\textbf{Prolog:}

\begin{lstlisting}[language=Prolog, caption={}]
% format/2 wird aus library(format) importiert
:- use_module(library(format)).

sichtbare_umgebung :-
    position(spieler, Raum),
    format("Du befindest dich in: ~w~n", [Raum]),
    forall(
        (position(G, Raum), gegenstand(G)),
        format("  Hier liegt: ~w~n", [G])
    ).

?- sichtbare_umgebung.
Du befindest dich in: eingang
  Hier liegt: laterne
\end{lstlisting}

\subsubsection{Formatierte Ausgabe und sortierte Listen}

\textbf{IF-spezifisches Problem:} Spielende brauchen in jedem Zug eine klare, konsistente Rückmeldung, zum Beispiel für Inventar und Rauminformationen.

\textbf{Formalisierung (Tupel/Logik):} Im Weltmodell $G=(V,E,W,S,\Sigma)$ wird aus dem Zustand $S$ eine geordnete Ansicht berechnet. Logisch bedeutet das, dass aus Relationen wie \texttt{inventar/2} die Terme gesammelt, sortiert und in eine Ausgabe überführt werden.

\textbf{Prolog:} In Scryer Prolog eignet sich \texttt{format/2} für strukturierte Ausgabe und \texttt{msort/2} fuer stabile Sortierung vor der Darstellung.

\begin{lstlisting}[language=Prolog, caption={}]
:- use_module(library(lists)).          % Fuer member/2
:- use_module(library(format)).         % Fuer format/2
% msort/2 ist ebenfalls in library(lists)

in_inventar :-
    findall(G, inventar(spieler, G), Items),
    msort(Items, Sortiert),
    format("Inventar:~n", []),
    forall(member(I, Sortiert), format("  - ~w~n", [I])).

inventar(spieler, schluessel).
inventar(spieler, laterne).
% ?- in_inventar.
\end{lstlisting}

\subsection{Game Loop: Kontrollstrukturen in Scryer Prolog}

\textbf{IF-spezifisches Problem:} Im Institut-Szenario geben Spielende fortlaufend Befehle ein (z. B. \texttt{umsehen}, \texttt{oeffne\_archiv}, \texttt{exit}). Das System soll nach jedem Zug weiterlaufen, bis ein Abbruchkommando kommt.

\textbf{Formalisierung (Tupel/STRIPS/Logik):} Pro Zug wird eine Aktion $a \in \Sigma$ verarbeitet und der Zustand $S$ ggf. aktualisiert. Die Schleife terminiert, sobald ein Stop-Zustand in $S$ markiert ist.

\textbf{Prolog:} In Scryer sind zwei Muster sinnvoll: \texttt{repeat/0} mit \texttt{fail/0} sowie rekursive Wiederholung. \texttt{call/1} dient als Dispatcher, \texttt{once/1} begrenzt unerwünschtes Backtracking bei Ja/Nein-Aktionen.

\texttt{repeat/0} ist ein eingebautes Prädikat, das bei Backtracking unbegrenzt oft \texttt{true} liefert. In Kombination mit \texttt{fail/0}, das immer scheitert, entsteht eine Schleife: \texttt{fail} erzwingt Backtracking, \texttt{repeat} liefert erneut \texttt{true}, der Schleifenkörper wird wieder ausgeführt. Die Schleife bricht erst, wenn eine Bedingung vor \texttt{fail} per Cut verhindert, dass Backtracking einsetzt.

\texttt{call/1} nimmt ein Ziel als Argument und führt es aus. Damit wird ein Befehl wie \texttt{umsehen} zur Laufzeit als Prolog-Term übergeben und als Prädikat aufgerufen -- \texttt{call(dispatch(umsehen))} ist äquivalent  zu \texttt{dispatch(umsehen)}. Das erlaubt es, Benutzereingaben direkt als Ziele zu behandeln, ohne jeden möglichen Befehl vorab explizit aufzuzählen.

\texttt{once/1} führt ein Ziel genau einmal aus und unterdrückt alle weiteren Lösungen durch einen internen Cut. \texttt{once(kann\_oeffnen(archivtuer))} liefert also entweder \texttt{true} oder \texttt{false} -- nie mehrere Antworten. Das ist wichtig, weil \texttt{kann\_oeffnen/1} mehrere Klauseln haben könnte (verschiedene Schlüssel, verschiedene Zustände) und die Ja/Nein-Frage nur eine einzige Antwort erwartet.
\newpage
\begin{lstlisting}[language=Prolog, caption={}]
:- dynamic stop/0.

dispatch(umsehen) :-
    sichtbare_umgebung.

dispatch(oeffne_archiv) :-
    (   once(kann_oeffnen(archivtuer))
    ->  format("Die archivtuer laesst sich oeffnen.~n", [])
    ;   format("Die archivtuer bleibt verschlossen.~n", [])
    ).

dispatch(exit) :-
    assertz(stop),
    format("Spiel beendet.~n", []).

dispatch(_) :-
    format("Unbekannter Befehl im Institut.~n", []).

% Beispiel-Eingabeadapter; spaeter durch Parser ersetzen.
lies_befehl(Befehl) :-
    read(Befehl).

% Muster 1: repeat/fail
command_loop_repeat :-
    retractall(stop),
    repeat,
        lies_befehl(B),
        call(dispatch(B)),
        (   stop
        ->  !
        ;   fail
        ).

% Muster 2: rekursiv
command_loop_rekursiv :-
    retractall(stop),
    loop_rekursiv.

loop_rekursiv :-
    lies_befehl(B),
    call(dispatch(B)),
    (   stop
    ->  true
    ;   loop_rekursiv
    ).
\end{lstlisting}
Fuer den Kurs reichen diese beiden Muster als Kern. \texttt{if\_/3}, Tabling und harte Prozessabbrueche mit \texttt{halt/0} sind Spezialfaelle und werden erst spaeter noetig.

Damit das Spiel auch immersiv und spielbar wird, braucht es ein Interface, über das die Spieler:innen mit der Welt interagieren können. In Prolog erfolgt die Interaktion über Anfragen, die den Zustand der Welt abfragen oder verändern. Die Steuerung umfasst die Mechaniken, die es den Spieler:innen ermöglichen, Aktionen auszuführen und damit die Geschichte voranzutreiben. Später werden wir noch eine Technik kennenlernen, die Prolog uns bietet, um natürlichsprachliche Eingaben zu verarbeiten. Mehr dazu im Abschnitt Definite Clause Grammars (DCG).

\subsubsection{ÜBRUNG Interaktion}
\begin{itemize}
    \item Definiert die Aktionen, die die Spieler:innen ausführen können sollen. 
    \item Was sollen die Aktionen bewirken? Verändern sie den Zustand der Welt, oder geben sie Informationen zurück, oder beides?
    \item Verändern diese sich im Laufe des Spiels?
    \item Modelliert die Bedingungen, unter denen diese Interaktion möglich sind.
\end{itemize}

\subsection{Regeln und Bedingungen}

\textbf{Spielmechanik:} Der Spieler versucht, die Archivtür zu öffnen. Das gelingt nur, wenn er sich im richtigen Raum befindet, den Schlüssel trägt und die Tür noch verschlossen ist.

\textbf{Formalisierung (STRIPS):}
\begin{itemize}
    \item $pre(\text{öffne}(T))$: \texttt{position(spieler, Raum)}, \texttt{tuer(Raum, \_, T)}, \texttt{zustand(T, verschlossen)}, \texttt{inventar(spieler, schluessel)}
    \item $del(\text{öffne}(T))$: \texttt{zustand(T, verschlossen)}
    \item $add(\text{öffne}(T))$: \texttt{zustand(T, offen)}
\end{itemize}

\textbf{Prolog:} Die Vorbedingungen werden als Regelrumpf formuliert.

\textbf{Beispiel:} 

Was muss erfüllt sein, damit eine Tür geöffnet werden kann? Es könnte eine Regel geben, die besagt, dass die Tür nur geöffnet werden kann, wenn der Spieler den Schlüssel besitzt und die Tür nicht verschlossen ist und sich der Spieler in der Nähe der Tür befindet. Diese Regel könnte in Prolog so aussehen:
\begin{lstlisting}[language=Prolog, caption={}]
% position/2, inventar/2 und zustand/2 bereits als dynamic deklariert.
% Spielzustand zu Beginn:
%   position(spieler, korridor).
%   inventar(spieler, schluessel).
%   zustand(archivtuer, verschlossen).

% Regel: generisch fuer jede Tuer T im aktuellen Raum des Spielers
kann_oeffnen(T) :-
    position(spieler, Raum),        % Spieler steht irgendwo
    tuer(Raum, _, T),               % T ist eine Tuer aus diesem Raum
    zustand(T, verschlossen),       % die Tuer ist noch verschlossen
    inventar(spieler, schluessel).  % und der Spieler hat den Schluessel

% Abfrage:
?- kann_oeffnen(archivtuer).
true.
\end{lstlisting}
\textbf{Polystatische Regeln}

Die Regel ist vollständig generisch: \texttt{T} steht für beliebige Türen. Entscheidend ist, dass der aktuelle Raum des Spielers per \texttt{position/2} abgefragt und dann über \texttt{tuer/3} alle Türen aus diesem Raum gesucht werden. Prolog prüft alle vier Teilbedingungen der Reihe nach. Scheitert auch nur eine, gibt Prolog \texttt{false} zurück und das Rätsel bleibt ungelöst. Diese Struktur nennen wir eine \emph{polystatische Regel}: Eine Regel, deren Anwendbarkeit von mehreren unabhängigen Zuständen der Spielwelt abhängt. Je mehr Bedingungen eine Regel hat, desto präziser und anspruchsvoller wird das Rätsel.

\subsubsection{Rätsel und Puzzles}
Auf diese Art können wir auch Rätsel und Puzzles modellieren, die auf logischen Bedingungen basieren. In IF sind Rätsel oft so gestaltet, dass sie bestimmte Informationen oder Gegenstände erfordern, um gelöst zu werden. Diese Anforderungen können als Regeln formuliert werden, die bestimmte Bedingungen erfüllen müssen, damit das Rätsel lösbar ist. Wie können wir die einzelnen Elemente eines Rätsels in Prolog modellieren? Was sind dabei die logischen Eigenheiten und Strukturen von bestimmten Rätseltypen? Was bedeutet es in Prolog, dass ein Rätsel lösbar ist, und wie formulieren wir das in Logik und Prolog?

Unser vorheriges Beispiel haben wir bereits gesehen, dass sich eine Tür nur öffnen lässt, wenn mehrere Bedingungen gleichzeitig erfüllt sind. Wir modellieren zunächst die Spielwelt als Fakten, Inventar, Positionen, Zustände und formulieren die Öffnungsbedingung dann als separate Regel.

\subsubsection{Rätsel mit mehreren Lösungswegen}

\textbf{Spielmechanik:} Das Schloss der Archivtür lässt sich auf verschiedenen Wegen knacken: mit dem richtigen Schlüssel, mit einem Dietrich und Fachwissen, oder mit Sprengstoff -- sofern keine Wache anwesend ist. Welche Option zieht, hängt davon ab, was die Spieler:in bisher gesammelt oder gelernt hat.

\textbf{Formalisierung:} Drei alternative Vorbedingungsmengen für dieselbe Aktion:
\begin{itemize}
    \item $pre_1$: \texttt{inventar(spieler, schluessel)}, \texttt{zustand(T, verschlossen)}
    \item $pre_2$: \texttt{inventar(spieler, dietrich)}, \texttt{wissen(spieler, schlosstechnik)}, \texttt{zustand(T, verschlossen)}
    \item $pre_3$: \texttt{inventar(spieler, sprengstoff)}, $\lnot\,\text{anwesend}(\text{wache})$, \texttt{zustand(T, verschlossen)}
\end{itemize}

\textbf{Prolog:} Jede Alternative wird als eigene Klausel desselben Prädikats modelliert., beispielsweise je nachdem, was die Spieler:in besitzt, weiß oder bereits getan hat. In Prolog lässt sich das durch mehrere Klauseln für dasselbe Prädikat ausdrücken. Prolog probiert sie der Reihe nach durch und verwendet die erste, die erfolgreich beweisbar ist:

\begin{lstlisting}[language=Prolog, caption={}]
% Spielzustand fuer dieses Beispiel:
inventar(spieler, dietrich).
wissen(spieler, schlosstechnik).  % Spieler hat Schlosstechnik erlernt
zustand(archivtuer, verschlossen).
% anwesend(wache) ist KEIN Fakt -> \+ anwesend(wache) ist true (CWA)

% Loesung 1: mit dem richtigen Schluessel
schloss_knackbar(T) :-
    inventar(spieler, schluessel),
    zustand(T, verschlossen).

% Loesung 2: mit Dietrich und Fachwissen
schloss_knackbar(T) :-
    inventar(spieler, dietrich),
    wissen(spieler, schlosstechnik),
    zustand(T, verschlossen).

% Loesung 3: mit Sprengstoff, wenn keine Wache da ist
schloss_knackbar(T) :-
    inventar(spieler, sprengstoff),
    \+ anwesend(wache),  % true, solange kein anwesend(wache)-Fakt existiert
    zustand(T, verschlossen).

?- schloss_knackbar(archivtuer).
true.  % Klausel 2 greift: dietrich + schlosstechnik + verschlossen
\end{lstlisting}

Jede Klausel beschreibt einen anderen Lösungsweg für dasselbe Ziel. Die Spieler:in muss nicht wissen, welche Wege existieren – Prolog findet selbstständig den ersten beweisbaren Pfad. Dieses Prinzip nennt man \emph{polysemantische Rätsel}: Rätsel, deren Lösung kontextabhängig ist und die unterschiedliche Spielverläufe ermöglichen.

\subsubsection{Rätsel als Weggabelungen}

Wenn verschiedene Lösungswege zu verschiedenen Konsequenzen führen, entstehen echte narrative Weggabelungen. Wir können die gewählte Lösungsmethode als Fakt speichern und später auswerten:

\begin{lstlisting}[language=Prolog, caption={}]
:- dynamic geloest_durch/2.

schloss_knacken(T) :-
    inventar(spieler, sprengstoff),
    \+ anwesend(wache),       % ISO: \+ statt not/1
    zustand(T, verschlossen),
    assertz(geloest_durch(T, sprengstoff)).

?- schloss_knacken(archivtuer).
true.

?- geloest_durch(archivtuer, Methode).
Methode = sprengstoff.
\end{lstlisting}

Am Ende des Spiels kann dann abgefragt werden, wie bestimmte Hindernisse überwunden wurden – und daraus folgt ein individuelles Ende. Prolog modelliert damit nicht nur das Rätsel selbst, sondern auch die Geschichte, die durch seine Lösung entsteht. Was zunächst wie formale Logik wirkt, wird zum Instrument für verzweigte, lebendige Erzählungen.

\subsubsection{Analogie: Strukturgleichheit als Verrätselung}

Zwei völlig verschiedene Systeme werden durch die gleiche abstrakte Struktur verbunden. Die Spieler:in muss das erkannte Muster von einer Domäne in eine andere übertragen.

Im Archiv des verlassenen Instituts stehen vier Sockel mit eingravierten Symbolen (Dreieck, Quadrat, Kreis, Stern). Ein Notizzettel, den die Spieler:in im Labor findet, zeigt die Symbole in einer bestimmten Reihenfolge. Die zugehörigen Hebel müssen in dieser Reihenfolge gezogen werden.

\begin{lstlisting}[language=Prolog, caption={}]
% Domane 1: Notizzettel (gefunden im Labor) -- Reihenfolge der Symbole
hinweis(1, dreieck).
hinweis(2, quadrat).
hinweis(3, kreis).
hinweis(4, stern).

% Domane 2: Sockel im Archiv -- jedes Symbol gehoert zu einem Hebel
hebel(hebel_a, dreieck).
hebel(hebel_b, quadrat).
hebel(hebel_c, kreis).
hebel(hebel_d, stern).

% Raetsel: Die Hebel muessen in der Reihenfolge der Symbole
% auf dem Notizzettel gezoegen werden (Analogie: Symbol-Reihenfolge
% auf Hebel-Reihenfolge uebertragen)
raetsel_geloest([H1, H2, H3, H4]) :-
    hinweis(1, S1), hebel(H1, S1),
    hinweis(2, S2), hebel(H2, S2),
    hinweis(3, S3), hebel(H3, S3),
    hinweis(4, S4), hebel(H4, S4).

% Abfrage: Welche Reihenfolge loest das Raetsel?
?- raetsel_geloest(Reihenfolge).
Reihenfolge = [hebel_a, hebel_b, hebel_c, hebel_d].
\end{lstlisting}

Die Analogie liegt darin, dass die Spieler:in die \textit{Struktur} des Notizzettels (Domäne 1) auf die Hebel (Domäne 2) übertragen muss. Der Notiz sagt: \textit{Dreieck zuerst, dann Quadrat, dann Kreis, dann Stern}. Die Hebel haben diese Symbole als Gravur – aber die Spieler:in muss erst beide Informationen \textit{verknüpfen}. Prolog \texttt{findet} durch Unifikation die korrekte Hebel-Reihenfolge, die Spieler:in muss nur die Analogie \emph{erkennen}. Dieses Muster funktioniert mit beliebigen Zuordnungen – wäre der Notizzettel \texttt{hinweis(1, quadrat)}, würde Prolog automatisch \texttt{hebel\_b} an erster Stelle liefern.

\subsubsection{Symbol mit Mehrfachcodierung: Funktion und Bedeutung}

Ein Objekt hat gleichzeitig eine mechanische Funktion und eine tiefere symbolische Bedeutung. Die Lösung erfordert, beide Ebenen zu verstehen.

\begin{lstlisting}[language=Prolog, caption={}]
% inventar/2 und zustand/2 bereits als dynamic deklariert (s. Weltblock).
:- dynamic dahinter/1, geoeffnet/1.

% Funktionale Ebene: Der verrostete Schluessel im Korridor oeffnet das Archiv
oeffnet(schluessel, archivtuer).

% Symbolische Ebene: Das Archiv steht fuer verdraengte Vergangenheit
bedeutung(schluessel, zugang_zum_vergangenen).
bedeutung(archivtuer, grenze_des_wissens).

% Spielaktion: Tuer oeffnen
% pre(a): inventar(spieler, schluessel), zustand(archivtuer, verschlossen)
% del(a): zustand(archivtuer, verschlossen)
% add(a): zustand(archivtuer, offen), geoeffnet(archivtuer), dahinter(akte)
archiv_oeffnen :-
    inventar(spieler, schluessel),
    zustand(archivtuer, verschlossen),
    oeffnet(schluessel, archivtuer),
    retract(zustand(archivtuer, verschlossen)),
    assertz(zustand(archivtuer, offen)),
    assertz(geoeffnet(archivtuer)),
    assertz(dahinter(akte)).

% Abfrage der symbolischen Ebene
einsicht_erlangt :-
    geoeffnet(archivtuer),
    dahinter(akte),
    bedeutung(schluessel, zugang_zum_vergangenen).

?- archiv_oeffnen.
true.

?- einsicht_erlangt.
true.

% Konsistenzcheck: Die Tuer ist jetzt offen, nicht mehr verschlossen
?- zustand(archivtuer, verschlossen).
false.
?- zustand(archivtuer, offen).
true.
\end{lstlisting}
Hier zeigt sich das Design-Prinzip: Der Schlüssel ist funktional (öffnet die Tür) und symbolisch (gewährt Zugang zum Vergangenen). \texttt{retract} stellt sicher, dass der alte Zustand vollständig entfernt wird, bevor \texttt{assertz} den neuen einträgt. Der Zustand der Tür ist damit zu jedem Zeitpunkt konsistent.

\subsubsection{Synekdoche: Teil impliziert Ganzes}

Ein kleines Detail oder Fragment lässt die Spieler:in auf ein größeres System schließen. Prolog nutzt hier \texttt{findall/3}, um verteilte Information zu rekonstruieren.
\begin{lstlisting}[language=Prolog, caption={}]
% Ein verlassenes Zimmer mit Fragmenten
fragment(seite_1, 'Namen einer Liste').
fragment(seite_2, 'Daten und Orte').
fragment(seite_3, 'Aufzeichnungen von Experimenten').
fragment(seite_4, 'Evakuierungsbefehl').

% Die Synekdoche: Einzelne Fragmente implizieren ein System
gesamtbild_erkannt :-
    findall(F, fragment(_, F), AlleFragmente),
    length(AlleFragmente, N),
    N >= 3,  % Mindestens 3 Fragmente noetig zum Erkennen
    write('Ein grosses, dunkles Projekt wird sichtbar: Menschenversuche und Vertuschung.').

% Abfrage: Hat die Spieler:in genug Fragmente gefunden?
?- gesamtbild_erkannt.
Ein grosses, dunkles Projekt wird sichtbar: Menschenversuche und Vertuschung.
true.
\end{lstlisting}
Die Kraft der Synekdoche in Prolog liegt darin, dass einzelne Fakten automatisch eine \textit{interpretierbare Gesamtheit} bilden, sobald genug davon gesammelt sind.

\subsubsection{Metarätsel: Einzelne Rätsel ergeben ein übergeordnetes Muster}

Mehrere kleine Rätsel sind nur zusammen bedeutsam. Ihre Lösungen kombiniert offenbaren ein größeres Geheimnis.
\begin{lstlisting}[language=Prolog, caption={}]
:- dynamic geloest/1.

% Raetsel 1: Der Code
raetsel_code(X) :- X = 1924.

% Raetsel 2: Der Name
raetsel_name(X) :- X = 'Dr. Voss'.

% Raetsel 3: Der Ort
raetsel_ort(X) :- X = 'Haus der Wissenschaft'.

% Metaraetsel: Erst wenn alle drei geloest sind,
% wird das wahre Geheimnis offenbar
wahrheit_erkannt :-
    geloest(code),
    geloest(name),
    geloest(ort),
    write('Alle Spuren fuehren zu einem verborgenen Forschungsprojekt: Diese Person traegt Verantwortung.').

% Spieler loest die einzelnen Raetsel
?- raetsel_code(X), assertz(geloest(code)), write('Code: '), write(X), nl.
Code: 1924
true.

?- raetsel_name(X), assertz(geloest(name)), write('Name: '), write(X), nl.
Name: Dr. Voss
true.

?- raetsel_ort(X), assertz(geloest(ort)), write('Ort: '), write(X), nl.
Ort: Haus der Wissenschaft
true.

% Nur wenn alle geloest sind:
?- wahrheit_erkannt.
Alle Spuren fuehren zu einem verborgenen Forschungsprojekt: Diese Person traegt Verantwortung.
true.
\end{lstlisting}
Das Metarätsel zeigt die Kraft der logischen Verknüpfung: Einzelne Lösungen sind unterhaltsam, aber ihre Kombination erzeugt Bedeutung. Für die Spieler:in entsteht so erst nachträglich die Erkenntnis, dass die Teile zusammenpassen – genau das ist Drama durch Struktur.

\subsubsection{Rekursion}

\textbf{Spielmechanik:} Ein Schlüssel liegt in einer Schatulle, die in einer Kiste steckt, die im Lagerraum steht. Der Spieler fragt: „Ist der Schlüssel irgendwo im Lagerraum?" -- was über mehrere Ebenen hinweg bejaht werden muss.

\textbf{Formalisierung:} Die Relation $\text{befindet\_sich\_in}(X, Y)$ ist rekursiv definiert:
$$\text{befindet\_sich\_in}(X, Y) \;\leftrightarrow\; \text{enthalten\_in}(X, Y) \;\lor\; \exists Z\colon \text{enthalten\_in}(X, Z) \land \text{befindet\_sich\_in}(Z, Y)$$

\textbf{Prolog:} Zwei Klauseln -- Basisfall und rekursiver Schritt.
\begin{lstlisting}[language=Prolog, caption={}]
% Direkte Enthaltensein-Beziehung
enthalten_in(schluessel, schatulle).
enthalten_in(schatulle, kiste).
enthalten_in(kiste, lagerraum).

% Rekursive Regel: X ist irgendwo in Y enthalten,
% entweder direkt oder ueber ein Zwischenobjekt
befindet_sich_in(X, Y) :-
    enthalten_in(X, Y).

befindet_sich_in(X, Y) :-
    enthalten_in(X, Z),
    befindet_sich_in(Z, Y).

?- befindet_sich_in(schluessel, kiste).
true.

?- befindet_sich_in(schluessel, lagerraum).
true.
\end{lstlisting}
Die erste Klausel von \texttt{befindet\_sich\_in/2} beantwortet den direkten Fall. Die zweite Klausel sucht ein Zwischenobjekt \texttt{Z}: Wenn \texttt{X} direkt in \texttt{Z} liegt und \texttt{Z} wiederum irgendwo in \texttt{Y}, dann liegt \texttt{X} auch in \texttt{Y}. Prolog wendet diese Regel so oft an, wie nötig, bis entweder ein direkter Fakt gefunden wird oder keine weiteren Alternativen mehr offen sind. So kann dieselbe Regel für zwei Ebenen und für zwanzig Ebenen funktionieren, ohne dass der Code sich ändert.

\subsubsection{Listen}

\textbf{Spielmechanik:} Der Spieler trägt mehrere Gegenstände bei sich. Er kann neue aufnehmen, vorhandene ablegen und das Inventar anzeigen lassen. Die Kapazität kann begrenzt sein.

\textbf{Formalisierung:} Das Inventar ist eine geordnete Menge $I = [g_1, g_2, \ldots, g_n]$. Aufnehmen: $I' = I \cup \{g\}$ (append). Ablegen: $I' = I \setminus \{g\}$ (select). Kapazitätsprüfung: $|I| < \text{max}$.

\textbf{Prolog:} Listen werden durch eckige Klammern notiert und können beliebig viele Elemente enthalten: \texttt{[schluessel, laterne, akte]}. Das erste Element heißt \emph{Kopf} (Head), der Rest \emph{Schwanz} (Tail). Prolog kann beides direkt durch Pattern Matching trennen: \texttt{[H|T]} bindet \texttt{H} an den Kopf und \texttt{T} an den Rest.

\textbf{Fakten oder Liste: zwei Alternativen für das Inventar}

Ein Inventar lässt sich auf zwei Arten modellieren. Mit einzelnen Fakten:

\begin{lstlisting}[language=Prolog, caption={}]
hat(spieler, schluessel).
hat(spieler, laterne).
hat(spieler, akte).
\end{lstlisting}
Oder als Liste in einem einzigen Fakt:

\begin{lstlisting}[language=Prolog, caption={}]
:- dynamic inventar/2.
inventar(spieler, [schluessel, laterne, akte]).
\end{lstlisting}
Einzelne Fakten sind einfach abzufragen und mit \texttt{assertz}/\texttt{retract} zu ändern. Listen sind besser geeignet, wenn Reihenfolge zählt, wenn die ganze Sammlung auf einmal übergeben werden soll oder wenn Listenoperationen wie die Länge der Liste, die Anzahl der Elemente oder das Durchsuchen von Listen nötig sind.

\textbf{Gegenstand im Inventar prüfen: \texttt{member/2}}

\begin{lstlisting}[language=Prolog, caption={}]
inventar(spieler, [schluessel, laterne, akte]).

hat_gegenstand(Spieler, Ding) :-
    inventar(Spieler, Liste),
    member(Ding, Liste).

?- hat_gegenstand(spieler, laterne).
true.

?- hat_gegenstand(spieler, dietrich).  % nicht im Inventar
false.
\end{lstlisting}

\textbf{Gegenstand aufnehmen: \texttt{append/3}}
\begin{lstlisting}[language=Prolog, caption={}]
aufnehmen(Spieler, Ding) :-
    inventar(Spieler, Alt),
    \+ member(Ding, Alt),
    append(Alt, [Ding], Neu),
    retract(inventar(Spieler, Alt)),
    assertz(inventar(Spieler, Neu)).

?- aufnehmen(spieler, dietrich).
true.

?- inventar(spieler, X).
X = [schluessel, laterne, akte, dietrich].
\end{lstlisting}

\textbf{Gegenstand ablegen: \texttt{select/3}}

\begin{lstlisting}[language=Prolog, caption={}]
:- use_module(library(lists)).

ablegen(Spieler, Ding) :-
    inventar(Spieler, Alt),
    select(Ding, Alt, Neu),
    retract(inventar(Spieler, Alt)),
    assertz(inventar(Spieler, Neu)).

?- ablegen(spieler, laterne).
true.

?- inventar(spieler, X).
X = [schluessel, akte, dietrich].
\end{lstlisting}
\textbf{Inventargröße prüfen: \texttt{length/2}}

Nützlich zum Beispiel, wenn eine Truhe oder ein Rucksack nur eine begrenzte Anzahl von Gegenständen tragen kann:
\begin{lstlisting}[language=Prolog, caption={}]
max_inventar(5).

inventar_voll(Spieler) :-
    inventar(Spieler, Liste),
    length(Liste, N),
    max_inventar(Max),
    N >= Max.

?- inventar_voll(spieler).
false.
\end{lstlisting}

\textbf{Besuchte Räume protokollieren}

Listen eignen sich auch, um den Weg des Spielers festzuhalten, etwa um Endnachrichten oder Rätselzustände von der Spielhistorie abhängig zu machen:
\begin{lstlisting}[language=Prolog, caption={}]
:- use_module(library(lists)).

:- dynamic besucht/2.
besucht(spieler, []).

raum_betreten(Spieler, Raum) :-
    besucht(Spieler, Alt),
    \+ member(Raum, Alt),
    append(Alt, [Raum], Neu),
    retract(besucht(Spieler, Alt)),
    assertz(besucht(Spieler, Neu)).

?- raum_betreten(spieler, eingang).
?- raum_betreten(spieler, korridor).
?- besucht(spieler, X).
X = [eingang, korridor].
\end{lstlisting}

\subsubsection{Sequenz-Rätsel: Listen als Zustandsgedächtnis}

Viele Rätsel in IF verlangen eine bestimmte \textit{Reihenfolge} von Aktionen – etwa das Ziehen von Hebeln in der korrekten Abfolge. Listen eignen sich hervorragend, um solche Sequenzen zu protokollieren und zu prüfen.

\textbf{Spielmechanik:} Ein Regal mit vier nummerierten Ordnern. Der Spieler muss sie in einer bestimmten Reihenfolge ziehen. Jeder Zug wird in der Wissensbasis gespeichert. Stimmt die Reihenfolge nicht, wird zurückgesetzt.

\begin{lstlisting}[language=Prolog, caption={Sequenz-Rätsel: Ordner in richtiger Reihenfolge ziehen}]
:- use_module(library(lists)).

:- dynamic gezogen/1.

% Die korrekte Reihenfolge (z.B. 17, 13, 7, 11)
regal_loesung(['17', '13', '7', '11']).

% Einen Ordner ziehen
entferne_ordner(Nr) :-
    assertz(gezogen(Nr)),
    findall(G, gezogen(G), Gezogen),
    regal_loesung(Loesung),
    length(Gezogen, Len),
    length(Praefix, Len),
    append(Praefix, _, Loesung),
    (   Praefix == Gezogen
    ->  (   Gezogen == Loesung
        ->  format("Regal geoeffnet!~n", [])
        ;   format("Klick. Der Ordner sitzt locker.~n", [])
        )
    ;   retractall(gezogen(_)),
        format("Falsche Reihenfolge! Die Ordner schnappen zurueck.~n", [])
    ).
\end{lstlisting}

Der Kern der Prüfung: \texttt{append(Praefix, \_, Loesung)} zerlegt die Lösung in einen Prefix der Länge der bereits gezogenen Ordner und den Rest. Stimmt der Prefix mit den gezogenen Ordnern überein, ist die Sequenz bisher korrekt. Ist die gesamte Lösung erreicht, ist das Rätsel gelöst. Bei Abweichung wird der gesamte Fortschritt zurückgesetzt (\texttt{retractall(gezogen(\_))}).

Dieses Muster – eine Liste von bisherigen Aktionen, eine Soll-Liste als Lösung und Prefix-Vergleich – lässt sich auf beliebige Sequenz-Rätsel übertragen: Hebel-Schalter, Kombinationsschlösser, Tanzschritte, Zaubersprüche.

\subsection{Prolog und der Zufall}

Da Prolog auf formaler Logik basiert liefert ein Prädikat mit denselben Argumenten grundsätzlich immer dasselbe Ergebnis. Zufall scheint damit zunächst ein Widerspruch zu sein, doch nicht nur in interaktive Fiction benötigen wir unvorhersehbare Ereignisse, zufällige Begegnungen oder zufällig wechselnde Ereignisse.

\subsubsection{Wie funktioniert Zufall im Computer?}

Echten Zufall kann ein Computer nicht erzeugen. Stattdessen verwenden Programme \emph{Pseudozufallszahlen}, eine mathematische Formel berechnet aus einem Startwert (\emph{Seed}) eine Folge von Zahlen, die zufällig \emph{wirken}, es aber nicht sind. Bei jedem Programmstart wird ein neuer Seed gesetzt – meist aus der aktuellen Uhrzeit – sodass die Folge praktisch nicht vorhersagbar ist. In Prolog erzeugt \texttt{random/1} eine Gleitkommazahl zwischen 0 und 1, \texttt{random\_between/3} eine ganze Zahl in einem definierten Bereich:
\begin{lstlisting}[language=Prolog, caption={}]
:- use_module(library(random)).

?- random(X).
X = 0.7362...

?- random_between(1, 6, X).   % Wuerfelwurf
X = 4.
\end{lstlisting}
\subsubsection{Zufällige Auswahl aus einer Liste}

Das Prädikat \texttt{random\_member/2} wählt bei jedem Aufruf ein zufälliges Element aus einer Liste. Es ist damit \emph{nicht deterministisch im logischen Sinne} , das dasselbe Ziel unterschiedliche Ergebnisse liefern kann und bildet eine bewusste Ausnahme vom streng logischen Modell Prologs:
\begin{lstlisting}[language=Prolog, caption={}]
:- use_module(library(random)).

zufaellige_begruessung(Begruessung) :-
    random_member(Begruessung, [
        "Der Wind fluestert deinen Namen.",
        "Stille. Dann ein Knacken im Gebusch.",
        "Ein Rabe beobachtet dich von oben."
    ]).
\end{lstlisting}

Jeder Aufruf von \texttt{zufaellige\_begruessung/1} liefert einen anderen Einstiegssatz – die Geschichte fühlt sich lebendig an, obwohl dieselbe Wissensbasis verwendet wird.

\subsubsection{Zufall mit Bedingungen kombinieren}

Zufall und Logik lassen sich verbinden: \texttt{findall/3} sammelt alle zufälligen Elemente, die zusätzlich eine Bedingung erfüllen. So entstehen gefilterte Zufallsauswahlen, etwa nur Gegenstände, die der Charakter tragen kann:
\begin{lstlisting}[language=Prolog, caption={}]
tragbar(schluessel).
tragbar(laterne).

% Zufaellig einen tragbaren Gegenstand aus einer Fundliste auswaehlen
zufaelliger_gegenstand(X) :-
    findall(G, (random_member(G, [schluessel, akte, laterne, staub]), tragbar(G)), Liste),
    random_member(X, Liste).
\end{lstlisting}
\subsubsection{Listen mischen mit \texttt{random\_permutation/2}}

Um eine zufällige Reihenfolge zu erzeugen, zum Beispiel für Ereigniskarten oder Begegnungstabellen, mischt \texttt{random\_permutation/2} alle Elemente einer Liste neu:

\begin{lstlisting}[language=Prolog, caption={}]
?- random_permutation([schluessel, laterne, akte, staub], Gemischt).
Gemischt = [laterne, staub, schluessel, akte].
\end{lstlisting}
\medskip
\noindent\textbf{Für das Design wichtig:} Zufall in Prolog ist ein pragmatisches Werkzeug, das den logischen Rahmen gezielt durchbricht. Als Designer:in entscheidest du, \emph{wo} Unvorhersehbarkeit die Erfahrung bereichert und \emph{wo} Konsistenz der Welt Glaubwürdigkeit verleiht.

\subsection{Lange Texte einlesen und ausgeben}

Um längere Texte, wie Raum- oder Charakterbeschreibungen, in Prolog zu verwalten, ist es oft praktischer, sie in externen Dateien zu speichern und bei Bedarf einzulesen. Das folgende Beispiel zeigt, wie man eine Textdatei öffnet, Zeile für Zeile liest und den Inhalt auf dem Bildschirm ausgibt. Am Ende wird die Datei geschlossen.
\begin{lstlisting}[language=Prolog, caption={}]
text_aus_datei(Datei) :-
    open(Datei, read, Datenstrom),
    repeat,
        read_line_to_codes(Datenstrom, Inhalt),
        (   Inhalt = end_of_file
        ->  !
        ;   atom_codes(Atom, Inhalt), write(Atom), nl, fail
        ),
    close(Datenstrom).
\end{lstlisting}

\subsubsection{Eingabe als Daten: Codes, Atome, Zahlen}

\textbf{IF-spezifisches Problem:} Spielbefehle kommen als freie Texteingabe und muessen in verarbeitbare Aktionen ueberfuehrt werden.

\textbf{Formalisierung (Tupel/STRIPS/Logik):} Die Eingabe ist zunaechst Teil von $\Sigma$ (moegliche Aktionen) als Zeichenfolge. Erst durch logische Zerlegung entstehen Symbole, die in Vorbedingungen und Effekten von Aktionen weiterverwendet werden koennen.

\textbf{Prolog:} In Scryer Prolog ist \texttt{read\_line\_to\_codes/2} ein robuster Einstieg, weil Texte zuerst als Codes gelesen und dann mit \texttt{atom\_codes/2}, \texttt{number\_codes/2} und \texttt{sub\_atom/5} schrittweise strukturiert werden.
\begin{lstlisting}[language=Prolog, caption={}]
lies_befehl_als_atom(Datei, Befehl) :-
    open(Datei, read, S),
    read_line_to_codes(S, Cs),
    close(S),
    atom_codes(Befehl, Cs).

token_ist_zahl(TokenAtom, N) :-
    atom_codes(TokenAtom, Cs),
    number_codes(N, Cs).

kurzform(Hauptbegriff, Eingabe) :-
    sub_atom(Hauptbegriff, 0, 3, _, Eingabe).

% ?- atom_codes(archivtuer, Cs).
% ?- kurzform(archivtuer, arc).
% ?- number_codes(X, "2").
\end{lstlisting}
So wird klar, wie aus Rohtext parserfähige Strukturen entstehen. Erst danach folgt mit DCG die Grammatikebene.

\subsection{Sprache und Parser}

Wie lässt sich menschliche Sprache so zerlegen, dass eine Maschine diese verarbeiten kann? Ein Parser ist ein Programm, das eine Texteingabe liest und sie gemäß den Regeln einer formalen Grammatik analysiert und in eine strukturierte Darstellung überführt. Das Ergebnis ist typischerweise ein Syntaxbaum oder eine Folge semantischer Aktionen, die von nachgelagerten Systemkomponenten weiterverarbeitet werden. Parser sind ein fundamentales Konzept der theoretischen Informatik und Computerlinguistik und bilden die Grundlage von Compilern und Interpretern und in jüngeren Entwicklungen, von neuronalen Sprachmodellen, deren Vorstufenarchitekturen auf grammatikbasierten Formalismen aufbauen.

Für Interactive Fiction ist der Parser das Interface, da sämtliche Interaktion mit Hilfe von Sprache erfolgt: \texttt{UNLOCK DOOR WITH IRON KEY}, \texttt{ASK WIZARD ABOUT MAP}. Das System muss aus solchen Eingaben den Verb-Objekt-Rahmen rekonstruieren, Bezüge auflösen, Mehrdeutigkeiten klären und den Satz auf Operationen abbilden. Nick Montfort beschreibt den Parser deshalb als das eigentliche epistemische Interface der Form: „The parser is the mechanism that creates the illusion of a world that responds to language" \citep{Montfort2003}. Die Konsequenz für das Design ist, dass jede Grenze des Parsers gleichzeitig auch eine Grenze der erfahrbaren Welt bedeutet. Daher 

Die historischen Lösungen für dieses Problem zeigen, wie stark die technische Architektur das Spielerlebnis formt. Der Parser von \textit{Colossal Cave Adventure} (Crowther/Woods, 1976) funktionierte als simples Lexikon-Lookup bei dem ein Verb und Objekt mit einer festen Tabelle abgeglichen wurde. \textit{Infocom} entwickelte für \textit{Zork} (1980) mit dem \textit{ZIL}-Compiler (\textit{Zork Implementation Language}) und der \textit{Z-Machine} eine deutlich differenziertere Architektur. Präpositionalphrasen, Pronomen, mehrdeutige Nomen konnten verarbeitet werden, mit Hilfe eines prozeduralen Paradigmas (MDL, einer LISP-Variante), das die Grammatiklogik als eindeutige Zustandsübergänge im Code definierte. \textit{Magnetic Scrolls} ging für \textit{The Pawn} (1985) noch weiter und implementierte in C einen Parser, der komplexere Satzstrukturen verstand als jeder zeitgenössische Konkurrent. Dabei setzte er stillschweigend voraus, dass Spieler:innen differenzierter formulieren würden, als es die Konventionen des Genres bis dahin erwartet hatten.

Ein besserer Parser verändert also nicht nur die technischen Möglichkeiten, sondern die Erwartungshaltung der Nutzenden. In imperativen Sprachen wie C, FORTRAN oder Assembler muss man jeden Schritt des Parsens von Hand programmieren, dadurch steckt die Grammatik im Programmablauf und liegt nicht als klar lesbare Regelbeschreibung vor. Das macht sie bei wachsender Komplexität schwer wartbar und kaum erweiterbar. Parsergeneratoren wie \textit{ANTLR} oder \textit{Yacc} lösen das durch eine Trennung von Grammatikspezifikation und ausgeführtem Code, ein Schritt in Richtung Lesbarkeit, der aber nach wie vor imperativen Maschinencode erzeugt, der die eigentliche Grammatikstruktur nicht mehr unmittelbar sichtbar macht.

Prolog nimmt gegenüber diesen Ansätzen eine andere Position ein. Weil Prolog-Klauseln selbst Relationen zwischen Termen beschreiben und die Auswertung als kontrollierter Beweisprozess mit Backtracking organisiert ist, lassen sich Grammatikregeln direkt als Prolog-Klauseln formulieren, ohne das ein separater Parser benötigt wird. Ein in \textbf{DCG-Notation} (\textit{Definite Clause Grammars}) geschriebenes Prolog-Programm ist gleichzeitig Grammatikspezifikation und ausführbarer Parser. Diese Eigenschaft von Prolog kommt nicht von ungefähr, denn Colmerauer, der Entwickler von Prolog, arbeitete ursprünglich an Systemen zur maschinellen Sprachverarbeitung, und die Verbindung zwischen logischer Programmierung und formaler Grammatik ist damit keine nachträgliche Anwendung, sondern historisch in der Entstehung der Sprache selbst verankert \citep{Colmerauer1996}. 

Für IF bedeutet das: Parser und Weltmodell lassen sich in derselben Sprache und derselben Denkform implementieren, ohne Kontextwechsel zwischen Paradigmen. Prolog stellt für diese Aufgabe eine eigene, kompakte Notation bereit, die bereits erwähnten \textbf{Definite Clause Grammars}.

\subsection{Definite Clause Grammars (DCG)}

\textbf{Definite Clause Grammars (DCGs)} sind eine Prolog-Notation für formale Grammatiken. Sie wurden eingeführt, um Sprachverarbeitung in Logikprogrammen präzise und zugleich lesbar zu machen. Statt Parsing als lange Folge von Listenoperationen zu implementieren, formulieren wir direkt Regeln der Form \textit{»Ein Befehl besteht aus Verb und Objekt«}.

Technisch sind DCGs \textit{Syntactic Sugar} für Prädikate mit Differenzlisten. Eine Regel wie
\begin{lstlisting}[language=Prolog, caption={}]
befehl --> verb, objekt.
\end{lstlisting}
wird intern in ein gewöhnliches Prolog-Prädikat mit zwei zusätzlichen Argumenten übersetzt (Eingabe und Rest).

Für unser natürlichsprachliches Textinterface ist das zentral: Wir können Spielereingaben parsen und sofort in interne Aktionen überführen.

\begin{lstlisting}[language=Prolog,caption={}]
:- use_module(library(dcgs)).

kommando(gehe(Richtung)) --> [geh, Richtung].
kommando(nimm(Objekt))   --> [nimm, Objekt].

parse_befehl(Woerter, Aktion) :-
    phrase(kommando(Aktion), Woerter).
\end{lstlisting}

Beispiel:
\begin{lstlisting}[language=Prolog, caption={}]
?- parse_befehl([nimm, laterne], A).
A = nimm(laterne).
\end{lstlisting}

DCGs können nicht nur analysieren, sondern auch Sequenzen erzeugen. Damit lassen sich auch ASCII-Ausgaben wie Kartenzeilen, Muster oder Labyrinthteile generieren.

\begin{lstlisting}[language=Prolog,caption={}]
:- use_module(library(dcgs)).

wand(0) --> [].
wand(N) --> ['#'], {N > 0, N1 is N - 1}, wand(N1).

ascii_zeile(N, Zeichen) :-
    phrase(wand(N), Zeichen).
\end{lstlisting}

Beispiel:
\begin{lstlisting}[language=Prolog, caption={}]
?- ascii_zeile(5, Z).
Z = [#,#,#,#,#].
\end{lstlisting}

Hinweis: In Scryer Prolog sind DCGs nicht automatisch aktiv. Ohne \texttt{:- use\_module(library(dcgs)).} funktionieren \texttt{-->} und \texttt{phrase/2} nicht. Zudem ist das Beispiel mit \texttt{is/2} und \texttt{>/2} nicht voll relational, es ist primär auf die Richtung \texttt{N -> Zeichenliste} ausgelegt.

Damit sind DCGs in IF doppelt nützlich, einmal als Parser für Spracheingaben und als Generator für strukturierte Textausgaben.

\subsubsection{Die vollständige Pipeline: von der Eingabe zur Aktion}

In einem spielbaren Textadventure durchläuft jede Benutzereingabe mehrere Verarbeitungsschichten. Am Beispiel \texttt{"gehe nach labor"} zeigen wir den gesamten Weg:

\begin{enumerate}
    \item \textbf{Zeichen einlesen:} \texttt{get\_line\_to\_chars/2} liefert die Eingabe als Liste von Character-Codes.
    \item \textbf{Tokenisierung:} \texttt{zeile\_zu\_woertern/2} zerlegt die Zeichenliste in einzelne Wörter (Atome).
    \item \textbf{Normalisierung:} Satzzeichen am Wortende werden entfernt.
    \item \textbf{DCG-Parsing:} \texttt{phrase/2} vergleicht die Tokenliste mit den Grammatikregeln und erzeugt eine strukturierte Aktion.
    \item \textbf{Dispatch:} Die Aktion wird als Prolog-Ziel ausgeführt.
\end{enumerate}

\begin{lstlisting}[language=Prolog, caption={Tokenisierung einer Eingabezeile}]
% Aus "gehe nach labor" wird die Liste [gehe, nach, labor].
% Aus "nimm die id_karte!" wird [nimm, die, id_karte].

leerzeichen_ueberspringen([Z|R], E) :-
    ist_leerzeichen(Z), !,
    leerzeichen_ueberspringen(R, E).
leerzeichen_ueberspringen(R, R).

wort_lesen([], [], []) :- !.
wort_lesen([Z|R], [], [Z|R]) :-
    ist_leerzeichen(Z), !.
wort_lesen([Z|R], [Z|W], E) :-
    wort_lesen(R, W, E).

zeile_zu_woertern([], []) :- !.
zeile_zu_woertern(Zeichen, Woerter) :-
    leerzeichen_ueberspringen(Zeichen, Rest1),
    (   Rest1 = []
    ->  Woerter = []
    ;   wort_lesen(Rest1, WZ, Rest2),
        atom_chars(Wort, WZ),
        zeile_zu_woertern(Rest2, Weitere),
        Woerter = [Wort|Weitere]
    ).
\end{lstlisting}

Die Normalisierung entfernt Satzzeichen, die als Teil eines Wortes eingelesen wurden:

\begin{lstlisting}[language=Prolog, caption={Normalisierung: Satzzeichen entfernen}]
entferne_abschluss(In, Out) :-
    append(Vorne, [Letztes], In),
    member(Letztes, ['.', ',', '!', '?']), !,
    entferne_abschluss(Vorne, Out).
entferne_abschluss(In, In).

normalisiere_token(Roh, Norm) :-
    atom_chars(Roh, Zeichen),
    entferne_abschluss(Zeichen, Bereinigt),
    (   Bereinigt = []
    ->  Norm = Roh
    ;   atom_chars(Norm, Bereinigt)
    ).
\end{lstlisting}

Erst dann kommt der DCG-Parser zum Einsatz. Ein vollständiger Parser für ein Textadventure erlaubt mehrere Schreibweisen für denselben Befehl:

\begin{lstlisting}[language=Prolog, caption={DCG mit Synonymen und Artikeln}]
:- use_module(library(dcgs)).

befehl(gehe(R))        --> [geh,   R].
befehl(gehe(R))        --> [gehe,  R].
befehl(gehe(R))        --> [gehe,  nach, R].

befehl(nimm(O))        --> [nimm,  O].
befehl(nimm(O))        --> [nimm,  die, O].
befehl(nimm(O))        --> [nimm,  den, O].

befehl(untersuche(O))  --> [untersuche, O].
befehl(untersuche(O))  --> [u,           O].

befehl(inventar)       --> [inventar].
befehl(inventar)       --> [i].

befehl(umsehen)        --> [umsehen].
befehl(umsehen)        --> [schau].

parse_befehl(Woerter, Aktion) :-
    phrase(befehl(Aktion), Woerter).
\end{lstlisting}

Prolog probiert dank Backtracking automatisch alle Alternativen durch – schlägt \texttt{[gehe, nach, labor]} fehl, weil nur zwei Tokens übergeben wurden, versucht es die nächste Klausel \texttt{[gehe, R]}. Der Parser ist damit gleichzeitig Spezifikation und ausführbarer Code.

Die Integration in den Game-Loop verbindet alle Schichten:

\begin{lstlisting}[language=Prolog, caption={Komplette Pipeline im Game-Loop}]
lies_befehl(Befehl) :-
    format("~n> ", []),
    get_line_to_chars(user_input, Zeichen, []),
    zeile_zu_woertern(Zeichen, Woerter),
    maplist(normalisiere_token, Woerter, Norm),
    (   parse_befehl(Norm, Befehl)
    ->  true
    ;   format("Unbekannter Befehl.~n", []),
        Befehl = unbekannt
    ).

loop_rekursiv :-
    lies_befehl(Befehl),
    dispatch(Befehl),
    (   stop -> true
    ;   loop_rekursiv
    ).
\end{lstlisting}

Damit wird aus der rohen Tastatureingabe \texttt{"nimm die id\_karte!"} Schritt für Schritt die Aktion \texttt{nimm(id\_karte)}. Jede Schicht lässt sich unabhängig testen und erweitern.

\subsection{Metaprogrammierung und Selbstreferenzialität}

Metaprogrammierung bezeichnet die Fähigkeit eines Programms, sich selbst zu analysieren, zu verändern oder sogar zu generieren. In Prolog ist das besonders mächtig, da die Sprache selbst auf Logik basiert und Regeln als Datenstrukturen behandelt werden können.

\subsubsection{Beispiel: Ein Meta-Interpreter}

\textbf{Metaprogrammierung} bedeutet hier, dass wir nicht nur über die Spielwelt sprechen, sondern auch über die \textit{Regeln}, mit denen die Spielwelt berechnet wird. Ein einfacher Meta-Interpreter in Prolog.
\begin{lstlisting}[caption={}]
solve(true).
solve((A,B)) :- solve(A), solve(B).
solve(A) :- clause(A, B), solve(B).

raum(eingang).
raum(korridor).
verbindet(eingang, korridor).
kann_gehen(X, Y) :- raum(X), raum(Y), verbindet(X, Y).

% Anfrage:
% ?- solve(kann_gehen(eingang, korridor)).
\end{lstlisting}

\texttt{solve/1} ist ein\textit{ Interpreter im Interpreter}. Statt ein Ziel direkt von Prolog lösen zu lassen, lassen wir unser eigenes Prädikat prüfen, ob es aus den vorhandenen Regeln folgt. Mit \texttt{clause/2} liest das Programm eine Regel aus der Wissensbasis und verarbeitet sie weiter. So wird sichtbar, wie Inferenz funktioniert.

\subsubsection{Wann und warum macht man das in Prolog?}

\begin{itemize}
    \item \textbf{Erklärbarkeit}: Wenn wir zeigen wollen, warum eine Schlussfolgerung zustande kommt.
    \item \textbf{Analyse und Anpassung}: Wenn Regeln zur Laufzeit untersucht oder gesteuert werden sollen, z. B. in einem adaptiven IF-System.
    \item \textbf{Eigene Inferenzstrategien}: Wenn man testen will, wie sich andere Such- oder Priorisierungsstrategien auf das Verhalten auswirken.
\end{itemize}

Das ermöglicht es, dass ein Prolog-Programm nicht nur die Spielwelt modelliert, sondern auch die Regeln des Spiels selbst dynamisch anpassen oder erweitern kann.

\subsection{Fortgeschrittene Spielmechaniken}

\subsubsection{Modularisierung, Savegame und testbare Ausgabe}

\textbf{IF-spezifisches Problem:} Groessere Spiele brauchen Kapitelstruktur, persistente Zustaende und reproduzierbare Tests der Textausgabe.

\textbf{Formalisierung (Tupel/STRIPS/Logik):} Die Welt bleibt ein logisches Modell, aber $S$ (aktueller Zustand) wird serialisiert und wieder eingelesen. Die Aktionsmenge $\Sigma$ bleibt gleich, waehrend Inhalte modular nachgeladen werden.

\textbf{Prolog:} In Scryer Prolog sind \texttt{consult/1} fuer modulare Wissensbasen, \texttt{read\_term/2} fuer strukturierte Zustandsdaten und \texttt{with\_output\_to/2} fuer testbare Ausgabetexte besonders nuetzlich.

\begin{lstlisting}[language=Prolog,caption={}]
lade_kapitel(Datei) :-
    consult(Datei).

lade_zustand(Datei, Zustand) :-
    open(Datei, read, S),
    read_term(S, Zustand, []),
    close(S).

status_als_text(Text) :-
    with_output_to(string(Text),
        (   format("Ort: labor~n", []),
            format("Tuer: archivtuer verschlossen~n", []),
            format("Ziel: akte bergen~n", [])
        )).

% ?- lade_kapitel('kapitel_institut.pl').
% ?- lade_zustand('save_institut.pl', Z).
% ?- status_als_text(T).
\end{lstlisting}

\subsubsection{Vollständige Spielstand-Serialisierung}

Für ein speicherbares Spiel müssen \textit{alle} dynamischen Prädikate gesichert werden – Positionen, Inventar, Zustände, besuchte Räume, gelesene Objekte, Rätselfortschritt. Der sauberste Weg: \texttt{findall/3} sammelt alle Fakten eines Prädikats in einer Liste, diese werden als ein einziger Term geschrieben und später wieder eingelesen.

\begin{lstlisting}[language=Prolog, caption={Spielstand als Term serialisieren}]
speichere_spielstand(Datei) :-
    findall(position(E,O),  position(E,O),   Pos),
    findall(inventar(S,G),  inventar(S,G),   Inv),
    findall(zustand(E,Z),   zustand(E,Z),    Zst),
    findall(besucht(S,L),   besucht(S,L),    Bes),
    findall(gelesen(S,O),   gelesen(S,O),    Gel),
    Zustand = spielstand(Pos,Inv,Zst,Bes,Gel),
    open(Datei, write, Stream),
    write_term(Stream, Zustand, [quoted(true)]),
    write(Stream, '.'),
    nl(Stream),
    close(Stream).

lade_spielstand(Datei) :-
    open(Datei, read, Stream),
    read_term(Stream, spielstand(Pos,Inv,Zst,Bes,Gel), []),
    close(Stream),
    retractall(position(_,_)),
    retractall(inventar(_,_)),
    retractall(zustand(_,_)),
    retractall(besucht(_,_)),
    retractall(gelesen(_,_)),
    maplist(assertz, Pos),
    maplist(assertz, Inv),
    maplist(assertz, Zst),
    maplist(assertz, Bes),
    maplist(assertz, Gel).
\end{lstlisting}

Der geschriebene Term sieht auf Dateiebene so aus:

\begin{lstlisting}[language=Prolog, caption={Serialisierter Spielstand (gekürzt)}]
spielstand(
    [position(ich, labor), position(brecheisen, eingang),
     position(id_karte, korridor), position(taschenlampe, labor)],
    [inventar(ich, id_karte), inventar(ich, brecheisen)],
    [zustand(labortuer, offen), zustand(archivtuer, verschlossen)],
    [besucht(ich, [eingang, korridor, labor])],
    [gelesen(ich, buch)]
).
\end{lstlisting}

\texttt{write\_term/3} mit \texttt{quoted(true)} stellt sicher, dass Atome mit Sonderzeichen korrekt escaped werden. Beim Laden rekonstruiert \texttt{maplist(assertz, \,\ldots)} den gesamten Zustand aus den Listen – ohne Schleife, ohne manuelles Iterieren.

Für produktive Spiele empfiehlt sich zusätzlich \texttt{catch/3} gegen fehlende oder korrupte Dateien:

\begin{lstlisting}[language=Prolog, caption={Fehlerbehandlung beim Laden}]
lade_spielstand(Datei) :-
    catch(
        (   open(Datei, read, Stream),
            read_term(Stream, spielstand(Pos,Inv,Zst,Bes,Gel), []),
            close(Stream),
            retractall(position(_,_)),
            ...
            maplist(assertz, Gel)
        ),
        _,
        format("Kein Spielstand gefunden.~n", [])
    ).
\end{lstlisting}

\textbf{OWA und nicht-monotone Logik:}In vielen IF-Welten gilt die Open World Assumption (OWA), was nicht bekannt ist, ist nicht automatisch falsch. Nicht-monotone Logik ergänzt das um revidierbare Schlüsse, sodass neue Informationen frühere Annahmen korrigieren können.

\begin{lstlisting}[language=Prolog,caption={}]
:- dynamic fact_true/1, fact_false/1.

status(Fakt, wahr)     :- fact_true(Fakt), !.
status(Fakt, falsch)   :- fact_false(Fakt), !.
status(_,    unbekannt).

% hat_spur/1: prueft ob ein Objekt eine Spur hinterlaesst.
% Muss als Spielregel definiert werden, z.B.:
% hat_spur(X) :- position(spur, R), position(X, R).
% Hier vereinfacht: die archivtuer hat eine sichtbare Spur.
hat_spur(archivtuer).

untersuchen(Objekt) :-
    (   hat_spur(Objekt)
    ->  assertz(fact_true(spur(Objekt)))
    ;   assertz(fact_false(spur(Objekt)))
    ).

% Anfrage vor Untersuchung:
% ?- status(spur(tuer), S).   % S = unbekannt
% Nach untersuchen(tuer):
% ?- status(spur(tuer), S).   % S = wahr oder falsch
\end{lstlisting}

Solange ein Raum nicht untersucht wurde, liefert \texttt{status/2} den Wert \texttt{unbekannt}. Erst nach expliziter Aktion wird die Tatsache in die Wissensbasis eingetragen und gilt als gesichert. Damit entsteht echte Ermittlungslogik statt sofortiger Gewissheit.

\textbf{Metaprozessierung und regelverändernde Welten:} Die Spielwelt kann nicht nur Zustände ändern, sondern auch ihre eigenen Regeln umschreiben, analysieren oder offenlegen. Dadurch entstehen Rätsel, bei denen das Verstehen des Systems selbst Teil der Lösung wird. Ein Meta-Rätsel könnte die Antwort aus der Codestruktur lesen:

\begin{lstlisting}[language=Prolog,caption={}]
:- meta_predicate geheimcode(?).
geheimcode(42).
geheimcode(17).
geheimcode(99).

anzahl_klauseln(Name/Arity, N) :-
    functor(H, Name, Arity),
    findall(_, clause(H, _), Cs),
    length(Cs, N).

raetsel_loesung(Code) :-
    anzahl_klauseln(geheimcode/1, Code).

% ?- raetsel_loesung(L).   % L = 3
\end{lstlisting}

Die Antwort auf das Rätsel ist nicht im Spieltext genannt, sondern in der Struktur des Codes selbst kodiert. Spielende müssen lernen, Prolog meta-prädikativ zu befragen: nicht \enquote{Was gilt?}, sondern \enquote{Wie viele Regeln hat dieses Prädikat?} Das macht das Codelesen zur spielerischen Handlung.

\textbf{Mehrere Handlungsstränge und Bedeutungsebenen:} Entscheidungen der Spielenden können unterschiedliche narrative Pfade und Interpretationen aktivieren. So entsteht Bedeutung nicht in einer einzigen Storyline, sondern im Zusammenspiel von Perspektive, Reihenfolge und Konsequenz. Perspektivabhängige Pfade:

\begin{lstlisting}[language=Prolog,caption={}]
:- dynamic flag/1.

pfad(start, archiv,  neutral).
pfad(start, krypta,  mystisch).
pfad(start, labor,   rational).

aktive_perspektive(mystisch)  :- flag(glaubt_an_omen), !.
aktive_perspektive(rational)  :- flag(vertraut_daten), !.
aktive_perspektive(neutral).

naechste_szene(Von, Nach) :-
    aktive_perspektive(P),
    pfad(Von, Nach, P).

% ?- assertz(flag(glaubt_an_omen)),
%    naechste_szene(start, N).   % N = krypta
\end{lstlisting}

Dieselbe Ausgangssituation führt zu unterschiedlichen Orten, abhängig davon, welche Flags gesetzt sind. Perspektive ist hier kein Stilmittel, sondern ein formales Prädikat, das Pfade steuert.

\textbf{Zeit- und Umweltdynamik:} Tag-Nacht-Zyklen, Wetter und ähnliche Prozesse verändern Erreichbarkeit, Sichtbarkeit und Verhalten von Figuren. Die Welt wirkt dadurch lebendig und reagiert nicht nur auf Befehle, sondern auch auf interne Dynamiken. Diskreter Zeittakt mit Tagesphase und Wetter:

\begin{lstlisting}[language=Prolog,caption={}]
:- dynamic zeit/1, wetter/1.
zeit(0).
wetter(klar).

tick :-
    retract(zeit(T0)),
    T1 is T0 + 1,
    assertz(zeit(T1)),
    update_wetter(T1).

tagesphase(tag)   :- zeit(T), M is T mod 24, M >= 6, M < 20.
tagesphase(nacht) :- zeit(T), M is T mod 24, (M < 6 ; M >= 20).

update_wetter(T) :-
    (   T mod 5 =:= 0 -> set_wetter(nebel)
    ;   T mod 7 =:= 0 -> set_wetter(regen)
    ;   set_wetter(klar)
    ).

set_wetter(W) :- retractall(wetter(_)), assertz(wetter(W)).

sichtbar(Objekt) :-
    tagesphase(tag),
    wetter(klar),
    nicht_versteckt(Objekt).
\end{lstlisting}
Jeder Spielzug kann \texttt{tick/0} aufrufen und damit die Weltzeit voranbewegen. Rätsel lassen sich so an Tageszeit oder Wetter knüpfen, sodass manche Objekte nur nachts sichtbar oder bestimmte Räume nur bei Nebel betretbar sind.

\textbf{Zufall und prozedurale Generierung:} Teile der Welt können bei jedem Durchlauf neu erzeugt oder stochastisch variiert werden. Das erhöht Wiederspielbarkeit und erzeugt emergente Situationen, verlangt aber klare Designregeln für Fairness und Nachvollziehbarkeit. Prozedurale Weltgenerierung mit Fairness-Bedingung:

\begin{lstlisting}[language=Prolog,caption={}]
:- use_module(library(random)).
% position/2 bereits als :- dynamic deklariert (s. Weltblock)

generiere_welt :-
    retractall(position(schluessel, _)),
    retractall(position(laterne, _)),
    random_member(SchluesselOrt, [eingang, korridor, labor]),
    assertz(position(schluessel, SchluesselOrt)),
    random_member(LaternOrt, [eingang, korridor, labor]),
    assertz(position(laterne, LaternOrt)),
    validiere_welt.

validiere_welt :-
    position(schluessel, S),
    position(laterne, L),
    S \= L.  % Schluessel und Laterne sollen nicht am selben Ort liegen

% ?- generiere_welt.
% ?- position(schluessel, Wo).   % wechselt bei jedem Aufruf
\end{lstlisting}
Schlüssel und Hinweis werden zufällig platziert, aber nie am selben Ort. \texttt{validiere\_welt/0} stellt als einfache Fairness-Bedingung sicher, dass die Welt lösbar bleibt. Komplexere Validierungen lassen sich als weitere Regeln ergänzen.

\section{Code als kulturelle Praxis}

Programmiersprachen sind nicht neutral. Der Linguist Noah Bubenhofer hat gezeigt, dass die Wahl einer Programmiersprache keine technische Entscheidung ist, sondern eine kulturelle \citep{Bubenhofer2015}. Man ordnet sich damit einem Denkstil unter, reproduziert ihn und grenzt sich von anderen ab. Der Begriff des Denkstils geht auf den Erkenntnistheoretiker Ludwik Fleck zurück \citep{Fleck1935}. Ein wissenschaftliches Gerät, also auch Software, ist niemals nur Werkzeug, sondern immer auch Träger der Denkkulturen, in denen es entstanden ist. Wie natürliche Sprachen tragen Programmiersprachen Konventionen, Vorannahmen, Weltbilder. Sie entscheiden mit, was sich ausdrücken lässt und was nicht.

Code tut dabei etwas, was kein anderes Schriftsystem kennt. Florian Cramer hat gezeigt \citep{Cramer2005}, dass er nicht nur beschreibt, sondern \textit{vollzieht},  er ist Text und Maschine, Notation und Ausführung zugleich. Er strukturiert wie wir Probleme lösen und wie wir sie überhaupt denken. In diesem Sinne ist Code ein Repräsentationssystem, das gesellschaftliche Werte naturalisiert und bestimmt, was sichtbar wird, welche Bilder entstehen und welche Logiken als gegeben gelten.

Prolog macht diese Eigenschaft besonders sichtbar. Als deklarative Sprache sagt sie nicht wie etwas geschehen soll, sondern was wahr ist und überlässt dem System, was daraus folgt. Das ist strukturell näher an der Arbeit der Gestaltung als jede imperative Sprache. Ein Designsystem legt keine Einzelentscheidungen fest, sondern die Regeln ihrer Beziehungen. Ian Bogost nennt das \textit{prozedurale Rhetorik}: Regeln und Systeme treffen Aussagen über die Welt, nicht durch Beschreibung, sondern durch ihr Verhalten \citep{Bogost2007}. Friedrich Kittler hat in „Es gibt keine Software" \citep{Kittler1993} darauf hingewiesen, dass Code das einzige Schriftsystem ist, das sich selbst ausführt und die Differenz zwischen Notation und Ausführung kollabiert, die alle anderen Schriften kennzeichnet. Wer in Prolog eine Welt entwirft, entwirft ihre Ontologie, Fakten als Sein, Regeln als Physik, Anfragen als Erfahrung.

Wer in diesem Kurs eine solche Welt gebaut hat, hat damit mehr getan als programmiert. Wir haben gelernt, Regeln so zu setzen, dass aus ihrem Zusammenspiel Bedeutung entsteht. Das ist Kulturtechnik als Form des Handelns, die tradiert, variiert und durch Praxis erworben wird, so wie Lesen, Schreiben, Rechnen. Seymour Papert hat gezeigt, dass Programmieren genau dann zum Denkwerkzeug wird, wenn es nicht als Technik, sondern als Sprache des Entwerfens begriffen wird: „Learning by Making" meint nicht das Abarbeiten von Aufgaben, sondern das Bauen von Objekten, an denen man denken kann \citep{Papert1980}. Code als Grundlagengestaltungsfeld, neben Bild und Text. Sie bedeutet, das wir nicht Werkzeuge bedienen, sondern eine Sprache sprechen, die Wirklichkeit nicht abbildet, sondern erzeugt.

Für das Kommunikationsdesign hat das Konsequenzen, die über diesen Kurs hinausreichen. Generative Bildsprache, adaptive Interfaces, prozedurale Identitäten, das sind keine Informatikprobleme, sondern Gestaltungsaufgaben. Sie verlangen eine Kompetenz, die hier geübt wurde, nämlich die \textit{Fähigkeit}, \textit{explizit zu formulieren, was gelten soll}. Große Sprachmodelle können überzeugend formulieren, aber sie kennen keine Wahrheitswerte. Sie wissen nicht, ob die Tür offen ist. Das logische Gerüst, die Fakten, Regeln, Weltkonsistenz, ist genau das, was sie nicht mitbringen. Wer es entwerfen kann, kann sie gestalterisch steuern, statt von ihnen gestaltet zu werden.

Christopher Manson hat 1985 mit MAZE ein Buch gebaut, das wie ein Prolog-Programm funktioniert: Raum und Zeit als Regelsystem, Lesende als Inferenz-Maschinen, die Bedeutung liegt in der Struktur selbst, nicht im Inhalt. Er wusste nichts von Prolog, aber er hat dasselbe gedacht. Das Labyrinth war die erste Metapher dieses Kurses und bleibt die treffendste: eine Struktur, die keine Übersicht gewährt, sondern nur den nächsten Schritt. Wer es durchquert hat, hat nicht Programmieren gelernt, sondern einen Denkstil erworben, der das Entwerfen regelbasierter Welten als gestalterische Praxis begreift. Und das ist, in einer Gegenwart, die von unzähliger solcher Welten durchzogen ist, eine der wichtigsten Kompetenzen im Kommunikationsdesign.

\section{Abschlusspräsentation}

Die Abschlusspräsentation findet in der letzten Kurseinheit statt. Jede Gruppe präsentiert ihr Projekt, erklärt die Logik hinter ihren Rätseln und diskutiert die gestalterischen Entscheidungen. Das Spiel wird live vorgeführt, damit wir die Welt erleben und die Interaktion nachvollziehen können.

\subsection{Inhalt der Präsentation}

In der Präsentation sollten folgende Aspekte beleuchtet werden, jeweils mit Begründung auf Basis der im Kurs behandelten Begriffe und Konzepte, euer Vorgehen und auch die Probleme, die es gab. Die folgenden Punkte sollen als Gerüst dienen, nicht jeder Punkt muss beantwortet werden

\subsubsection{Literarische und philosophische Überlegungen}

\begin{itemize}
    \item Wie lautet der Plot und Grundidee der Geschichte?
    \item Wie weit hält sie sich an die Konventionen des IF und wo bricht es sie?
    \item Welche Bedeutungsebenen habt ihr verwendet?
    \item Wie ist die Welt erfahrbar? 
\end{itemize}

\subsubsection{Gestalterische und narrative Entscheidungen}

\begin{itemize}
    \item Ist die Welt erfahrbar? 
    \item Wie habt ihr eure Welt strukturiert? 
    \item Welche Räume, Charaktere und Objekte existieren in der Welt? 
    \item Welche Rolle spielen sie?
    \item Wie hängen sie zusammen? 
    \item Welche Geschichte erzählt eure Welt? 
    \item Welche Atmosphäre wollt ihr schaffen? 
    \item Wie spiegelt sich das in der Gestaltung der Räume, Charaktere und Rätsel wider?
    \item Welche Spielmechaniken habt ihr eingebaut?
    \item Welche Herausforderungen erwarten die Spieler:innen?
    \item Wie funktionieren die Rätsel und sind sie fair konstruiert?
\end{itemize}

\subsubsection{Interaktion}

Welche Befehle können die Spieler:innen nutzen? Wie reagieren die Räume und Charaktere auf die Aktionen der Spieler:innen?

\begin{itemize}
    \item Wie interagiert textbasiert man mit der Welt? 
    \item Wie kann ich den Zustand der Weltverändern?
    \item Gibt es ein Inventar oder ähnliche Mechanismen?
    \item Welche Befehle gibt es?
\end{itemize}

\subsubsection{Logik}

Wie habt ihr die Logik der Welt in Prolog umgesetzt?

\begin{itemize}
\item Welche logischen Strukturen habt ihr verwendet?
\item Auf welche Weise beeinflussen sie die Erzählung und die Spielmechaniken?
\item Auf welche Konzepte der Logik habt ihr zurückgegriffen?
\item Auf welche Probleme seid ihr dabei gestoßen? 
\item Gab es logische Paradoxien oder Aporien in eurer Welt?
\end{itemize}

\subsubsection{Technische Umsetzung mit Prolog}
\begin{itemize}
\item Wie habt ihr die Welt in Prolog umgesetzt?
\item Welche Fakten, Regeln und Datenstrukturen habt ihr verwendet und warum?
\item Wie funktioniert der Game Loop?
\item Wie wird der Weltzustand und der Charakterzustand verwaltet?
\item Wie werden Aktionen verarbeitet und wie verändern sie die Welt?
\item Welche Mechanismen habt ihr für die Rätsel implementiert? Wie funktionieren sie?
\item Habt ihr einen Parser oder eine eigene DSL für die Interaktion entwickelt? Wie funktionieren diese?
\item Welche Herausforderungen gab es bei der technischen Umsetzung und wie habt ihr sie gelöst?
\end{itemize}

\subsection{Ende und Ausblick}

Die vier Fragen vom Anfang, was existiert, wie hängt es zusammen, was darf passieren, was ändert sich wann, könnt ihr jetzt beantworten. Nicht theoretisch, sondern als Entwurfspraxis.

Dieser Kurs war eine verwinkelte Reise durch ein Labyrinth aus Sprache, Logik, Philosophie und Informatik. Dabei haben wir das Labyrinth als Strukturen des Nicht-Wissens, in denen Orientierung, Entscheidung und Bedeutung erst im Handeln entstehen, betreten und Räume des sprachlichen und logischen Wissens entdeckt, die uns herausgefordert haben, unsere Perspektive zu erweitern, um neue Wege zu entdecken. 

Ich hoffe, dass dieser Kurs genau das erreicht hat und euch nicht nur die Werkzeuge an die Hand gegeben hat, um neue eigene Welten zu denken, sondern auch ihre logischen Strukturen zu verstehen, sie zu formulieren und so zum Leben zu erwecken, egal wie phantastisch sie sind. Das war harte Arbeit, aber genau darin liegt ja auch der gestalterische Kern ergodischer Literatur udn Interactive Fiction. Eben nicht beim kleinsten Hindernis aufzugeben, sondern dran zu bleiben.

In Zukunft könnten große Sprachmodelle (Large Language Models) diese logischen Gerüste ergänzen, indem sie, anstelle des Parsers eine nahezu grenzenlose narrative Freiheit und dynamische Interaktion in natürlicher Sprache ermöglichen. Um so mehr müssen die Grenzen und Regeln definiert werden, die die Modelle gestalterisch lenken und einschränken, um nicht in der Beliebigkeit zu enden. 

Die Forschung steht hier erst am Anfang, und die Präzision der Logik mit den Möglichkeiten großer Sprachmodelle zu verbinden. Diese Verbindung könnte ein interessantes zukünftiges Forschungsfeld im Kommunikationsdesign, an den Schnittstellen von Sprache, Literatur, Informatik und Spielwissenschaft sein.

\bibliographystyle{plainnat}
\bibliography{literature}

% =============================================================================
\chapter{Bonus: Dynamische Welten und prozedurale Generierung}
% =============================================================================

Das Textadventure \textit{Das Institut} arbeitet mit einem \textbf{statischen Weltmodell}: alle Räume, Verbindungen und Türzustände sind als unveränderliche Fakten deklariert. Diese Entscheidung war bewusst, denn das statische Modell lässt sich vollständig mit den Grundkonzepten von Prolog erklären, die im Kurs behandelt wurden: Fakten, Regeln, Unifikation, Backtracking, \texttt{assertz}/\texttt{retract} auf Gegenständen und Zuständen.

\textbf{Prozedurale Generierung} erfordert darüber hinaus: dynamische Topologie (also \texttt{assertz} auf \texttt{verbindet/2} und \texttt{tuer/3} zur Laufzeit), Graphsuche zur Erreichbarkeitsgarantie, Zufallsauswahl aus Pools sowie Constraint-Prüfung bei der Platzierung. Jedes dieser Themen ist ein eigenes, kursumfangswürdiges Gebiet. Wer dieses Kapitel liest, hat den Kurs abgeschlossen und kann die hier vorgestellten Konzepte als Ausgangspunkt für eigene Forschung und Projekte nutzen.

% -----------------------------------------------------------------------------
\section{Einfache dynamische Raumgenerierung}
% -----------------------------------------------------------------------------

Der einfachste Schritt weg vom statischen Modell: Raumverbindungen werden beim Spielstart aus einem \textit{Pool} zufällig zusammengesetzt, anstatt fest deklariert zu sein.

\begin{lstlisting}[language=Prolog, caption={Einfache prozedurale Weltgenerierung (Scryer-kompatibel)}]
:- use_module(library(random)).

raum_pool([hoehle, bibliothek, werkstatt, garten, keller]).

% Scryer stellt kein random_permutation/2 bereit (SWI-Prolog-Praedikat).
% Wir implementieren es ueber random_member/2 und select/3:
% (benoetigt library(random), s. use_module-Direktive oben)
zufalls_permutation([], []).
zufalls_permutation(Liste, [Wahl|Rest]) :-
    random_member(Wahl, Liste),  % benoetigt library(random)
    select(Wahl, Liste, Verbleibend),
    zufalls_permutation(Verbleibend, Rest).

% Hinweis: Nur R1-R3 werden verbunden; R4 und R5 bleiben isoliert.
% generiere_bis_erreichbar/1 muesste fuer dieses Beispiel daher haeufig
% neu generieren. Fuer ein vollstaendiges Modell sollte der gesamte Pool
% zu einer zusammenhaengenden Kette verbunden werden.
generiere_welt :-
    raum_pool(Pool),
    zufalls_permutation(Pool, [R1, R2, R3 | _]),
    assertz(verbindet(R1, R2)),
    assertz(verbindet(R2, R1)),
    assertz(verbindet(R2, R3)),
    assertz(verbindet(R3, R2)).
\end{lstlisting}

\texttt{zufalls\_permutation/2} wählt per \texttt{random\_member/2} iterativ Elemente aus der Liste; \texttt{assertz} schreibt die erzeugten Verbindungen in die Datenbank. Nach \texttt{generiere\_welt/0} gilt \texttt{verbindet(R1,R2)} genau wie ein statisch deklariertes Fakt. \textit{SWI-Prolog-Code kann stattdessen direkt \texttt{random\_permutation/2} aus \texttt{library(random)} nutzen -- in Scryer ist dieses Prädikat nicht enthalten.}

\textbf{Offene Frage:} Wie stellt man sicher, dass das Spielziel nach der Generierung erreichbar ist? Eine Breitensuche mit \texttt{findall} und Fixpunktiteration über \texttt{verbindet/2} liefert die Menge aller erreichbaren Räume vom Startpunkt aus. Ist das Zielzimmer nicht darin enthalten, wird die Welt neu generiert:

\begin{lstlisting}[language=Prolog, caption={Erreichbarkeitsgarantie per Fixpunktiteration}]
erreichbare_menge(Von, Menge) :-
    erreichbare_menge([Von], [], Menge).

erreichbare_menge([], Gesehen, Gesehen).
erreichbare_menge([H|T], Gesehen, Menge) :-
    (   member(H, Gesehen)
    ->  erreichbare_menge(T, Gesehen, Menge)
    ;   findall(N, verbindet(H, N), Nachbarn),
        append(T, Nachbarn, Queue),
        erreichbare_menge(Queue, [H|Gesehen], Menge)
    ).

generiere_bis_erreichbar(Ziel) :-
    generiere_welt,
    erreichbare_menge(eingang, Menge),
    (   member(Ziel, Menge)
    ->  true
    ;   retractall(verbindet(_, _)),
        generiere_bis_erreichbar(Ziel)
    ).
\end{lstlisting}

Prolog-Konzepte: \texttt{assertz}/\texttt{retractall} auf Topologie, \texttt{library(random)}, Graphsuche als Akkumulator-Rekursion.

% -----------------------------------------------------------------------------
\section{Der Wave Function Collapse Algorithmus}
% -----------------------------------------------------------------------------

Der \textbf{Wave Function Collapse Algorithmus} (kurz WFC) wurde 2016 von Maxim Gumin auf GitHub veröffentlicht \citep{Gumin2016}. Er ist analog zum gleichnamigen Ereignis in der Quantenmechanik benannt und löst \textit{Constraint Satisfaction Problems} für Kachelgitter: Gegeben ein Tileset mit allen möglichen Nachbarschaftsbeziehungen, erzeuge ein konsistentes Bild, in dem jede Kachel zu ihren Nachbarn passt. Nach seiner Veröffentlichung gewann der Algorithmus durch seine organisch wirkenden Ergebnisse schnell Einfluss auf die Indie-Spieleentwicklung und wird seitdem im Bereich der \textit{Procedural Content Generation} (PCG) eingesetzt, etwa in \textit{Proc Skater 2016} und \textit{Caves of Qud}.

\subsection*{Das Grundprinzip}

\begin{enumerate}
    \item Jede Zelle des Gitters besitzt eine \textbf{Menge möglicher Zustände} (die \enquote{Wellenfunktion}).
    \item Wähle die Zelle mit der geringsten \textbf{Entropie} (kleinste Menge).
    \item \textbf{Kollabiere} sie: wähle zufällig einen Zustand aus ihrer Menge.
    \item \textbf{Propagiere} die Constraints: benachbarte Zellen streichen alle Zustände, die mit dem neuen Zustand unverträglich sind.
    \item Wiederhole, bis alle Zellen kollabiert sind. Bei Widerspruch: zurücksetzen und neu versuchen.
\end{enumerate}

In Prolog lässt sich der Zustand einer Zelle als dynamisches Faktum mit einer Zustandsliste darstellen:

\begin{lstlisting}[language=Prolog, caption={WFC-Grundgerüst in Prolog}]
:- dynamic zelle/2.  % zelle(Pos, Moeglichkeiten)

% Constraint: welche Kacheln duerfen nebeneinander stehen?
vertraeglich(flur, tuer).
vertraeglich(tuer, raum).
vertraeglich(wand, wand).
vertraeglich(raum, flur).

% Entropie: Laenge der Moeglickeiten-Liste
entropie(Pos, N) :-
    zelle(Pos, Ms), length(Ms, N).

% Schritt: Zelle mit minimaler Entropie kollabieren.
% N > 1 schliesst bereits kollabierte Zellen (Laenge 1) von der Auswahl aus.
wfc_schritt :-
    findall(N-Pos, (entropie(Pos, N), N > 1), Paare),
    sort(Paare, [_-Ziel | _]),        % kleinste Entropie zuerst
    zelle(Ziel, Moeglichkeiten),
    random_member(Wahl, Moeglichkeiten),
    retract(zelle(Ziel, _)),
    assertz(zelle(Ziel, [Wahl])),     % kollabiert
    propagiere(Ziel, Wahl).

% Propagierung: direkte Nachbarzellen aktualisieren (eine Ebene).
% Achtung 1: Propagierung ist nicht rekursiv -- aendert sich eine Nachbarzelle,
%   werden deren Nachbarn nicht automatisch aktualisiert (kein echtes WFC).
% Achtung 2: Falls include/3 eine leere Liste liefert (Widerspruch), schlaegt
%   forall fehl -- bereits ausgefuehrte retract/assertz fuer fruehere Nachbarn
%   bleiben jedoch unwiderruflich in der Datenbank (Seiteneffekte).
propagiere(Pos, Wahl) :-
    forall(
        nachbar(Pos, NachbarPos),
        (   zelle(NachbarPos, Alt),
            include(vertraeglich(Wahl), Alt, Neu),
            Neu \== [],               % Widerspruch: Neu ist ground nach include/3
            retract(zelle(NachbarPos, Alt)),
            assertz(zelle(NachbarPos, Neu))
        )
    ).
\end{lstlisting}

\textbf{Prolog-Konzepte:} \texttt{findall} zur Entropieberechnung, \texttt{sort/2} für die Minimalauswahl, \texttt{forall} für die Constraint-Propagierung, \texttt{include/3} zum Filtern verträglicher Zustände.

\textbf{Zwei wichtige Einschränkungen dieses Grundgerüsts:}
\begin{itemize}
    \item \textbf{Propagierung ist nicht rekursiv.} Echter WFC propagiert Constraints bis zur Fixpunktkonvergenz (Arc Consistency): Ändert sich die Menge einer Nachbarzelle, müssen auch deren Nachbarn aktualisiert werden. Das erfordert eine Queue-basierte Schleife -- analog zu \texttt{erreichbare\_menge/3} aus Abschnitt~1.
    \item \textbf{\texttt{assertz}/\texttt{retract} sind extralogische Seiteneffekte:} Prolog-Backtracking macht sie \textit{nicht} rückgängig. Ist eine Zelle einmal in die Datenbank geschrieben, bleibt das so -- auch wenn spätere Schritte zu einem Widerspruch führen. Für echtes Rücksetzen bei Widersprüchen braucht man einen expliziten Checkpoint: den Zustand per \texttt{findall} als Term sichern, mit \texttt{retractall} löschen und per \texttt{maplist(assertz,\,\ldots)} neu aufbauen. Das ist exakt das Muster von \texttt{speichere\_spielstand}/\texttt{lade\_spielstand} aus \textit{Das Institut}.
\end{itemize}

Eine vollständige Implementierung des Grundprinzips für Raumkarten findet sich im Kursarchiv unter \texttt{archive/WavefunctionCollapse.pl}.

% -----------------------------------------------------------------------------
\section{Blue Prince: Drafting-Mechanik und Koordinatengitter}
% -----------------------------------------------------------------------------

\textit{Blue Prince} (2025) kombiniert zwei Konzepte, die sich direkt in Prolog-Idiome übersetzen lassen.

\subsection*{Drafting}

Anstatt eine feste Welt zu bauen, wählt der Spieler aus drei zufällig gezogenen Räumen aus, welcher als nächster an der aktuellen Tür platziert wird. Der Spieler \textit{entwirft} die Welt beim Spielen. In Prolog:

\begin{lstlisting}[language=Prolog, caption={Drafting-Mechanik: Spieler wählt Raum aus drei Kandidaten}]
:- dynamic platziert/2.  % platziert(Raum, pos(X,Y))

draft_schritt(Kandidaten, Wahl, NeuePos) :-
    member(Wahl, Kandidaten),          % Spieler trifft Wahl
    assertz(platziert(Wahl, NeuePos)),
    verbinde_mit_nachbarn(Wahl, NeuePos).
\end{lstlisting}

\subsection*{Koordinatengitter und Constraint-Prüfung}

Jeder Raum hat eine \texttt{(X,Y)}-Position im unsichtbaren Gitter. Türen haben Richtungen (\texttt{nord}, \texttt{sued}, \texttt{ost}, \texttt{west}). Beim Platzieren eines Raums wird geprüft, ob benachbarte platzierte Räume passende Anschlüsse haben:

\begin{lstlisting}[language=Prolog, caption={Verbindung nach Koordinatenlogik}]
% Nachbarposition berechnen
nachbar_pos(pos(X,Y), nord, pos(X,Y1)) :- Y1 is Y + 1.
nachbar_pos(pos(X,Y), sued, pos(X,Y1)) :- Y1 is Y - 1.
nachbar_pos(pos(X,Y), ost,  pos(X1,Y)) :- X1 is X + 1.
nachbar_pos(pos(X,Y), west, pos(X1,Y)) :- X1 is X - 1.

% Gegenrichtung
gegenueber(nord, sued). gegenueber(sued, nord).
gegenueber(ost,  west). gegenueber(west, ost).

% Beim Platzieren: alle Nachbarraeume verbinden, wo Tueren passen.
% Hinweis: forall + assertz, falls die Action-Seite fuer einen Nachbarn
% fehlschlaegt, bleiben frueherre assertz-Aufrufe in der Datenbank
% (Seiteneffekte). Da hier nur Fakten hinzugefuegt (kein retract) werden,
% ist das Risiko geringer als im WFC-Code, aber strukturell dasselbe.
verbinde_mit_nachbarn(Raum, Pos) :-
    forall(
        (   tuer_richtung(Raum, Richtung),
            nachbar_pos(Pos, Richtung, NachbarPos),
            platziert(NachbarRaum, NachbarPos),
            tuer_richtung(NachbarRaum, Gegenrichtung),
            gegenueber(Richtung, Gegenrichtung)
        ),
        (   assertz(verbindet(Raum, NachbarRaum)),
            assertz(verbindet(NachbarRaum, Raum))
        )
    ).
\end{lstlisting}

\textbf{Prolog-Konzepte:} Arithmetik auf Koordinatentermen, \texttt{forall} für Constraint-Checks, dynamische Topologie, Nutzung des vorhandenen \texttt{verbindet/2}-Musters aus \textit{Das Institut}.

\subsection*{Zustandsexplosion}

\textit{Blue Prince} operiert pro Spielsitzung mit einer großen Zahl expliziter Zustandsgrößen, darunter verbleibende Schritte, Ressourcen wie Schlüssel und der aktuelle Erkundungsfortschritt. Damit folgt das Modell strukturell dem \texttt{zustand/2}-Prinzip aus \textit{Das Institut}, jedoch in einer deutlich größeren und stärker verzweigten Faktenbasis. Zugleich ist das System nicht als reine Gegenüberstellung von statischem und dynamischem Design zu beschreiben: Raum-Blueprints und Regeln sind vorentworfen (statische Ebene), während konkrete Auswahl, Reihenfolge und räumliche Anordnung erst im Spielverlauf durch Drafting und Entscheidungen entstehen (dynamische Ebene). Für die Persistenz gilt daher eine präzise Einschränkung: Der relevante Weltzustand muss an den vom Spiel vorgesehenen Übergängen konsistent fortgeschrieben werden, aber nicht notwendigerweise zu jedem beliebigen Zeitpunkt innerhalb eines laufenden Tages.

\subsection*{Weiterführende Literatur}

\begin{itemize}
    \item \citet{Shaker2016}: \textit{Procedural Content Generation in Games}. Umfassendes Open-Access-Lehrbuch zu PCG, Constraint-Satisfaction und Levelgenerierung.
    \item Gumin, M. (2016). \textit{WaveFunctionCollapse}. \url{https://github.com/mxgmn/WaveFunctionCollapse}
    \item \citet{Russell2020}: \textit{Artificial Intelligence: A Modern Approach}. Kapitel zu Constraint Satisfaction Problems und Suchstrategien.
\end{itemize}

\end{document}